% DER_XIII-a: Grothendieck's text with the commentary set in, then the overview, inside a group that applies the chunk's own definitions
\dgrchapter{XIII}{Chapter XIII}{Catégories de modèles (1)}
\dgrpart{0}{XIII-a}{Part 1 of 2, manuscript pages 1--60}{XIII-a\enspace\textperiodcentered\enspace Chapter XIII, part 1 of 2\enspace\textperiodcentered\enspace manuscript pages 1--60}{Chapter XIII \textperiodcentered\ \textup{XIII-a}}
\dgrchunkbegin
% DER_XIII-a: the definitions of the chunk's preamble, applied inside its group
\let\printedpage\relax\newcommand{\printedpage}[1]{\marginpar{\footnotesize p.~#1}}
\let\msp\relax\newcommand{\msp}[1]{\marginpar{\footnotesize [#1]}}
\let\Fl\relax\newcommand{\Fl}{\mathrm{Fl}}
\let\M\relax\newcommand{\M}{\mathcal{M}}
\let\cM\relax\newcommand{\cM}{\mathcal{M}}
\let\Cof\relax\newcommand{\Cof}{\mathrm{Cof}}
\let\Fib\relax\newcommand{\Fib}{\mathrm{Fib}}
\let\Ob\relax\newcommand{\Ob}{\mathrm{Ob}}
\let\circled\relax\newcommand{\circled}[1]{\tikz[baseline=(c.base)]\node[draw,circle,inner sep=1pt](c){#1};}
\let\univ\relax\newcommand{\univ}{\mathrm{univ}}
\let\cercle\relax\newcommand{\cercle}[1]{\mbox{\tikz[baseline=(c.base)]\node[draw,circle,inner sep=1pt](c){#1};}}
\let\clarrow\relax\newcommand{\clarrow}{\mathrel{\leftarrow\mkern-4mu{+}\mkern-4mu\rightarrow}}
\let\limind\relax\newcommand{\limind}{\mathop{\underrightarrow{\mathrm{lim}}}}
\let\limproj\relax\newcommand{\limproj}{\mathop{\underleftarrow{\mathrm{lim}}}}
\let\HO\relax\newcommand{\HO}{\mathrm{HO}}

% ---- from Der-ChXIII-25-01-04_p001-005.tex
\printedpage{--} % no page number printed on this page
% head: (none)

\begin{center}
\vspace*{5cm}
{\LARGE\sffamily LES DÉRIVATEURS}

\vspace{1.5cm}
{\large\textbf{ALEXANDRE GROTHENDIECK}}

\vspace{2.5cm}
{\large\textbf{Chapitre XIII}}

\vspace{0.5cm}
{\large\textbf{Catégories de modèles (1)}}
\end{center}

\clearpage

\printedpage{--} % no page number printed on this page
% head: (none)

\vspace*{6cm}
\noindent Ce texte a été déchiffré et transcrit en \LaTeXe{} par M. Künzer. Il a été édité par M. Künzer, J. Malgoire, et G. Maltsiniotis. La transcription est aussi fidèle que possible au manuscrit. Pour les quelques corrections évidentes, ou rares commentaires des éditeurs, ainsi que pour la numérotation originale des pages du manuscrit, les caractères de machine à écrire \texttt{[typewriter]} entre crochets sont utilisés. Un point d'interrogation entre crochets signifie que l'on n'est pas sûr du mot qui précède. Les index, la liste des axiomes, le glossaire, et les références bibliographiques à la fin du chapitre ont été compilés par les éditeurs. Les numéros des pages dans les index sont ceux de l'original.

\noindent Cette édition est provisoire. Les remarques, commentaires et corrections sont bienvenus. Envoyer un message à :

\noindent\textbf{[e-mail address removed]}

\hfill G. Maltsiniotis\hspace*{2cm}

\clearpage

\printedpage{1}
% head: A. Grothendieck : Dérivateurs, Ch. XIII, Édition des Universités de Montpellier II et Paris VII
% foot: Version 25 janvier 2004

\msp{page 1}

\section*{1\quad Couples de Quillen $(\Phi, \Psi)$}

Je veux d'abord donner une variante affaiblie de la notion de triple de Quillen (``model category'' de Quillen). Ce ne sera pas une notion tout-à-fait autoduale -- ainsi il y aura \textit{deux} variantes affaiblies, duales l'une de l'autre.

\dgrblock{dgrb666}{Two calculi behind a model structure}
\begin{dcomment}{exposition}{p1:Je veux d'abord donner}{Two calculi behind a model structure}
Grothendieck begins by separating the machinery that constructs replacements from the equivalences those replacements are meant to respect. This separation will make it possible to ask whether the homotopy theory belongs to the pair $(\mathcal M,W)$ itself. The first section therefore contains no weak equivalences: its subject is the elementary calculus of lifting and factorization.

For arrows $i:A\to B$ and $p:X\to Y$, write $i\mathrel{\square}p$ when every commutative square with these vertical sides admits a diagonal $h:B\to X$, satisfying $hi=a$ and $ph=b$, where $a:A\to X$ and $b:B\to Y$ are its horizontal sides. No uniqueness is required. In current notation Grothendieck's $\Phi_*$ is $\Phi^{\square}$, and $\Psi^*$ is ${}^{\square}\Psi$. Consequently
\[
 \Phi\subseteq{}^{\square}\Psi
 \quad\Longleftrightarrow\quad
 \Psi\subseteq\Phi^{\square}.
\]
This is an order-reversing Galois connection between classes of arrows. The arrows in the two displayed boxes operate on classes, rather than on individual morphisms; the editors' correction is mathematically essential, while the missing parenthesis in that correction is merely typographical. For a large category we work relative to universes, or speak of classes instead of forming an unrestricted set of all its subclasses.

There are already two different closure notions in view. A class is \emph{left lifting closed} when it has the form ${}^{\square}\Psi$. Later, a class will be enlarged by isomorphisms, pushouts, transfinite composites and retracts; this is its \emph{weak saturation} in current model-category language. These operations agree for the closure generated by a set in a locally presentable category, by the small object argument. Their identification requires that theorem; it is not a consequence of the Galois connection alone. Conversely, Grothendieck's phrase ``$W$ faiblement saturée'' on manuscript p.~7 means two-out-of-three together with the isomorphisms. Keeping these uses separate prevents a genuine ambiguity in the conjectural argument below.
\end{dcomment}


Mais avant encore, un sorite sur les LLP (left lifting property) et RLP (right lifting property) (\footnote[1]{donner définition en forme, avec renvoi à \textsc{Quillen}.}). Soit $\M$ une catégorie quelconque (le plus souvent, $\M$ sera grande). On considère des parties
\[
\Phi,\ \Psi\ \subset\ \Fl(\M)
\]
et on définit des applications
\[
\boxed{\begin{array}{ccc}
\Phi & \longmapsto & \Phi_* \\[1ex]
\Psi & \longmapsto & \Psi^*
\end{array}}
\qquad
\Fl(\M) \longrightarrow \Fl(\M)
\]
\texttt{[plutôt} $\mathfrak{P}(\Fl(\M) \longrightarrow \mathfrak{P}(\Fl(\M))$\texttt{,}\newline\texttt{où pour tout ensemble} $E$\texttt{,} $\mathfrak{P}(E)$ \texttt{désigne l'ensemble des parties de} $E$\,\texttt{]} par % [sic] \mathfrak{P}(\Fl(\M) -> \mathfrak{P}(\Fl(\M)), a closing parenthesis is missing in the first term
\[
\begin{aligned}
\Phi_* \ &= \ \{\text{ensemble des } p \in \Fl(\M) \text{ qui ont la RLP par rapport à } \Phi\} \\
\Psi^* \ &= \ \{\text{ensemble des } i \in \Fl(\M) \text{ qui ont la LLP par rapport à } \Psi\},
\end{aligned}
\]
et ensuite
\[
\boxed{\begin{array}{ccc}
\Phi & \longmapsto & \widetilde{\Phi} \\[1ex]
\Psi & \longmapsto & \overline{\Psi}
\end{array}}
\qquad
\Fl(\M) \longrightarrow \Fl(\M)
\]
\texttt{[plutôt} $\mathfrak{P}(\Fl(\M) \longrightarrow \mathfrak{P}(\Fl(\M))$\texttt{]} par % [sic] \mathfrak{P}(\Fl(\M) -> \mathfrak{P}(\Fl(\M)), a closing parenthesis is missing in the first term
\[
\begin{aligned}
\widetilde{\Phi} \ &= \ (\Phi_*)^* \\
\overline{\Psi} \ &= \ (\Psi^*)_* \ .
\end{aligned}
\]

\msp{page 2}

Je résume les propriétés les plus immédiates de ces opérations dans une

\noindent\textbf{Proposition 1.}

\begin{enumerate}[label={}]
\item[a)] $\Phi \longmapsto \Phi_*$, $\Psi \longmapsto \Psi^*$ \textit{sont des applications décroissantes de} $\Fl(\M)$ \texttt{[plutôt} $\mathfrak{P}(\Fl(\M))$\texttt{]} \textit{dans lui-même, pour la relation d'ordre} $\subset$\textit{, et} $\Phi \longmapsto \widetilde{\Phi}$\textit{,} $\Psi \longmapsto \overline{\Psi}$ \textit{des applications croissantes.}

\printedpage{2}
% head: A. Grothendieck : Dérivateurs, Ch. XIII, Édition des Universités de Montpellier II et Paris VII
% foot: Version 25 janvier 2004

\item[b)] \textit{On a}
\[
\Phi\ \subset\ \widetilde{\Phi},
\]
\textit{et pour qu'on ait $\Phi = \widetilde{\Phi}$, il faut et il suffit que $\Phi$ soit de la forme $\Psi^*$. On dit alors que $\Phi$ est} \textup{clos à gauche} (\footnote[2]{Donner définition en forme.})\textit{. L'ensemble des parties de} $\Fl(\M)$ \textit{closes à gauche est stable par intersections quelconques} (\footnote[3]{de façon précise, $\bigcap_i(\Psi_i^*) = (\bigcup_i \Psi_i)^*$}), \textit{et $\widetilde{\Phi}$ est la plus petite partie close de} $\Fl(\M)$ \textit{qui contienne $\Phi$ et la plus grande partie $\Phi'$ de} $\Fl(\M)$ \textit{telle que $\Phi'_* = \Phi_*$ (ou telle que $\Phi'_* \supset \Phi_*$).}

\textit{Énoncés symétriques pour $\Psi$ : on a}
\[
\Psi\ \subset\ \overline{\Psi},
\]
\textit{on a $\Psi = \overline{\Psi}$ si et seulement si $\Psi$} \texttt{[est]} \textit{de la forme $\Phi_*$, on dit alors que $\Psi$ est} \textup{clos à droite.} \textit{L'ensemble des parties de} $\Fl(\M)$ \textit{closes à droite est stable par intersections quelconques} (\footnote[4]{de façon précise, on a $\bigcap_i(\Phi_{i\,*}) = (\bigcup_i \Phi_i)_*$}), \textit{et $\overline{\Psi}$ est la plus petite partie close à droite contenant $\Psi$, et aussi la plus grande partie $\Psi'$ de} $\Fl(\M)$ \textit{telle que $\Psi'^* = \Psi^*$ (ou telle que $\Psi'^* \supset \Psi^*$).}

\item[] \hspace*{-2.5em}\msp{page 3}

\item[c)] \textit{Les applications}
\[
\Phi \ \longmapsto\ \Phi_*, \qquad\quad \Psi\ \longmapsto\ \Psi^*
\]
\textit{établissent des bijections inverses l'une de l'autre entre l'ensemble des parties $\Phi$ closes à gauche de} $\Fl(\M)$\textit{, et l'ensemble des parties $\Psi$ closes à droite. C'est un anti-isomorphisme d'ensembles ordonnés. On a alors}
\[
(\Phi \cap \Phi')_* \ = \quad
\begin{array}{l}
\textit{partie close à droite de } \Fl(\M) \ \textit{``engendrée par } \Phi_*,\ \Phi'_*\textit{'',} \\
\text{i.e. \texttt{[la]} \textit{plus petite partie close à droite contenant} } \Phi_*, \\
\Phi'_*, \text{ i.e. } \overline{\Phi_* \cup \Phi'_*}
\end{array}
\]
\texttt{[pour} $\Phi$\texttt{,} $\Phi'$ \texttt{parties closes à gauche]}\textit{, plus généralement}
\[
(\bigcap_i \Phi_i)_* \ = \ \overline{\bigcup_i (\Phi_{i\,*})}
\]
\textit{pour toute famille $(\Phi_i)_{i \in I}$ de parties closes à gauche de} $\Fl(\M)$\textit{. Formules duales} \texttt{[pour des parties closes à droite]}
\[
\left\{
\begin{aligned}
(\Psi \cap \Psi')^* \ &= \ \widetilde{\Psi^* \cup \Psi'^*} \\
(\bigcap_j \Psi_j)^* \ &= \ \widetilde{\bigcup_j (\Psi_j^*)}\ .
\end{aligned}
\right.
\]

\item[d)] \textit{Si} $\Phi, \Psi \subset \Fl(\M)$\textit{, alors on a}
\[
\Phi \subset \Psi^* \iff \Psi \subset \Phi_*, \qquad\quad \textit{on écrit alors } \Phi \longleftrightarrow \Psi,
\]
\textit{cela signifie que $\Phi$ a la LLP pour $\Psi$, ou que $\Psi$ a la RLP pour $\Phi$.}
\end{enumerate}

\printedpage{3}
% head: A. Grothendieck : Dérivateurs, Ch. XIII, Édition des Universités de Montpellier II et Paris VII
% foot: Version 25 janvier 2004

\msp{page 4}

\dgrblock{dgrb667}{Closure, transfinite lifting and the retract argument}
\begin{dcomment}{exposition}{p3}{Closure, transfinite lifting and the retract argument}
Grothendieck wants the formal part of Quillen's method to be available before choosing any examples. Proposition~1 provides exactly this preparation. To see its content directly, let $L(R)={}^{\square}R$ and $R(L)=L^{\square}$. Antitonicity follows by adding lifting problems. Every arrow of $L$ lifts against every arrow of $R(L)$, so $L\subseteq L(R(L))$. Applying antitonicity gives $R(L(R(L)))\subseteq R(L)$, while the reverse inclusion is the same elementary observation with the two roles exchanged. Thus $RLR=R$ and $LRL=L$. These identities prove idempotence, describe the fixed points and show that the two fixed-point posets are oppositely isomorphic. For instance,
\[
 R\left(\bigcap_iL_i\right)
   =\overline{\bigcup_iR(L_i)}
\]
when each $L_i$ is left lifting closed. The hypothesis matters: the right side is the join in the lattice of right lifting closed classes, not their ordinary union.

The stability assertions in Proposition~2 express how to transport a chosen solution. Isomorphisms lift by inversion. A lifting problem for a composite is solved successively. A lifting problem for a pushout is equivalent, by its universal property, to one for the original map. For a continuous ordinal chain, begin with the prescribed map on its initial object, extend at each successor, and take the unique induced map at each limit. Transfinite recursion then supplies the diagonal against the composite. The existence of the colimits appearing in the chain is the only colimit hypothesis used here; no presentability has yet entered.

The retract assertion deserves an explicit formulation, since it is the bridge from a constructed factorization to a lifting characterization. Suppose $f:A\to B$ factors as $A\xrightarrow{i}E\xrightarrow{p}B$, and suppose $f\mathrel{\square}p$. Apply this lifting property to the square with top $i$ and bottom $1_B$. The resulting $s:B\to E$ satisfies $sf=i$ and $ps=1_B$. In the arrow category, the morphisms $(1_A,s):f\to i$ and $(1_A,p):i\to f$ have composite $1_f$. Hence $f$ is a retract of $i$. In particular, a factorizing pair $(L,R)$ with $L\mathrel{\square}R$ becomes a weak factorization system as soon as both classes are retract closed. This is the content of the standard retract argument; compare \citet[Lemmas~11.1.4 and~11.2.3]{Riehl2014}. It is this system of two classes of arrows that the manuscript calls a \emph{couple de Quillen}. A modern \emph{Quillen adjunction} is a different notion, involving two model categories and an adjoint pair of functors.
\end{dcomment}


\noindent\textbf{Proposition 2.} \textit{Soit $\Phi$ une partie close à gauche de} $\Fl(\M)$\textit{. Alors $\Phi$ satisfait les conditions} a) b) c) \textit{ci-dessous :}

\begin{enumerate}[label={}]
\item[a)] \textit{$\Phi$ contient les isomorphismes, et est stable par composition, et par cochangement de base (}i.e. \textit{si on a}
$\vcenter{\hbox{\begin{tikzcd}
A \arrow[r, "i"] \arrow[d] & B \arrow[d] \\
A' \arrow[r, "i'"] & B'
\end{tikzcd}}}$
\textit{carré cocartésien dans $\M$, avec $i \in \Phi$, on a $i' \in \Phi$).}

\item[b)] \textit{$\Phi$ stable par facteurs directs (en tant que partie de l'ensemble des objets de} $\underline{\Fl}(\M)$\textit{).}

\item[c)] \textit{$\Phi$ ``stable par limites inductives ordinales'', dans le sens suivant : Soit $I$ un ensemble bien ordonné, $(A_i)_{i \in I}$ un système inductif dans $\M$, tel que les $u_{ji} : A_i \longrightarrow A_j$ soient dans $\Phi$, et $A_{i_0} \xleftarrow{\ \sim\ } \varinjlim\limits_{i<i_0} A_i$ si $i_0 \in I$ est un élément limite. Supposons que $A = \varinjlim\limits_{I} A_i$ existe. Alors $\forall\ i_0 \in I$, la flèche canonique $A_{i_0} \longrightarrow A$ est dans $\Phi$.}
\end{enumerate}

Il y a un énoncé dual pour une partie $\Psi$ close à droite; elle contient les isomorphismes, est stable par composition et par changement de base, par facteurs directs, et enfin stable par limites projectives ordinales.

\msp{page 5}

On aura une réciproque partielle dans les catégories accessibles. C'est pour l'énoncé de cette réciproque que j'ai mis à part la condition \textit{b}) de stabilité par facteurs directs, au lieu de la bloquer avec \textit{a}).

\noindent\textbf{Définition 1.}

\begin{enumerate}[label={}]
\item[\textit{a})] (Pour mémoire.) RLP, LLP, la relation $\Phi \longleftrightarrow \Psi$ entre parties de $\Fl(\M)$ (relation non symétrique!).

\item[\textit{b})] (Pour mémoire.) Partie de $\Fl(\M)$ close à gauche, close à droite.

\item[\textit{c})] Soit $\Phi \longleftrightarrow \Psi$. On dit que le couple $(\Phi, \Psi)$ a la \textit{propriété de factorisation de Quillen}, si \textit{toute} flèche $f$ de $\M$ se factorise en
\[
f = pi, \qquad\quad i \in \Phi,\ p \in \Psi.
\]
On dit que $(\Phi, \Psi)$ est un \textit{couple de Quillen} si c'est un couple associé par \texttt{[les]} propriétés LLP et RLP ($\Phi \longleftrightarrow \Psi$), et \texttt{[si]} elles \texttt{[les parties} $\Phi$ \texttt{et} $\Psi$\texttt{]} sont closes (\textit{i.e.} $\Phi = \Psi^*$, $\Psi = \Phi_*$), et si elles ont la propriété de factorisation de Quillen.
\end{enumerate}

J'énonce dès maintenant le résultat technique essentiel, assurant l'existence de beaucoup de couples de Quillen :

\dgrblock{dgrb668}{A weak factorization system in sets}
\begin{dcomment}{examples}{p5:J'énonce dès maintenant}{A weak factorization system in sets}
The elementary example of injections and surjections makes the preceding abstraction concrete. We work in sets with the axiom of choice. Every injection $i:A\hookrightarrow B$ lifts against every surjection $p:X\twoheadrightarrow Y$: retain the given map on $A$, and choose a preimage in $X$ of the prescribed image of each point of $B\setminus A$. Conversely, lifting against $\varnothing\to\{*\}$ characterizes surjections. If a map $i$ has the left lifting property against all surjections, factor it as
\[
 A\xrightarrow{\iota} A\amalg B
     \xrightarrow{[i,1_B]} B.
\]
The second arrow is surjective, so the retract argument makes $i$ a retract of the injection $\iota$, hence an injection. Thus $(\mathrm{Mono},\mathrm{Surj})$ is a weak factorization system, generated by the single arrow $\varnothing\to\{*\}$.

The factorization displayed above is functorial: a square $(u,v):i\to i'$ acts on its middle object by $u\amalg v$. The diagonal in an arbitrary lifting problem need not be unique or functorially chosen. This is why functorial factorization and a specified coherent family of fillers are additional kinds of structure, even though both are useful in later developments. Nor should this factorization be confused with the familiar image factorization, whose left class is surjections and whose right class is injections. The same two classes occur in opposite order, with a different lifting relation. The asymmetry noted in Definition~1(a) is already visible in this example.
\end{dcomment}


% ---- from Der-ChXIII-25-01-04_p006-010.tex
\setcounter{footnote}{4}

\printedpage{4}
% head: A. Grothendieck : Dérivateurs, Ch. XIII, Édition des Universités de Montpellier II et Paris VII
% foot: Version 25 janvier 2004
\msp{page 6}\textbf{Théorème 1.} \textit{Supposons que $\cM$ soit} accessible, \textit{et stable par (petites) limites inductives. (Il doit suffire stabilité par sommes amalgamées, et par $\varinjlim$ filtrantes, voire ordinales \dots).} \textit{Soit $\Phi_0$ une} petite \textit{partie de} $\Fl(\cM)$\textit{, $\Psi=(\Phi_0)_*$, et $\Phi$ une partie de} $\Fl(\cM)$ \textit{telle que}
\[
\Phi_0\ \subset\ \Phi\ \subset\Psi^*\ (=\widetilde{\Phi}_0).
\]
\textit{Alors :}

\begin{itemize}[leftmargin=2.5em]
\item[a)] \textit{Pour que $\Phi$ soit close à gauche,} i.e. $\Phi=\widetilde{\Phi}_0$\textit{, il faut et il suffit qu'elle satisfasse les conditions} a) b) c) \textit{de la proposition 2. Quand il en est ainsi, le couple $(\Phi,\Psi)$ est un couple de Quillen.}

\item[b)] \textit{Supposons seulement que $\Phi$ satisfasse aux conditions} a), c) \textit{de la proposition 2 (à l'exclusion de la condition} b) \textit{des facteurs directs) – donc on ne peut affirmer que $\Phi$ soit close à gauche. Néanmoins, $(\Phi,\underbrace{\Psi}_{=\,\Phi_*})$ satisfait à la propriété de factorisation de Quillen.}
\end{itemize}

\dgrblock{dgrb669}{The small object theorem and its precise input}
\begin{dcomment}{exposition}{p6:on ne peut affirmer que}{The small object theorem and its precise input}
Theorem~1 is sound under its stated cocompleteness and accessibility hypotheses, understood in the usual locally small sense: an accessible cocomplete category is locally presentable. Its proof explains why the small generating family and the separate retract condition have been singled out. Let $I=\Phi_0$, choose a regular cardinal $\kappa$ such that the domains of all arrows in $I$ are $\kappa$-presentable, and fix $f:X\to Y$. Start with $X_0=X$. At stage $\alpha$, take every square
\[
\begin{tikzcd}
A\arrow[r]\arrow[d,"j"'] &X_\alpha\arrow[d]\\
B\arrow[r]&Y
\end{tikzcd}
\qquad(j\in I)
\]
and attach one copy of $B$ along its specified map from $A$. This is a set of squares because $I$ is small and the category is locally small. Their coproduct and pushout give $X_\alpha\to X_{\alpha+1}$ over $Y$. At limit stages take colimits. After $\kappa$ stages put $E=\mathop{\rm colim}_{\alpha<\kappa}X_\alpha$.

Every map $A\to E$ with $A$ a chosen domain factors through some $X_\alpha$, since the ordinal $\kappa$ is $\kappa$-filtered. A square $j\to(E\to Y)$ is therefore one of the squares attached at the next stage; its attaching cell supplies a filler into $E$. Hence $E\to Y$ belongs to $I^{\square}=\Psi$. The arrow $X\to E$ is an $I$-cell complex. The construction is functorial if every square is attached, without making choices of representatives depending on $f$.

For the theorem's slightly unusual formulation, coproducts of attaching maps cause no problem. A simultaneous attachment indexed by a set can be replaced, after well ordering that set, by successive pushouts. Thus closure under pushouts and continuous transfinite composition already puts every $I$-cell complex in $\Phi$. This proves part~(b). Moreover
\[
 I\subseteq\Phi\subseteq{}^{\square}(I^{\square})
 \quad\Longrightarrow\quad \Phi^{\square}=I^{\square}.
\]
If $\Phi$ is also retract closed, any $g\in{}^{\square}(I^{\square})$ is a retract of the left factor of its cellular factorization. Consequently $g\in\Phi$ and $\Phi={} ^{\square}(I^{\square})$, proving part~(a). This proof supplies the theorem only stated here; a modern account is \citet[Theorem~12.2.2 and Corollary~12.2.4]{Riehl2014}.

The parenthetical weakening of cocompleteness can also be made precise. For this construction it suffices to have the needed pushouts, ordinal colimits, local smallness and one regular cardinal giving the required smallness of the generating domains relative to the cellular chains. Pushouts together with ordinal colimits construct each set of attachments successively from the existing object $X_\alpha$, so a separately postulated empty coproduct is unnecessary. What cannot be dropped is the smallness condition guaranteeing that a lifting problem has appeared before the construction ends.
\end{dcomment}


\textbf{Remarque.} Dans la démonstration, l'existence de $\Phi_0\subset\Phi$, qui joue le rôle d'un petit ensemble de “générateurs” (à gauche) pour $\Phi$, joue un rôle essentiel. Il y aurait lieu de développer la notion de \textit{parties accessibles}
\msp{page 7}
de $\Fl(\cM)$, qui impliquerait l'existence d'un petit ensemble $\Phi_0\subset\Phi$ tel que $\Phi\subset\widetilde{\Phi}_0$, et qui aurait la propriété que
\[
\begin{aligned}
&\Phi\ \text{accessible}\quad\Longrightarrow\quad\Phi_*\ \text{accessible}\\
&\Psi\ \text{accessible}\quad\Longrightarrow\quad\Psi^*\ \text{accessible.}
\end{aligned}
\]
Tout couple $(\Phi,\Psi)$ de parties closes associées ($\Phi=\Psi^*$, $\Psi=\Phi_*$) et accessible, serait donc un couple de Quillen.

\dgrblock{dgrb670}{Accessibility does not supply generating arrows}
\begin{dcomment}{examples}{p7:serait donc un couple de Quillen.}{Accessibility does not supply generating arrows}
Grothendieck's proposed notion of an accessible class is meant to make the small object theorem reusable without displaying its generators. Modern terminology separates three properties: accessibility of a full subcategory of the arrow category, accessibility of a chosen factorization functor, and generation of its right lifting class by a set of arrows. The following known example shows why these cannot be substituted for one another; see \citet[Introduction and Examples~3.7]{Rosicky2021Remarks}.

Let $\mathcal M=\mathrm{Pos}$, and let $L$ consist of the order embeddings. This is a left lifting closed class. Indeed, a monotone map from a subposet $A\subseteq B$ to the two-element chain $2=\{0<1\}$ extends to $B$ by taking the upward closure of the inverse image of $1$. Thus every order embedding lifts against $2\to1$. Conversely, if $u:A\to B$ is not an order embedding, there exist $a,a'\in A$ with $u(a)\leq u(a')$ but $a\not\leq a'$. This also detects failure of injectivity by choosing the appropriate ordered pair of distinct points with the same image. The characteristic function of the principal upper set $\{x\mid a\leq x\}$ takes $a$ to $1$ and $a'$ to $0$, and cannot extend over $u$. Hence
\[
 L={} ^{\square}\{2\to1\}.
\]
One can even construct the associated factorization. Embed $A$ into the complete lattice $2^{\operatorname{Up}(A)}$ by recording membership in all its upper sets, and factor $f:A\to B$ through its graph in $B\times2^{\operatorname{Up}(A)}$. Projection to $B$ lifts against order embeddings, because a monotone map into a complete lattice extends by taking joins of the images of elements below the new point. Thus $(L,L^{\square})$ is a weak factorization system.

As a full subcategory of $\mathrm{Pos}^{[1]}$, $L$ is accessible and accessibly embedded. Filtered colimits preserve order reflection: any relation between images is witnessed at a finite stage, where it already reflects. Every embedding is a filtered union of its finite subembeddings; such finite arrows are finitely presentable in the arrow category and remain so in this full subcategory.

Nevertheless no set $I\subseteq L$ satisfies $I^{\square}=L^{\square}$. Choose an infinite regular $\kappa$ larger than the cardinalities of all domains and codomains in $I$, and let $P$ be the poset of subsets of $\kappa$ of cardinality less than $\kappa$. A map $h:A\to P$ extends along every $i:A\hookrightarrow B$ in $I$ by
\[
 \bar h(b)=\bigcup_{i(a)\leq b}h(a).
\]
Regularity makes this union an element of $P$, and order reflection gives $\bar h(i(a))=h(a)$. Therefore $P\to1$ belongs to $I^{\square}$. Now embed the discrete poset on $\kappa$ into the poset obtained by adjoining a top element. The map $a\mapsto\{a\}$ cannot extend to that top, since its image would have to contain all of $\kappa$. Thus $P\to1\notin L^{\square}$. The proposed generating-set conclusion fails even for an accessible, left lifting closed class having factorizations. This does not refute the possibility of defining the stronger notion Grothendieck wants; it identifies what such a definition would have to add.
\end{dcomment}


\section*{2\quad Catégories de Quillen faibles, et construction d'icelles (Théorème-scholie conjectural)}

Dans toute la suite, on se donne une \textit{catégorie de modèles}
\[
\left\{
\begin{array}{ll}
\mbox{$(\cM,W)$, où $W\subset\Fl(\cM)$} & \\[3ex]
\mbox{\textit{$W$ faiblement saturée} (\footnotemark) :} &
\begin{array}[c]{@{}l@{\quad}p{6.5cm}@{}}
a) & $W$ contient les isomorphismes.\\[1ex]
b) & Si deux parmi $f$, $g$, $gf$ \texttt{[sont]} dans $W$, le troisième aussi.
\end{array}
\end{array}
\right.
\]
\footnotetext{C'est l'axiome M 5 de \textsc{Quillen}.}

\printedpage{5}
% head: A. Grothendieck : Dérivateurs, Ch. XIII, Édition des Universités de Montpellier II et Paris VII
% foot: Version 25 janvier 2004
On veut décrire des données supplémentaires un peu plus faibles que celles d'une catégorie de modèles de Quillen. On suppose données des parties
\[
C,\ F,\ TC,\ TF\ \subset\ \Fl(\cM),
\]
avec les propriétés suivantes. Comme axiome préliminaire, je pose

\begin{itemize}[leftmargin=3em,labelsep=0.6em]
\item[\textbf{M 0}] Stabilité de $\cM$ par $\varprojlim$ finies et $\varprojlim$ \texttt{[plutôt $\varinjlim$]} finies.

\item[\textbf{M 1}]\msp{page 8}Propriété de relèvement-prolongement
\[
C\longleftrightarrow TF,\qquad TC\longleftrightarrow F.
\]

\item[\textbf{M 2}] Propriété de factorisation pour $(C,TF)$, $(TC,F)$.

\item[\textbf{M 3}] Stabilité de $F$ par composition et par changement de base, de $C$ par composition et cochangement de base, $F$ et $C$ contiennent les isomorphismes.

\item[\textbf{M 4}] Stabilité de $TF$ par composition et changement de base, de $TC$ par composition et cochangement de base, et $TF$ et $TC$ contiennent les isomorphismes.

\item[\textbf{M}$'$] $TC\subset C\cap W$, $TF\subset F\cap W$.

\item[\textbf{M}$''$] $C\subset\ \underbrace{\Cof_W}_{\text{$W$-cofibrations}}\ ,\ F\subset\ \underbrace{\Fib_W}_{\text{$W$-fibrations}}\ .$
\end{itemize}

\texttt{[Un morphisme $p:X\longrightarrow Y$ de $\cM$ est une $W$-fibration si pour tout diagramme de carrés cartésiens de la forme :}
\[
\begin{tikzcd}
X \ar[d,"p"'] & X' \ar[l] \ar[d] & X'' \ar[l,"t"'] \ar[d] \\
Y & Y' \ar[l] & Y'' \ar[l,"s"] \ ,
\end{tikzcd}
\]
\texttt{si $s\in W$, alors $t\in W$. La notion de $W$-cofibration est la notion duale.]}

Dans le cas de \textsc{Quillen} (Q-\textit{catégories de modèles}), M$'$ est renforcé en
\[
TC=C\cap W,\qquad TF=F\cap W,
\]
et il n'y a pas d'axiome M$''$. L'axiome M$''$ équivaut à la propriété de “propriety”, correspondant aux Q-catégories de modèles \textit{propres.}

\dgrblock{dgrb671}{Relative equivalences and properness}
\begin{dcomment}{exposition}{p8:correspondant aux Q-catégories de modèles}{Relative equivalences and properness}
The additional four classes serve two purposes: they furnish replacements and make ordinary base change behave homotopically. In modern language $(\mathcal M,W)$ is first a relative category with two-out-of-three. It is not yet a Quillen model category. Axiom M~2 supplies both factorizations, and M~1 supplies the corresponding lifting relations; without retract closure these pairs need not yet be weak factorization systems. The strict inclusions allowed in M$'$ deliberately leave open whether the chosen trivial cofibrations exhaust $C\cap W$, and similarly for fibrations.

The intrinsic class $\mathrm{Fib}_W$ has a stronger universal quantifier than the assertion that one pullback preserves weak equivalences. A map $p:X\to Y$ belongs to it precisely when, after every base change $Y'\to Y$, the functor
\[
 (p')^*:\mathcal M/Y'\longrightarrow\mathcal M/(X\times_Y Y')
\]
takes weak equivalences to weak equivalences. Its dual is $\mathrm{Cof}_W$. For an ordinary model structure, pullbacks of fibrations are again fibrations. Therefore right properness says exactly that $F\subseteq\mathrm{Fib}_W$; left properness similarly gives $C\subseteq\mathrm{Cof}_W$. The identification in the manuscript is correct, with the customary completeness and retract conventions supplied when translating to a modern closed model category.

There is a useful unconditional inclusion
\[
 W_{\mathrm{univ}}\subseteq\mathrm{Fib}_W\cap W,
 \qquad
 W^{\mathrm{univ}}\subseteq\mathrm{Cof}_W\cap W.
\]
For the first, put a universally weak equivalence vertically in a cartesian square. Both vertical sides are in $W$. If the bottom horizontal arrow is in $W$, two-out-of-three puts the top arrow in $W$; apply the same argument after any further base change. The reverse inclusions require the factorization argument of Proposition~4. This logical dependence explains why Proposition~3 first invokes the displayed equality and then proves the result which justifies it. There is no circularity: Proposition~4 uses factorization and the intrinsic fibrations only, not Proposition~3's saturation construction.
\end{dcomment}


\textbf{Proposition 3.} \textit{Supposons les conditions} M 0 \textit{(pour mémoire),} M 1, M 2, M$'$, M$''$ \textit{satisfaites. Alors en remplaçant $C$, $TC$ par les saturés pour les propriétés de stabilité de contenir les isomorphismes, d'être stable par composition et par cochangement de base, et $F$, $TF$ par celles de contenir les isomorphismes et d'être stable par composition et par changement de base, on trouve $C'$, $TC'$, $F'$, $TF'$ satisfaisant}
\printedpage{6}
% head: A. Grothendieck : Dérivateurs, Ch. XIII, Édition des Universités de Montpellier II et Paris VII
% foot: Version 25 janvier 2004
\msp{page 9}\textit{à tous les axiomes} M 1 à M 4, M$'$, $M''$. % [sic] $M''$ (italic M) -> M$''$ (roman M), the axiom is printed with roman M everywhere else

\textsc{Démonstration}. La proposition 2 nous assure que M 1 continue à être satisfait, et pour M 2 c'est trivial. M 3 et M 4 seront valables par construction. D'autre part, M$''$ reste valable, car $\Cof_W$ a les stabilités requises pour $C'$, et $\Fib_W$ celles requises pour $F'$. Reste à prouver M$'$, \textit{i.e.}
\[
TC'\ \subset\ C'\cap W,\qquad TF'\subset F'\cap W,
\]
et pour ceci il suffit de voir que $C'\cap W$ resp. $F'\cap W$ a les stabilités requises pour $TC'$ resp. $TF'$. Cela résulte du

\textbf{Lemme.} \textit{Soit $C$ une partie de} $\Cof_W$\textit{. Si $C$ contient les isomorphismes, il en est de même de $C\cap W$. Si de plus $C$ est stable par composition resp. par cochangement de base, il en est de même} \texttt{[de]} \textit{$C\cap W$. Énoncé dual pour une partie $F$ de} $\Fib_W$\textit{.}

Cela provient du fait que
\[
\boxed{\ \Cof_W\cap W=W^{\mathrm{univ}}\ (\mbox{\footnotemark})\ }
\stackrel{\text{déf}}{=}
\begin{array}{l}
\text{ensemble des flèches $A\longrightarrow B$ telles que pour tout}\\
A\longrightarrow A',\ A'\longrightarrow B'=A'\vee_A B\ \ \text{est dans $W$.}
\end{array}
\]
\footnotetext{flèches \textit{coüniversellement} dans $W$, ou $W$-\textit{équivalences coüniverselles}.}
est lui-même stable par composition et par cochangement de base, et que l'on a bien sûr, si $C\subset\Cof_W$,
\[
C\cap W\ =\ C\cap(\Cof_W\cap W)\ =\ C\cap W^{\mathrm{univ}}.
\]

\msp{page 10}Donc il reste à prouver le

\textbf{Corollaire 1.} \textit{Supposons qu'il existe des parties $C$, $F$, $TC$, $TF$ $\subset$} $\Fl(\cM)$ \textit{satisfaisant} M 2, M$'$ (\footnote{On n'a pas besoin de tout M$'$, seulement de $TC\subset W$, $TF\subset W$.}), M$''$. \textit{Alors}
\[
\begin{aligned}
\Cof_W\cap W\ &=\ W^{\mathrm{univ}}\\
\Fib_W\cap W\ &=\ W_{\mathrm{univ}}\ \ (\mbox{\footnotemark}).
\end{aligned}
\]
\footnotetext{Flèches $f:X\longrightarrow Y$ \textit{universellement} dans $W$, \textit{i.e.} telles que pour tout $Y'\longrightarrow Y$, $f':X'=X\times_Y Y'\longrightarrow Y'$ soit dans $W$.}

Pour voir ceci, montrons p. ex. le deuxième (le premier s'en déduira par dualité).
\[
\begin{tikzcd}
X \ar[d,"f"'] & X' \ar[l] \ar[d,"f'"] \\
Y & Y' \ar[l]
\end{tikzcd}
\qquad\qquad\qquad
f\in\Fib_W,\ f\in W\Longrightarrow f'\in W.
\]

\printedpage{7}
% head: A. Grothendieck : Dérivateurs, Ch. XIII, Édition des Universités de Montpellier II et Paris VII
% foot: Version 25 janvier 2004
On a mieux :

\textbf{Proposition 4.} \textit{Supposons que toute flèche de $\cM$ se factorise en $pi$, avec $i\in W$ et $p\in$} $\Fib_W$\textit{. Alors, si $X$, $Y$ sont dans} $\underline{\Fib}_W S\subset\cM/S$ \texttt{[sous-catégorie pleine de $\cM/S$ formée des $r:Z\longrightarrow S$, où $r\in\Fib_W$]}\textit{, alors toute $f:X\longrightarrow Y$ dans $\cM/S$ qui est dans} $W_S$ \texttt{[formé des flèches au-dessus de $S$ qui sont dans $W$]} \textit{est dans $W^{\mathrm{u}}_S$,} i.e. \textit{elle reste dans $W$ après tout changement de base $S'\longrightarrow S$.}
\[
\begin{tikzcd}[column sep=small]
X \ar[rr,"f"] \ar[dr,"p"'] & & Y \ar[dl,"q"] \\
& S &
\end{tikzcd}
\qquad\qquad\qquad
p,q\in\Fib_W
\]

\textit{Énoncé dual pour les flèches de} $S\backslash\underline{\Cof}_W\subset S\backslash\cM$ \texttt{[sous-catégorie pleine de $S\backslash\cM$ formée des $i:S\longrightarrow Z$, où $i\in\Cof_W$]}\textit{, une telle flèche qui est dans $W$ est $S$-coüniversellement dans $W$, si toute flèche de $\cM$ se factorise en $qj$, avec $j\in$} $\Cof_W$ \textit{et $q\in W$.}

\msp{page 11}\textsc{Démonstration}.
\[
\begin{tikzcd}[column sep=small]
X \ar[dr,"f\in W"] \ar[dd] & & \overline{X'} \ar[ll] \ar[dr,"\overline{f'}\in W"] \ar[dd] & & X' \ar[ll,"i_X"'] \ar[dr,"f'"] \ar[dd] & \\
& Y \ar[dl] & & \overline{Y'} \ar[ll,crossing over] \ar[dl] & & Y' \ar[ll,"i_Y"',crossing over] \ar[dl] \\
S & & \overline{S'} \ar[ll,"F"] & & S' \ar[ll,"i"',"TC"] &
\end{tikzcd}
\]
On factorise $S'\longrightarrow S$ en $S'\overset{i}{\longrightarrow}\overline{S'}\overset{p}{\longrightarrow}S$, avec $i\in W$, $p\in\Fib_W$, par M 2. Donc $f\in W\Longrightarrow\overline{f'}\in W$. Comme $X,Y\in\Ob\,\underline{\Fib}_W S$, on a $\overline{X'},\overline{Y'}\in\Ob\,\underline{\Fib}_W\overline{S'}$, et comme $i:S'\longrightarrow\overline{S'}$ est dans $W$, donc aussi $i_X$, $i_Y$, on a $f'\in W$ par M 5 (saturation faible de $W$).

\dgrblock{dgrb672}{Why weak equivalences between fibrations survive base change}
\begin{dcomment}{exposition}{p11:par M 5 (saturation faible}{Why weak equivalences between fibrations survive base change}
Proposition~4 turns the factorization machinery into a statement about $W$ alone. Its proof is short, but each use of the intrinsic fibration condition has a different role. Given a weak equivalence $f:X\to Y$ over $S$, with $X\to S$ and $Y\to S$ in $\mathrm{Fib}_W$, factor a proposed base change as
\[
 S'\xrightarrow{i}\overline S\xrightarrow{q}S,
 \qquad i\in W,\quad q\in\mathrm{Fib}_W.
\]
Pulling $f$ back along $q$ gives $\bar f\in W$: it is a pullback of the weak equivalence $f$ along the base change $Y\times_S\overline S\to Y$ of $q$. Next, the maps $X\times_S\overline S\to\overline S$ and $Y\times_S\overline S\to\overline S$ remain intrinsic fibrations. Pulling $i$ back along them gives weak equivalences
\[
 X\times_SS'\longrightarrow X\times_S\overline S,
 \qquad
 Y\times_SS'\longrightarrow Y\times_S\overline S.
\]
The commutative square with these sides and $\bar f$ then places $f\times_SS'$ in $W$ by two-out-of-three. No lifting property, functorial choice of factorization, or retract argument is used.

For Corollary~1 apply this assertion in the slice over the codomain of $f\in\mathrm{Fib}_W\cap W$, comparing $f$ with the identity fibration. The converse inclusion was proved above. Duality gives the cofibration equality. Now Proposition~3 can be completed: adjoining the specified composition and base-change operations preserves lifting by Proposition~2 and preserves available factorizations because the four classes only grow. The intrinsic classes have the same finite closure properties. Finally $C'\cap W=C'\cap W^{\mathrm{univ}}$ is stable under pushout and composition, and therefore contains the closure of $TC$; the fibration side is dual. The strengthening by retracts on manuscript p.~11 is valid if $W$ is retract closed, since pullback and pushout transport retract diagrams to retract diagrams.
\end{dcomment}


\dgrblock{dgrb673}{The factorization hypothesis seen on three objects}
\begin{dcomment}{development}{p11:par M 5 (saturation faible}{The factorization hypothesis seen on three objects}
Grothendieck's intrinsic classes make it possible to isolate the role of factorization without any topology. The following small calculation shows exactly what would go wrong if Corollary~1 were read as a formal property of two-out-of-three alone. We checked the comparisons with relative-category and factorization-system treatments in \citet{Riehl2014,Rosicky2021Remarks}; we did not find this particular three-object calculation there. Its purpose is to clarify the hypothesis, not to claim a new general phenomenon.

\textbf{A weak equivalence which is an intrinsic fibration but is not universal.}\ours
Let $\mathcal M$ be the chain $0<1<2$, regarded as a category, and let $W$ consist of the identities and the arrow $w:0\to2$. Then $W$ has two-out-of-three, $w\in\mathrm{Fib}_W$, but $w\notin W_{\mathrm{univ}}$.

\emph{Proof.} The only composite of two nonidentity arrows is $0\to1\to2$. Only its composite is in $W$, so two-out-of-three imposes no further weak equivalence; composites involving identities satisfy the condition. Pullback in this chain is minimum. Every base change of $w$ is $0\to y$ for $y\in\{0,1,2\}$. Below $y=0$ or $1$, the only weak equivalences are identities. Below $2$, the sole additional test is $0\to2$, whose pullback along $0\to2$ is $1_0$. Thus all base changes of $w$ preserve weak equivalences and $w\in\mathrm{Fib}_W$. Its pullback along $1\to2$ is $0\to1$, which is not in $W$. This proves the assertion.

The missing factorization can also be located. The map $1\to2$ is not in $\mathrm{Fib}_W$, since it pulls $w$ back to $0\to1$. A factorization of $1\to2$ with first factor in $W$ must have first factor $1_1$, leaving precisely this nonfibration as second factor. Thus the hypothesis of Proposition~4 fails for a specified arrow. This finite category has all small limits and colimits; filtered diagrams of arrows have colimits attained in their finite directed image. Hence $W$ is even stable under filtered colimits. Those size and closure properties do not replace factorization in the corollary.\hfill$\square$
\end{dcomment}


\textbf{Commentaire sur Proposition 3.} On peut être encore plus exigeant, en demandant dans M 3, M 4 également la stabilité par facteurs directs, et (quand seuls M 1, M2, M$'$, M$''$ % [sic] M2 -> M 2, printed without the space used everywhere else
sont satisfaits) en saturant $C$, $TC$, $F$, $TF$ également pour cette stabilité supplémentaire. La proposition 2 assure encore que M 1 continue à être vérifié (et M 2, M 3, M 4 trivialement), on voit qu'il en est ainsi encore de M$'$, M$''$ à condition qu'on renforce l'axiome de saturation faible M 5 sur $W$ par
\[
\textit{stabilité de $W$ par facteurs directs,}
\]
ce qui implique en effet que $\Cof_W$, $\Fib_W$ sont stables par facteurs directs, ainsi que $W^{\mathrm{univ}}$ et $W_{\mathrm{univ}}$.

\msp{page 12}Il ne semble pas que les axiomes M 0 à M 4, plus M$'$ et M$''$ suffisent à développer les suites exactes de suspension et de cosuspension. Il faut un petit chouia de plus :

\printedpage{8}
% head: A. Grothendieck : Dérivateurs, Ch. XIII, Édition des Universités de Montpellier II et Paris VII
% foot: Version 25 janvier 2004
\textbf{Définition 2.} J'appelle \textit{pré-Q-catégorie de modèles à gauche,} la donnée de $\cM$ et de $W$, $C$, $F$, $TC$, $TF$ $\subset$ $\Fl(\cM)$, satisfaisant les axiomes M 0 à M 5, M$'$, M$''$, plus la forme renforcée de la première inclusion dans M$'$
\[
\text{M}'_0\ \text{(g)}\qquad\qquad\qquad TC=C\cap W\qquad\quad(\text{non seulement } TC\subset C\cap W).
\]
De même, on définit la \textit{pré-Q-catégorie de modèles à droite,} par la condition
\[
\text{M}'_0\ \text{(d)}\qquad\qquad\qquad\qquad\qquad\qquad TF=F\cap W.
\]

NB Quand les deux conditions sont satisfaites à la fois, alors la donnée de $TC$, $TF$ résulte déjà de celle de $W$, $C$, $F$, et on a simplement une Q-catégorie de modèles “propre” (au sens de \textsc{Quillen}).

\dgrblock{dgrb674}{The purpose of the one-sided equalities}
\begin{dcomment}{exposition}{p12:Quand les deux conditions sont satisfaites à la fois}{The purpose of the one-sided equalities}
The one-sided conditions are intended to recover exactness of homotopy fibres or cofibres while leaving the auxiliary choices flexible. They should not be read as a claim that every relative category has such a structure, or that a left structure and a right structure must use the same four classes. Grothendieck explicitly allows different choices. The assertion about a proper Quillen category when both equalities hold uses Quillen's original conventions; for a modern closed model category one also requires the retract conditions, or replaces the classes by their retract closures when the stated hypotheses allow it.

The historical motivation is coherent with the earlier reduction to categories of fibrant objects. Brown's formulation develops homotopy using weak equivalences, fibrations and path objects, without specifying a full model structure; see \citet[\S1, pp.~420--421]{Brown1973}. Grothendieck is pursuing a further independence statement: the resulting homotopy operations should be attached to $W$, with $(C,TC,F,TF)$ used to prove that they exist and have exactness properties. Later in these pages, the same issue reappears for all diagram categories, where objectwise replacement alone does not automatically construct a coherent derived Kan extension.

The claims on manuscript p.~13 about exactness are explicitly provisional: the source says that they still require checking. They are therefore recorded as conjectural assertions in the audit. The completion proved in the commentary on manuscript p.~32 settles their conclusion under the printed hypotheses: either one-sided equality implies exact factorization of weak equivalences, after which the two designated trivial classes can be enlarged to $C\cap W$ and $F\cap W$. The resulting structure is biexact with the same $W,C,F$, though it need not be closed. Thus one-sided structures retain flexibility in the designated trivial maps, but their existence already yields a possibly nonclosed Quillen structure. This is a later verification of the assertion, not a proof supplied on p.~13. The cross-reference to ``proposition~3'' on p.~13 points in fact to Proposition~4: the property cited there is the universal base-change property just proved, together with Corollary~1.
\end{dcomment}


\textbf{Scholie.} Pour moi, l'introduction de $C$, $TC$, $F$, $TF$ est purement auxiliaire – ce qui m'importe, ce sont les propriétés de $W$. Ainsi, l'existence de $C$, $TC$, $F$, $TF$ satisfaisant aux axiomes faibles M 0 à M 4, M$'$, M$''$, du fait qu'elle donne les factorisations
\msp{page 13}
dont il est question dans \texttt{[la]} proposition 3, entraîne la propriété remarquable de $W$ énoncée dans cette proposition, et dans le corollaire qui la précède. De même, ces propriétés de factorisation (\footnote{sous une forme légèrement renforcée il est vrai – il faut une factorisation \textit{fonctorielle} pour $f$ \dots}) assurent l'existence d'une bonne théorie des fibres et cofibres $W$-homotopiques sans besoin de référence dans cette théorie à $C$, $TC$, $F$, $TF$ (dont l'introduction donne tout au plus des petits compléments \dots). Mais je ne vois pas comment, à partir des seules hypothèses de factorisation, prouver la propriété cruciale d'exactitude pour les carrés $W$-cartésiens, et $W$-cocartésiens, nécessaire pour pouvoir développer la suite \textit{exacte} de cosuspension et de suspension. C'est pour avoir cette exactitude que l'introduction auxiliaire, et transitoire, de l'attirail $(C,TC,F,TF)$ est utile, l'axiome supplémentaire à gauche $TC=C\cap W$ assurant (sauf erreur) l'exactitude des carrés $W$-cartésiens, et l'axiome à droite $TF=F\cap W$ celle des carrés $W$-cocartésiens. (Chose à vérifier avec plus de soin que je ne l'ai fait à présent.) Une fois ceci obtenu, et éventuellement itou pour les $(\cM(I),W(I))$ pour toute catégorie-diagramme $I$, on peut
\msp{page 14}
totalement oublier ces ensembles auxiliaires – j'ai envie en tous cas de les oublier.

\dgrblock{dgrb675}{The auxiliary status of the Quillen classes}
\begin{dcomment}{context}{p14:totalement oublier ces ensembles auxiliaires}{The auxiliary status of the Quillen classes}
Grothendieck wants the eventual constructions to belong to the pair
$(\mathcal M,W)$, even when their existence is established by choosing
cofibrations and fibrations. His wish to forget those choices is therefore a
mathematical requirement: two suitable presentations should give the same
homotopy theory of diagrams, including its Kan extensions and their comparison
maps. It does not assert that the choices are dispensable in a proof.

The letter to Thomason of 2 April 1991 places this passage in its immediate
setting. Grothendieck begins with the accessibility theory he had been
developing from SGA~4, then distinguishes the intrinsic localizer from the
additional classes used to calculate with it. Concerning closed Quillen
triples he writes, ``J'avais cru pouvoir m'en passer, mais finalement n'y suis
pas parvenu'' \citep[pp.~3--4]{Grothendieck1991Letter}. The contrast he draws
there with Brown's one-sided methods explains why the present chapter keeps
trying weaker combinations of axioms. The asymmetry is deliberate: the
existence of homotopy fibres and that of homotopy cofibres need not initially
be established by the same auxiliary structure.

There are two distinct tasks in such a programme. One must first construct
the localized categories of diagrams, with sufficient control over their
morphisms to form derived operations. One must then prove that the operations
satisfy comparison laws as the indexing category varies. An equivalence
between two ordinary homotopy categories settles neither task by itself.
This is why the apparently elementary question whether a lifting problem
can be solved naturally in a diagram will become decisive on manuscript
page~41. A family of componentwise solutions need not be a natural
transformation.

The later literature realizes these two tasks in several settings, with
hypotheses that must be retained. Cisinski's construction of homotopy right
Kan extensions starts from a closed Quillen model category with small
limits; it does not require an injective model structure on every diagram
category \citep[Theorem~6.11]{Cisinski2003}. His subsequent theory of derivable
categories provides a further framework for one-sided and exact-category
examples \citep[Introduction and \S1]{Cisinski2010}. These results give a
precise continuation of the programme, but they should not be substituted
for a proof of the particular conjectural recognition criterion now being
formulated. Here the outstanding issue is the smallness of the classes from
which the factorization systems are to be built.
\end{dcomment}


Mon espoir, c'est que pour une $W$ donnée, même si on n'arrive pas à trouver une structure de triple de Quillen $(W,C,F)$ associée (p. ex. parce qu'il n'y en aurait pas), on peut trouver à la fois des quintuplets $(W,C,F,TC,TF)$ qui soient des pré-Quillen à gauche, et des quintuplets (d'\textit{autres} quintuplets!) qui soient des pré-Quillen à droite.

Voici le genre d'énoncé en forme, pour le moment conjectural voire totalement heuristique, auquel je songe :

\textbf{Théorème-scholie conjectural :} \textit{On suppose que $\cM$ est stable par $\varinjlim$ “quelconques”, et par $\varprojlim$ finies, et que $W$ est stable par $\varinjlim$ filtrantes (en plus de la saturation faible).}

% ---- from Der-ChXIII-25-01-04_p011-015.tex
\printedpage{9}
% head: A. Grothendieck : Dérivateurs, Ch. XIII, Édition des Universités de Montpellier II et Paris VII

\begin{itemize}[leftmargin=2.5em]
\item[\circled{1}] \textit{Alors} $\Cof_W$ \textit{et} $W^{\univ}$ \textit{ont les propriétés de stabilité} a) \textit{et} c) \textit{de la proposition} 2 (p.~4), \textit{et aussi la propriété de stabilité} b) \textit{par facteurs directs, pourvu que} $W$ \textit{lui-même l'a.}

\item[\circled{2}] \textit{Supposons} $\mathcal{M}$ accessible, \textit{de sorte qu'on dispose du théorème} 1. \textit{Conditions équivalentes sur} $W$ \textit{:}

\msp{page 15}

\begin{itemize}[leftmargin=2.5em]
\item[a)] \textit{Il existe une partie petite} $C_0$ \textit{de} $\Cof_W$\textit{, et} $TC_0$ \textit{de} $W^{\univ}$\textit{, telles que}
\[ TF \overset{\text{déf}}{=} C_{0*} \subset W, \qquad (TC_0)_* \overset{\text{déf}}{=} F \subset \Fib_W. \]

\item[a$'$)] \textit{Il existe un quadruplet} $(C,TC,F,TF)$ \textit{de parties de} $\Fl(\mathcal{M})$\textit{, satisfaisant les conditions} M~1 \textit{à} M~4, M$'$, M$''$ \textit{de la page} 6 (\footnote[10]{avec $C$, $TC$ clos à gauche, $F$, $TF$ clos à droite.}).

\item[a$''$)] \textit{Comme} a$'$)\textit{, en supposant} M$'$ \textit{sous forme renforcée avec} M$'_0$ (g), i.e. \textit{on a une structure de pré-Q-catégorie de modèles à gauche (}$TC = C \cap W$\textit{).}

\item[b)] \textit{On a}
\[ (\Cof_W)_* \ \subset\ W, \qquad (W^{\univ})_* \ \subset\ \Fib_W \]
\textit{(donc} $(\Cof_W)_* \subset \Fib_W \cap W$ \textit{puisque} $(\Cof_W)_* \subset (W^{\univ})_*$ (\footnote[11]{$W^{\univ} \subset \Cof_W$.})\textit{).}

\item[b$'$)] \textit{On a}
\[ (\Cof_W)_* \ \subset\ W_{\univ}, \qquad (W^{\univ})_* \ \subset\ \Fib_W. \]
\end{itemize}

\item[\circled{3}] \textit{Pour qu'il existe un quadruplet qui, avec} $W$\textit{, définisse une pré-Q-catégorie de modèles à droite (}$TF = F \cap W$\textit{)} (\footnote[12]{On suppose toujours $\mathcal{M}$ accessible dans 3° et 4°, avec en plus $F$, $TF$ (et $C$, $TC$) closes à droite (resp. à gauche).})\textit{, il faut et il suffit qu'on ait} b$'$) \textit{ci-dessus, ainsi que}
\[ (W^{\univ})_* \cap W_{\univ} \ \ \text{\itshape est clos à droite.} \]

\msp{page 16}

\item[\circled{4}] \textit{Les hypothèses préliminaires sur} $(\mathcal{M}, W)$\textit{, y inclus l'accessibilité, ainsi que les conditions équivalentes dans}\circled{2}\textit{, sont stables par passage de} $(\mathcal{M}, W)$ \textit{à} $(\mathcal{M}(I), W(I))$\textit{. Mais qu'en est-il de la condition dans}\circled{3} \textit{? Est-elle conséquence des conditions dans} \circled{2} \textit{(}i.e. \textit{leurs est-elle équivalente) ?} % [sic] leurs est-elle -> lui est-elle, grammar
\end{itemize}

Je vais esquisser la démonstration du théorème-scholie, pour cerner le travail technique qui reste à faire.

\dgrblock{dgrb676}{The closure part of the conjectural theorem}
\begin{dcomment}{exposition}{p16:Je vais esquisser la démonstration}{The closure part of the conjectural theorem}
The theorem-scholie combines a correct closure calculation with a much more substantial assertion that small generating data exist. It helps to finish the closure calculation before attempting the existence argument. Assume the colimit and weak-equivalence hypotheses stated on manuscript p.~14. Then $W^{\mathrm{univ}}$ and $\mathrm{Cof}_W$ contain isomorphisms and are stable under pushout and composition, by pasting pushout squares and using two-out-of-three where appropriate.

For transfinite composition, let $(A_\alpha)_{\alpha\leq\lambda}$ be continuous. Suppose first that every successor map is a universal weak equivalence. Push the chain out along an arbitrary $A_0\to T$, obtaining $T_\alpha=T\amalg_{A_0}A_\alpha$. Each successor is in $W$. Transfinite induction shows that $T\to T_\beta$ is in $W$: the successor step uses composition, while a limit step is the filtered colimit of the arrows $T\to T_\alpha$. Thus $A_0\to A_\lambda$ lies in $W^{\mathrm{univ}}$.

For $\mathrm{Cof}_W$, start with a weak equivalence $T\to T'$ under $A_0$. The two pushed-out chains satisfy $T_0\to T'_0\in W$. The successor step preserves this relation because the corresponding base change of $A_\alpha\to A_{\alpha+1}$ is a $W$-cofibration; at a limit step use filtered-colimit closure of $W$. It follows that pushing out along $A_0\to A_\lambda$ preserves weak equivalences. Repeating after arbitrary cobase change establishes the universal quantifier in the definition of $\mathrm{Cof}_W$.

Retract closure is available here without a further assumption on $W$. If $f$ is a retract of $g\in W$ in the arrow category, write the retract maps as $s:f\to g$ and $r:g\to f$, with $rs=1_f$. Then $f$ is the colimit of
\[
 g\xrightarrow{sr}g\xrightarrow{sr}g\xrightarrow{sr}\cdots.
\]
Indeed a compatible cocone $h:g\to z$ satisfies $h(sr)=h$ and factors uniquely as $hs:f\to z$ followed by $r$. Since this is a filtered diagram of arrows in $W$, its colimit $f$ belongs to $W$. Pullback and pushout preserve retract diagrams, so the intrinsic and universal classes inherit retract closure. This is the elementary mechanism behind \citet[Proposition~1.19]{Beke2000}; the stronger accessibility hypothesis is not needed when filtered-colimit closure is already assumed.

These arguments establish item~\circled{1}. They do not establish that either whole class is defined by a lifting property, nor do they produce a small generating family. The distinction becomes decisive in the next paragraph of the manuscript.
\end{dcomment}


1°) Est de la vérification de routine.

2°) Voici les implication $\pm$ tautologiques entre les conditions envisagées : % [sic] les implication -> les implications, agreement
\[ a'' \Longrightarrow a' \overset{?}{\Longrightarrow} a \Longrightarrow b \Longleftarrow b', \tag*{$(*)$} \]
% foot: Version 25 janvier 2004
\printedpage{10}
% head: A. Grothendieck : Dérivateurs, Ch. XIII, Édition des Universités de Montpellier II et Paris VII
mais je vois que $a' \Longrightarrow a$ n'est nullement tautologique, et peut-être même faux après tout~! (\footnote[13]{C'est pourquoi je l'ai modifié, en renforçant $a'$) par \texttt{[la condition de la note 10]}.}). Il faudrait supposer que dans \texttt{[le]} quadruplet $(C,F,TC,TF)$, $F$ et $TF$ sont clos à droite, et plus précisément, que $F = (TC)_*$, $TF = C_*$, ce qui sera le cas p.~ex. si on a une structure de Q-catégorie de modèles “close” (“closed model category”), au sens de \textsc{Quillen}. Les autres trois implications dans $(*)$ sont bel et bien tautologiques.

$a \Longrightarrow a'$. On peut dans $a$) supposer
\[ TC_0 \ \subset\ C_0 \]
(quitte à remplacer $C_0$ par $C_0 \cup TC_0$), puis, grâce à
\msp{page 17}
1°) ($W^{\univ}$ et $\Cof_W$ sont clos à gauche), posant
\[ \overline{TC_0} = TC, \qquad \overline{C_0} = C, \qquad \text{\texttt{[plutôt}}\ \ \widetilde{TC_0} = TC, \qquad \widetilde{C_0} = C, \text{\texttt{]}} \]
on aura $TC \subset W^{\univ}$, $C \subset \Cof_W$. Passant de $TC_0$ à $TC$, de $C_0$ à $C$, on a perdu seulement que ces parties de $\Fl(\mathcal{M})$ soient petites. Mais par le théorème 1, posant
\[ F \ =\ (TC)_*, \qquad TF \ =\ C_*, \]
les couples de parties closes (à gauche -- à droite)
\[ (C,\, TF), \qquad (TC,\, F) \]
satisfont à la condition de factorisation de Quillen. On aura alors les propriétés M~1 à M~4, plus les inclusions $TC \ \subset\ C \cap W$ (car $TC \ \subset\ W^{\univ} \ \subset\ W$, et $TC \ \subset\ C$ par construction), et $C \subset \Cof_W$ de M$'$ et M$''$. Il reste à vérifier $TF \subset F \cap W$, $F \subset \Fib_W$, mais $F$ $(= (TC)_*)$ $\ \subset \Fib_W$ par hypothèse, et $TF$ $(= C_*)$ $\ \subset W$ aussi~; d'autre part $TF \subset F$ résulte de $TC \subset C$.

\dgrblock{dgrb677}{From generators to the two factorizations}
\begin{dcomment}{exposition}{p17:par construction), et}{From generators to the two factorizations}
The implication $(a)\Rightarrow(a')$ can be proved completely without the unsupported assertion that $\mathrm{Cof}_W$ and $W^{\mathrm{univ}}$ themselves are left lifting closed. Replace $C_0$ by $C_0\cup TC_0$; this preserves the required inclusions and makes $TC_0\subseteq C_0$. Put
\[
 C=\operatorname{cof}(C_0),\qquad
 TC=\operatorname{cof}(TC_0),\qquad
 TF=C_0^{\square},\qquad F=TC_0^{\square},
\]
where $\operatorname{cof}$ means pushouts, transfinite composites and retracts. The small object theorem identifies these left classes with the corresponding double lifting closures and supplies both factorizations. The closure calculation in the preceding block gives
\[
 C\subseteq\mathrm{Cof}_W,\qquad
 TC\subseteq W^{\mathrm{univ}},\qquad TC\subseteq C.
\]
Condition $(a)$ supplies $TF\subseteq W$ and $F\subseteq\mathrm{Fib}_W$; antitonicity gives $TF\subseteq F$. All of M~1--M~4, M$'$ and M$''$ follow.

This repairs the proof, not by assuming the desired global closure statement, but by applying the theorem to the actual small sets that condition $(a)$ provides. The reverse implication remains a separate issue: a right lifting closed class need not be generated by a set on the left. The order-embedding example above explains why the footnote adding lifting closedness to $(a')$ does not, by itself, supply that missing smallness. Thus $(a'')\Rightarrow(a')$, $(a)\Rightarrow(b)$ and $(b')\Rightarrow(b)$ are immediate, and $(a)\Rightarrow(a')$ is now proved; the remaining arrows in the proposed equivalence require additional arguments.
\end{dcomment}


$a' \Longrightarrow a''$. On veut donc, en plus des conditions M~1 à M~4 et M$'$, M$''$, la condition
\[ TC \ =\ C \cap W\ (= C \cap W^{\univ}). \]
L'idée, c'est de remplacer le couple
\msp{page 18}
\[ C,\ TC \qquad (C \subset \Cof_W,\ TC \subset C \cap W), \]
donnant naissance à
\[ F = (TC)_*, \qquad TF = C_*, \]
par
\[ C,\ C \cap W\ (\supset TC), \]
% foot: Version 25 janvier 2004
\printedpage{11}
% head: A. Grothendieck : Dérivateurs, Ch. XIII, Édition des Universités de Montpellier II et Paris VII
ce qui ne change pas le couple $(C, TF)$, mais bien le couple $(TC, F)$, remplacé par
\[ \underbrace{(TC)'}_{\supset TC} \ =\ \underbrace{C \cap W}_{= C\cap W^{\univ}}, \qquad F' \ =\ \underbrace{(TC')_*}_{\subset F} \]
(on a remplacé $F$ par un $F'$ plus petit) (\footnote[14]{NB $TC'$ sera encore clos à gauche comme intersection de deux parties closes à gauche.}). Le seul problème, c'est si $(TC', F')$ satisfait encore à la condition de factorisation de Quillen. Ce sont des parties closes (à gauche pour $TC'$, à droite pour $F'$) “duales” l'une de l'autre,
\[ TC' = F'^*, \qquad F' = TC'_* \ , \]
mais comme on ne sait pas si $TC'$ est de la forme $\widetilde{TC'_0}$, avec $TC'_0$ petite, on n'est pas sûr que la conclusion du théorème 1 s'applique. Mais si on suppose dans $a'$) que $C$ est une partie \textit{accessible} de $\Fl(\mathcal{M})$ (\textit{i.e.}, étant déjà close à gauche, que $C = \widetilde{C_0}$, $C_0$ petit),
\msp{page 19}
alors je présume qu'on pourra prouver qu'il en est de même de $C \cap W$ ($W$ elle-même étant accessible, grâce à la condition de stabilité par petites $\varinjlim$ filtrantes), donc le théorème 1 s'applique bel et bien à $TC' = C \cap W$. Mais si on garde $a'$) telle quelle, alors il faudrait sans doute prouver que $a'$) \textit{implique} $a'$) sous la forme plus forte, avec $C$ \textit{accessible.} Mais comme on a $a' \Longrightarrow a$, et qu'on vient d'établir que $a \Longrightarrow a'$ sous cette forme plus forte, on gagne (modulo le fait conjectural admis sur les parties admissibles \texttt{[plutôt accessibles]} de $\Fl(\mathcal{M})$).

\dgrblock{dgrb678}{The solution-set repair and Smith's theorem}
\begin{dcomment}{later}{p19:donc le théorème 1 s'applique}{The solution-set repair and Smith's theorem}
Grothendieck is trying to enlarge $TC$ to $C\cap W$ while keeping a factorization theorem. The successful modern recognition theorem places the smallness assumption on the weak equivalences, where it can be checked independently. Smith's theorem, published with proof by \citet[Theorem~1.7]{Beke2000}, assumes a locally presentable category, a set $I$ of arrows, two-out-of-three and retract closure for $W$, the inclusion $I^{\square}\subseteq W$, closure of $\operatorname{cof}(I)\cap W$ under pushouts and transfinite composition, and a solution-set condition for maps from each $i\in I$ into $W$ in the arrow category. It then gives a cofibrantly generated model structure with cofibrations $\operatorname{cof}(I)$ and fibrations $(\operatorname{cof}(I)\cap W)^{\square}$. Accessibility and accessible embeddedness of the full weak-equivalence subcategory imply the required solution-set condition; this is Beke's Proposition~1.15.

Here is the exact consequence for these pages. Retain their hypotheses and condition $(a)$, and assume in addition that the full subcategory of weak equivalences is accessible and accessibly embedded in $\mathcal M^{[1]}$. Use $I=C_0\cup TC_0$ and the two factorizations constructed in the preceding block. Corollary~1 gives $\mathrm{Cof}_W\cap W=W^{\mathrm{univ}}$. Hence
\[
 C\cap W=C\cap W^{\mathrm{univ}}
\]
is closed under pushouts, transfinite composites and retracts. Smith's hypotheses hold, so $C\cap W$ is generated by a set and the desired replacement of $TC$ has factorizations. In fact the result is stronger than the proposed left pre-Q structure: it is a full combinatorial model structure. It is proper because $C\subseteq\mathrm{Cof}_W$ and
\[
 (C\cap W)^{\square}\subseteq TC_0^{\square}
       \subseteq\mathrm{Fib}_W.
\]
The first inclusion uses $TC_0\subseteq C\cap W$.

This application is a repair under an explicit additional hypothesis, not a proof of the original theorem-scholie. Closure under filtered colimits says that colimits of diagrams already in $W$ stay in $W$; accessibility additionally supplies a set of sufficiently presentable arrows from which all of $W$ is built. Nor is the manuscript's parenthetical definition $C=\widetilde C_0$ with $C_0$ small equivalent to accessibility of the full arrow subcategory. The example in $\mathrm{Pos}$ separates those notions even when $C$ is left lifting closed.
\end{dcomment}


Il reste donc à prouver
\[ b \ \Longrightarrow\ a, \qquad b \ \Longrightarrow\ b', \]
et il suffit de prouver la première implication, car on sait que $a$) $\Longrightarrow$ $a'$), et que $a'$) implique que
\[ \Fib_W \cap W \ =\ W_{\univ} \]
(cor.~1, page 10), donc $(\Cof_W)_* \subset W$ et $(W^{\univ})_* \subset \Fib_W$ (\textit{i.e.} $b$)) impliquent, comme $(\Cof_W)_* \subset (W^{\univ})_*$ (car $W^{\univ} \subset \Cof_W$), $(\Cof_W)_* \subset W \cap \Fib_W = W_{\univ}$, ce qui est la partie à établir de $b'$) (l'autre $(W^{\univ})_* \subset \Fib_W$ étant tautologiquement contenue dans $b$).

\msp{page 20}

Notons maintenant que $b$) n'est autre que la variante affaiblie de $a$), quand on omet de spécifier que $C_0$, $TC_0$ doivent être petits. Donc le seul problème, pour $b \Longrightarrow a$, c'est de voir ceci :

\textbf{Lemme conjectural} (\footnote[15]{Ce lemme a l'air un peu faux, mais ça doit pouvoir s'arranger \dots}).

\begin{itemize}[leftmargin=2.5em]
\item[\circled{1}] \textit{Si on a une partie close à gauche} $\Phi$ \textit{de} $\Fl(\mathcal{M})$\textit{, telle que} $\Phi_* \subset \Fib_W$ \textit{(resp.} $\Phi_* \subset W$\textit{), alors il existe une partie} petite $\Phi_0 \subset \Phi$\textit{, telle que} $\Phi_{0*} \subset \Fib_W$ \textit{(resp.} $\Phi_0 \subset W$\textit{).}

\item[\circled{2}] \textit{Si} $C$\textit{,} $C'$ \textit{sont deux parties closes à gauche de} $\Fl(\mathcal{M})$\textit{, si} $C$ \textit{et} $C'$ \textit{sont accessibles, il en est de même de} $C \cap C'$\textit{.}
\end{itemize}
% foot: Version 25 janvier 2004

\printedpage{12}
% head: A. Grothendieck : Dérivateurs, Ch. XIII, Édition des Universités de Montpellier II et Paris VII
Je crois que c'est bel et bien vrai sous les hypothèses faites sur $\mathcal{M}$, $W$ -- qui doivent impliquer que dans un sens convenable, $W$ est une partie accessible de l'ensemble $\Fl(\mathcal{M})$ des flèches de la catégorie accessible $\mathcal{M}$. Il doit en résulter qu'une partie de $\mathcal{M}$ construite explicitement en termes de $W$, telle $\Fib_W$, est elle-même accessible, et le lemme doit sortir de là.

Ainsi, 2° est prouvé, modulo ce lemme d'accessibilité, qu'il me faudra prouver $\pm$ en même temps que le théorème 1 (que je n'ai fait ici qu'énoncer).

\msp{page 21}

\dgrblock{dgrb679}{The intersection lemma has a later affirmative answer}
\begin{dcomment}{later}{p20}{The intersection lemma has a later affirmative answer}
The second part of the conjectural lemma has a precise modern resolution when its word ``accessible'' is read in the sense specified on manuscript p.~18, namely set-generated under left lifting closure. In a locally presentable category, if
\[
 C=\operatorname{cof}(I),\qquad C'=\operatorname{cof}(J)
\]
for sets $I,J$, then $C\cap C'$ is again $\operatorname{cof}(K)$ for a set $K$. This follows from \citet[Theorem~3.2 and Remark~3.8]{MakkaiRosicky2014Cellular}: combinatorial cellular structures are stable under the relevant pseudopullbacks. Apply their result to the two identity functors from $(\mathcal M,C)$ and $(\mathcal M,C')$ into $\mathcal M$ with every arrow cellular. The resulting structure has exactly the intersection as its cellular class. This resolves the intersection assertion; it does not imply that the entire intrinsic class $W^{\mathrm{univ}}$ has such generators.

The first part of the lemma asks for a different kind of compactness: from $\Phi^{\square}\subseteq T$, where $T$ is $W$ or $\mathrm{Fib}_W$, it seeks a set $\Phi_0\subseteq\Phi$ with $\Phi_0^{\square}\subseteq T$. Neither the Galois connection nor the preceding intersection theorem provides this reduction. The final parenthesis in the edition has $\Phi_0\subseteq W$; its intended expression is $\Phi_{0*}\subseteq W$, as the surrounding reduction from $(b)$ to $(a)$ requires. The printed slip is retained and recorded separately.

We found the applicable positive results in the small object theorem, Smith's recognition theorem, and the cellular-category theorem just cited. They do not prove this small-subfamily assertion under precisely the preliminary assumptions of the manuscript. We therefore leave that assertion unresolved in this audit rather than attach a claim of current openness to a loosely specified accessibility principle. In a small category the selection is immediate by taking $\Phi_0=\Phi$; its entire content here is the passage from a proper class to a set. Modern accessible weak factorization systems also require care: accessibility refers to a factorization functor, and does not guarantee that its full left or right class is an accessible subcategory. The corrected statements and the distinction from set generation are recorded in \citet[\S3, especially Remarks~3.2, 3.5 and Theorem~3.9]{Rosicky2021Remarks}.
\end{dcomment}


Prouvons 3°. Supposons d'abord qu'il existe un quadruplet “clos” $(C, TC, F, TF)$, s'insérant dans un “pré-quintuplet de Quillen clos” qui soit une Q-catégorie de modèles à droite, \textit{i.e.}
\[ TF \ =\ F \cap W \ =\ F \cap W_{\univ}. \]
Prouvons alors \texttt{[que]}
\[ \underbrace{(W^{\univ})_* \cap W_{\univ}}_{A_W} \qquad \text{est clos à droite.} \tag*{$(*)$} \]
En effet, on a
\[ TC \ \subset\ W^{\univ}, \]
d'où
\[ (W^{\univ})_* \ \subset\ (TC)_* \ =\ F, \]
donc
\[ A = (W^{\univ})_* \cap W_{\univ} \ \subset\ F \cap W_{\univ} \ \overset{\text{(par hyp.)}}{=}\ TF, \]
or $TF$ est clos à droite et contenu dans $W_{\univ}$, d'où
\[ \overline{A} \ \subset\ TF \ \subset\ W_{\univ}, \]
et comme $(W^{\univ})_*$ est clos à droite, on a aussi
\[ \overline{A} \ \subset\ (W^{\univ})_*, \]
d'où $\overline{A} \subset W_{\univ} \cap (W^{\univ})_* = A$, donc $A$ est clos à droite, q.e.d.

Prouvons la réciproque, \textit{i.e.} supposons satisfaites les conditions de 2°, \textit{plus} la condition $(*)$ ci-dessus. On doit construire le quadruplet.

\msp{page 22}

\dgrblock{dgrb680}{A closure obstruction before the existence problem}
\begin{dcomment}{exposition}{p21}{A closure obstruction before the existence problem}
Grothendieck is trying to extract a property of the weak equivalences which is independent of the auxiliary Quillen structure. The class
\[
 A_W=(W^{\univ})_*\cap W_{\univ}
\]
is designed for this purpose: its first condition is a lifting condition, its second says that every pullback is a weak equivalence. There is a useful separation here between a necessary condition, which is established without an accessibility argument, and a construction of factorizations, which still needs one.

Suppose that a closed weak structure is right exact. Since $TC\subset W^{\univ}$, we have $(W^{\univ})_*\subset F$. Consequently
\[
 A_W= (W^{\univ})_*\cap TF.
\]
Indeed, the forward inclusion uses $TF=F\cap W_{\univ}$, while the reverse inclusion uses $TF\subset W_{\univ}$. This writes $A_W$ as an intersection of two right lifting classes. It is therefore right closed. This short proof is equivalent to the closure computation in the manuscript and shows exactly why no generation hypothesis enters the necessary condition.

The converse is more delicate. Set $L=W^{\univ}$, $F=L_*$, $T=F\cap W_{\univ}$ and $C=T^*$. If $T$ is right closed, then $C_*=T$, and $L\subset C$. These are statements about orthogonality; they do not factor an arbitrary arrow. To obtain the proposed structure, one must additionally know that $(L,L_*)$ and $(T^*,T)$ admit factorizations and that $T^*\subset\Cof_W$. If also $L$ is left closed, these conditions give the two weak factorization systems. Since $L\subset C\cap W$, $T\subset F\cap W$, and $F\subset\Fib_W$, they give a closed weak structure; the right exactness follows from $F\cap W=F\cap W_{\univ}$, using the universal-equivalence criterion already proved on manuscript pp.~10--11.

The sentence on the next page which attributes left lifting closure of $\Cof_W$ and $W^{\univ}$ to part~1 of the conjectural theorem passes over precisely this distinction. Part~1 establishes stability under certain constructions. Left lifting closure is a stronger assertion unless a generation theorem applies. The earlier block \emph{From generators to the two factorizations} supplies the appropriate repair when generators have been specified. We read the diagrams on pp.~22--24 as the formal organization of a proposed construction, with the remaining factorization hypotheses still visible; they are not an additional proof of those hypotheses.
\end{dcomment}


Notons d'abord, pour nous y reconnaître, que les relations $b'$) de 2°), et les relations transposées (compte tenu que $\Cof_W$ et $W^{\univ}$ sont clos à gauche, par 1°), savoir
\[ \left\{ \begin{array}{lllllll}
(\Cof_W)_* & \subset & W_{\univ} & , & (W^{\univ})_* & \subset & \Fib_W \\
(W_{\univ})^* & \subset & \Cof_W & , & (\Fib_W)^* & \subset & W^{\univ},
\end{array} \right. \]
donnant (\footnote[16]{compte tenu de $W^{\univ} \subset \Cof_W$, $W_{\univ} \subset \Fib_W$ et des inclusions transposées.}) les deux diagrammes carrés d'inclusions
% foot: Version 25 janvier 2004
\printedpage{13}
% head: A. Grothendieck : Dérivateurs, Ch. XIII, Édition des Universités de Montpellier II et Paris VII
\[
\begin{tikzcd}[column sep=small, row sep=small]
 & W^{\univ} \arrow[dr, hook] & \\
(\Fib_W)^* \arrow[ur, hook] \arrow[dr, hook] & & \Cof_W \\
 & (W_{\univ})^* \arrow[ur, hook] &
\end{tikzcd}
\qquad\qquad (\text{\footnotemark[17]})
\tag*{(I)}
\]
\footnotetext[17]{\textbf{NB} On a ces carrés d'inclusions, chaque fois qu'on a un système de Quillen $(W, C, F, TC, TF)$ \textit{clos.}}
et
\[
\begin{tikzcd}[column sep=small, row sep=small]
 & (W^{\univ})_* \arrow[dl, hook'] & \\
\Fib_W & & (\Cof_W)_*\ , \arrow[ul, hook'] \arrow[dl, hook'] \\
 & W_{\univ} \arrow[ul, hook', "\alpha"'] &
\end{tikzcd}
\tag*{(II)}
\]
dont le premier est transposé (par $\Phi \longmapsto \Phi^*$) du second, et où les 8 ensembles de flèches envisagés, à l'exception de $\Fib_W$ et $W_{\univ}$ (reliés par $\alpha$ dans le carré II) sont clos (les quatre ensembles de I clos à gauche, et les deux ensembles de II non intéressés par $\alpha$, clos à droite).

\msp{page 23}

Les “pré-quintuplets de Quillen clos”, associés à $W$, correspondent au choix de parties closes à gauche $TC$ et $C$, s'insérant dans le diagramme (I) comme montré dans I$'$ ci-dessous, et donnent lieu par transposition à $F$ et $TF$ comme dans II$'$ ci-dessous,
\[
\begin{tikzcd}[column sep=small, row sep=small]
 & & W^{\univ} \arrow[dr, hook] & \\
 & TC \arrow[ur, hook] \arrow[dr, hook] & & \Cof_W \\
(\Fib_W)^* \arrow[ur, hook] \arrow[dr, hook] & & C \arrow[ur, hook] & \\
 & (W_{\univ})^* \arrow[ur, hook] & &
\end{tikzcd}
\tag*{(I$'$)}
\]
\[
\begin{tikzcd}[column sep=small, row sep=small]
 & & (W^{\univ})_* \arrow[dl, hook'] & \\
 & F \arrow[dl, hook'] & & (\Cof_W)_* \arrow[ul, hook'] \arrow[dl, hook'] \\
\Fib_W & & TF \arrow[ul, hook'] \arrow[dl, hook'] & \\
 & W_{\univ}\ , \arrow[ul, hook'] & &
\end{tikzcd}
\tag*{(II$'$)}
\]
le “hic” étant donc de choisir $TC$ et $C$ assez grandes pour qu'on ait
\[ \underbrace{(TC)_*}_{= F} \ \subset\ \Fib_W, \qquad \underbrace{C_*}_{= TF} \ \subset\ W_{\univ} \qquad \text{--} \]
% foot: Version 25 janvier 2004

% ---- from Der-ChXIII-25-01-04_p016-020.tex
\printedpage{14}
% head: A. Grothendieck : Dérivateurs, Ch. XIII, Édition des Universités de Montpellier II et Paris VII
% foot: Version 25 janvier 2004

\msp{page 24}
\noindent et les inclusions stipulées dans 2$^\circ$, $b'$) assurent qu'on peut bel et bien trouver $TC$ et $C$, quitte à les prendre les plus grandes possibles, $TC = W^{\univ}$, $C = \Cof_W$ -- auquel cas, d'ailleurs, on aura
\[
TC \;=\; C \cap W^{\univ},
\]
\textit{i.e.} on a une pré-Q-catégorie de modèles à gauche. Mais il s'agit d'assurer maintenant, par contre, que
\[
TF \;=\; F \cap W_{\univ},
\]
ce qui implique, on vient de le voir, (en notant $(W^{\univ})_* \subset F$) que
\[
\overline{(W^{\univ})_* \cap W_{\univ}} \;\subset\; \overline{F \cap W_{\univ}} = \overline{TF} = TF \subset W_{\univ}. \tag*{$(*)$}
\]
Inversement, si $(*)$ est satisfait, prenons
\[
F = (W^{\univ})_*
\]
le plus petit possible, \textit{i.e.} $C = W^{\univ}$ \texttt{[plutôt} $TC = W^{\mathtt{univ}}$\texttt{]} le plus grand possible. On doit donc prendre
\[
TF \;=\; F \cap W_{\univ} \;=\; (W^{\univ})_* \cap W_{\univ},
\]
et cet ensemble est bien clos à droite par hypothèse.

\textbf{Corollaire.} \textit{Supposons la condition} b$'$) \textit{du théorème-scholie satisfaite. Les conditions suivantes} a), b) \textit{sur $W$ sont équivalentes (et impliquent la condition de} 3$^\circ$ \textit{du théorème-scholie).}

\msp{page 25}
\begin{itemize}[leftmargin=3em]
\item[\cercle{a}] $(W^{\univ})_* \cap \overline{W_{\univ}} = (\Cof_W)_*$.

\item[\cercle{b}] $\widetilde{W^{\univ} \cup (W_{\univ})}^{*} = \Cof_W$, i.e. $\Cof_W$ \textit{est la partie close à gauche engendrée par} $W^{\univ}$ \textit{et} $(W_{\univ})^*$\textit{,} % doubt resolved: printed tilde ends before the final star; preserve the edition reading, see corrections.csv (PDF leaf 16).
\end{itemize}

\noindent\textit{et ces conditions impliquent}

\begin{itemize}[leftmargin=3em]
\item[\cercle{c}] \textit{qu'il existe une structure de} triple clos de Quillen $(W, C, F)$ \textit{incluant} $W$ \textit{(donc avec} $TC = W \cap C$, $TF = W \cap F$\textit{) (lequel sera clos au sens de Quillen si et seulement si} $W$ \texttt{[est]} \textit{stable par facteurs directs -- car $C$ et $F$ sont déjà stables spontanément).}
\end{itemize}

\noindent Il suffit en effet de prendre alors
\[
TC \;=\; W^{\univ}, \qquad\quad C \;=\; \Cof_W .
\]
Mais je doute que l'existence d'un triple clos de Quillen $(W, C, F)$ suffise à impliquer les conditions \textit{a}), \textit{b}) ci-dessus, ni même la condition plus faible
\[
(W^{\univ})_* \cap W_{\univ} \;=\; (\Cof_W)_*
\]
qui signifie \textit{exactement} que $C = \Cof_W$, $TC = W^{\univ}$, et leurs transposés $TF = (\Cof_W)_*$, $F = (W^{\univ})_*$ forment une catégorie de Quillen close.

\clearpage
\printedpage{15}
% head: A. Grothendieck : Dérivateurs, Ch. XIII, Édition des Universités de Montpellier II et Paris VII
% foot: Version 25 janvier 2004

\msp{page 26}
\noindent Conditions sur $(\mathcal{M}, W)$ avec
\begin{itemize}[leftmargin=3em]
\item[\textit{a})] $\mathcal{M}$ stable par $\varprojlim$ finies et $\varinjlim$ (petites),
\item[\textit{b})] $W$ stable par $\varinjlim$ filtrantes,
\item[\textit{c})] $\mathcal{M}$ accessible.
\end{itemize}

\dgrblock{dgrb681}{The two closures in the maximal construction}
\begin{dcomment}{exposition}{p25}{The two closures in the maximal construction}
The corollary seeks a particularly economical structure: take every universal weak cofibration as a trivial cofibration and every $W$-cofibration as a cofibration. Its strength comes from this maximal choice, and the two different closures in the formulas measure its cost.

For arbitrary classes $L,R$ of arrows, the lifting calculus gives
\[
 \bigl(L\cup R^*\bigr)_*=L_*\cap(R^*)_*
       =L_*\cap\overline R.
\]
Thus the right orthogonal of $W^{\univ}\cup(W_{\univ})^*$ is
$(W^{\univ})_*\cap\overline{W_{\univ}}$, with a bar on the second factor. It is not, without another hypothesis, $(W^{\univ})_*\cap W_{\univ}$. If $\Cof_W$ is left closed, applying the Galois correspondence proves the equivalence
\[
 (W^{\univ})_*\cap\overline{W_{\univ}}=(\Cof_W)_*
 \quad\Longleftrightarrow\quad
 \widetilde{\,W^{\univ}\cup(W_{\univ})^*\,}=\Cof_W.
\]
The superscript star must belong to the second summand inside the tilde. The printed formula on manuscript p.~25 places it outside the visible tilde; the accompanying prose and the summary on p.~26 identify the intended expression. We preserve the edition's formula and record the discrepancy separately.

Under condition $b'$ one has $(\Cof_W)_*\subset W_{\univ}$. Hence the stronger displayed equality implies
\[
 (W^{\univ})_*\cap W_{\univ}=(\Cof_W)_*.
\]
For the inclusion from right to left, use $W^{\univ}\subset\Cof_W$ and $b'$; for the reverse inclusion, enlarge $W_{\univ}$ to its right closure and apply the stronger equality. The latter equality is exactly the compatibility needed by the maximal two pairs, once their factorizations exist. It does not itself supply factorizations.

This explains the hierarchy I--IV on p.~26. At the level of formal implications, IV imposes a closure identity stronger than the identity in III, and a maximal structure in III is a special instance of a structure in II. The alleged equivalences inside I and the construction clauses inside II retain the accessibility and factorization problems from pp.~17--20. The hierarchy is a research programme with several proved implications, not four unconditional recognition theorems. In particular, the correct formula in II(c) is $C=(TF)^*$; the printed subscript star is a variance slip.
\end{dcomment}


\[
\mathrm{I}\ \left\{
\begin{array}{@{}l@{\quad}p{11cm}@{}}
\cercle{a} & existence d'un système de Quillen faible $(W, C, TC, F, TF)$ (clos si on y tient) \\[1ex]
\Updownarrow & \\[1ex]
\cercle{b} & existence d'un système de Quillen faible exact à gauche ($TC = C \cap W$). \\[1ex]
\Updownarrow & \\[1ex]
\cercle{c} & $\underbrace{(\Cof_W)_* \subset W,\ (W^{\univ})_* \subset \Fib_W}_{(*)}$ ($\Longleftrightarrow$ $(\Cof_W)_* \subset W_{\univ}$, $(W^{\univ})_* \subset \Fib_W$).
\end{array}\right.
\]

\noindent$\Uparrow$
\[
\mathrm{II}\ \left\{
\begin{array}{@{}l@{\quad}p{11cm}@{}}
\cercle{a} & existence d'un système de Quillen faible exact à droite ($TF = F \cap W$) \\[1ex]
\Updownarrow & \\[1ex]
\cercle{d} & On a $(*)$, et de plus $(W^{\univ})_* \cap W_{\univ}$ clos à droite. \\[1ex]
\Updownarrow & \\[1ex]
\cercle{b} & existence d'un système de Quillen clos (closed model category), $TC = C \cap W$, $TF = F \cap W$, $F = (TC)_*$, $TF = C_*$, $TC = F^*$, $C = (TF)^*$. \\[1ex]
\Updownarrow & \\[1ex]
\cercle{c} & Posant $TC = W^{\univ}$, $F = (TC)_*$, $TF = F \cap W_{\univ}$, $C = (TF)_*$, on trouve avec $W$ un système de Quillen clos. % [sic] $C = (TF)_*$ -> $C = (TF)^*$, cf. the line b) above where $C = (TF)^*$
\end{array}\right.
\]

\noindent$\Uparrow$
\[
\mathrm{III}\ \left\{
\begin{array}{@{}p{11.5cm}@{}}
$C = \Cof_W$ et $TC = W^{\univ}$ font partie (avec $W$) d'un système de Quillen clos (donc $F = (W^{\univ})_*$, $TF = (\Cof_W)_*$, $W = TF \circ TC$). \\[1ex]
$\Updownarrow$ \\[1ex]
\texttt{[On a]} $(*)$, et de plus $(W^{\univ})_* \cap W_{\univ} = (\Cof_W)_*$.
\end{array}\right.
\]

\noindent$\Uparrow$
\[
\mathrm{IV}\ \left\{
\begin{array}{@{}p{11.5cm}@{}}
On a $(*)$ et $\Cof_W = \widetilde{(W_{\univ})^* \cup W^{\univ}}$\ \ . \\[1ex]
$\Updownarrow$ \\[1ex]
On a $(*)$ et $(W^{\univ})_* \cap \overline{W_{\univ}} = (\Cof_W)_*$.
\end{array}\right.
\]

\clearpage
\printedpage{16}
% head: A. Grothendieck : Dérivateurs, Ch. XIII, Édition des Universités de Montpellier II et Paris VII
% foot: Version 25 janvier 2004

\noindent Cas $\mathcal{M} = \mathrm{Cat}$, $W = \mathbf{W}_\infty$. On sait que II est satisfait, par \textsc{Thomason}. Donc on a
\[
(\Cof_W)_* \subset W_{\univ}, \qquad\quad (W^{\univ})_* \subset \Fib_W, \qquad\quad (W^{\univ})_* \cap W_{\univ} \ \text{ clos à droite.}
\]
J'ai pu prouver directement les deux premières inclusions (\textit{i.e.} I), mais pas que $(W^{\univ})_* \cap W_{\univ}$ soit clos à droite. (Je devrais au moins essayer!) (\footnote[18]{Il suffirait de savoir que $\overline{W_{\univ}} \cap \Fib_W = W_{\univ}$, ce qui impliquerait que III et IV sont équivalents (\footnotemark[19]).}\footnotetext[19]{Mais c'est faux, p. ex. si $X$ \texttt{[est une]} catégorie discrète non vide, $X \longrightarrow e$ est dans $\overline{W_{\univ}} \cap \Fib_W$.}). On n'a pas IV, \textit{i.e.} on a
\[
(W^{\univ})_* \cap \overline{W_{\univ}} \;\neq\; (\Cof_W)_*
\]
car
\[
(W^{\univ})_* \cap \overline{W_{\univ}} \not\subset W_{\univ}
\]
(contient les $X \longrightarrow e$ avec $X$ discrète), donc on a
\[
(W^{\univ})_* \cap \overline{W_{\univ}} \;\;\underset{\neq}{\supset}\;\; \underbrace{(W^{\univ})_* \cap W_{\univ}}_{\text{n'est \textit{pas} l'orthogonal de } W^{\univ} \cup (W_{\univ})^*} \quad . % doubt resolved: printed strict-superset sign confirmed; see corrections.csv (PDF leaf 18).
\]

\msp{page 27}

\dgrblock{dgrb682}{Discrete categories detect the extra closure}
\begin{dcomment}{development}{p26}{Discrete categories detect the extra closure}
Grothendieck's example deserves a proof because it distinguishes two conditions that otherwise look nearly identical. We take $W$ in $\mathrm{Cat}$ to be the functors inducing weak homotopy equivalences of nerves. Thomason's theorem gives a proper model structure with this $W$; its existence does not imply the maximal closure condition IV \citep[Theorem~4.9 and Corollary~5.5]{Thomason1980}.

\noindent\textbf{A discrete obstruction to condition IV.}\ours\quad
For every nonempty discrete category $D$, its terminal map belongs to
\[
 (W^{\univ})_*\cap\overline{W_{\univ}}\cap\Fib_W.
\]
If $D$ has at least two objects, this map does not belong to $W_{\univ}$.

\emph{Proof.} We first examine a functor $i:A\to B$ having the left lifting property against every universally weakly equivalent functor. We claim that $\pi_0(i):\pi_0A\to\pi_0B$ is injective. Suppose that two distinct components of $A$ have images connected in $B$ by a zigzag of length $m$. Choose a finite poset $J$ whose Hasse diagram is an alternating linear zigzag with more than $m$ edges. Its nerve is an interval subdivided into edges, so $J\to e$ is a weak equivalence. It is universally so: its pullback over any category $E$ is the projection $J\times E\to E$, and the nerve preserves products; a product with this finite contractible simplicial set is a weak homotopy equivalence.

Send the first chosen component of $A$ to one endpoint of $J$, the second to the other endpoint, and every other component to either endpoint. This defines a functor $a:A\to J$, constant on components. A lift against $J\to e$ would extend $a$ to $B\to J$. The zigzag of length $m$ in $B$ would then connect the two endpoints of $J$ by at most $m$ edges. Every nonidentity arrow of the alternating poset joins adjacent vertices, so this is impossible. This proves the claim.

Now a functor $A\to D$ is just a function $\pi_0A\to D$. If $i\in(W_{\univ})^*$, injectivity of $\pi_0(i)$ lets us extend this function to $\pi_0B$: use its prescribed values on the image and one fixed object of $D$ on all remaining components. Thus $D\to e$ has the right lifting property against $(W_{\univ})^*$, which is exactly membership in $\overline{W_{\univ}}$.

Every functor in $W^{\univ}$ is in $W$, so induces a bijection on components. The same extension argument, now with a unique possible value on every component, gives $D\to e\in(W^{\univ})_*$. Finally, any base change of $D\to e$ is a coproduct of copies of an identity map. Pulling a weak equivalence back along it gives the corresponding coproduct of that weak equivalence; the nerve and geometric realization verify that the result is again a weak equivalence. Therefore $D\to e\in\Fib_W$. If $|D|\geq2$, the map fails already to induce a bijection on components, so is not in $W$, and a fortiori not in $W_{\univ}$.\hfill$\square$

This also explains the qualification required in the footnote: a nonempty discrete category supplies the displayed membership, but a counterexample requires at least two components. The empty category is not included in the lifting argument because extending a function to a newly appearing component would then be impossible. The proof expands the manuscript's example; we checked Thomason's construction and the cited modern treatments of lifting and localization, without finding this precise zigzag proof there. The marker records a proof supplied for this commentary, not a claim of priority.
\end{dcomment}


\section*{3\quad Catégorie de Quillen (deuxième mouture)}

On se donne toujours un couple
\[
(\mathcal{M}, W), \qquad\quad W \subset \mathrm{Fl}(\mathcal{M}), \qquad\quad W \text{ faiblement saturé,}
\]
$\mathcal{M}$ satisfait M 0 (stabilité par $\varprojlim$ finies et $\varinjlim$ finies).

\noindent\textbf{Définition 1 :} Une \textit{structure de Quillen faible} associée à la catégorie de modèles $(\mathcal{M}, W)$ est la donnée de quatre parties $(C, TC, F, TF)$ de $\mathrm{Fl}(\mathcal{M})$, satisfaisant les conditions suivantes.

\begin{itemize}[leftmargin=3em,labelsep=0.5em,align=left]
\item[\textbf{M 1}] $C \longleftrightarrow TF$, $TC \longleftrightarrow F$ (\textit{i.e.} $TF \subset C_*$, $F \subset (TC)_*$).
\item[\textbf{M 2}] Propriété de factorisation de Quillen pour $(C, TF)$ et $(TC, F)$.
\item[\textbf{M 3}] Stabilité de $C$ par composition et par cochangement de base, de $F$ par composition et changement de base -- $C$ et $F$ contiennent les isomorphismes.
\item[\textbf{M 4}] Stabilité similaire pour $TC$ et $TF$.
\item[$\mathbf{M}'$] $TC \subset C \cap W$, $TF \subset F \cap W$.
\item[$\mathbf{M}''$] $C \subset \Cof_W$, $F \subset \Fib_W$.
\end{itemize}

\noindent On dit que la structure de Quillen est \textit{exacte} si elle satisfait à l'axiome supplémentaire

\clearpage
\printedpage{17}
% head: A. Grothendieck : Dérivateurs, Ch. XIII, Édition des Universités de Montpellier II et Paris VII
% foot: Version 25 janvier 2004

\noindent\textbf{M 2 bis}\qquad $W = TF \circ TC$.

\noindent On dit qu'elle est \textit{exacte à gauche} si elle satisfait l'axiome supplémentaire

\noindent\textbf{M$'_0$ (g)}\qquad $TC = C \cap W$ (au lieu de $TC \subset C \cap W$ seulement),

\noindent\textit{exacte à droite} si elle satisfait

\noindent\textbf{M$'_0$ (d)}\qquad $TF = F \cap W$ (au lieu de $TF \subset F \cap W$ seulement) (\footnote[20]{\textbf{NB} Si $C$, $F$, $W$ \texttt{[sont]} stables par facteurs directs (donc aussi $C \cap W$ et $F \cap W$), alors l'exactitude à gauche $TC = C \cap W$ \textit{équivaut} à l'exactitude à droite $TF = F \cap W$, et implique que $C$, $TC$, $F$, $TF$ sont \textit{clos} (à gauche pour $C$, $TC$, à droite pour $F$, $TF$).}).

\msp{page 28}
\noindent\textbf{Remarque.}
\begin{itemize}[leftmargin=3em]
\item[a)] Une structure de Quillen exacte à gauche ou exacte à droite est exacte.
\item[b)] Si on a une donnée $(C, TC, F, TF)$ satisfaisant M 1, M 2, M$'$, $M''$, quand on sature $C$, $TC$ à gauche et $F$, $TF$ à droite, pour les stabilités précisées dans M 2, M 3, on trouve une structure de Quillen. (Tous les axiomes sont évidents, sauf M$'$, pour lequel je renvoie à p. 9, 10, notamment prop. 4.) Si on avait M 2 bis, cela sera encore vérifié (a fortiori) après saturation, \textit{i.e.} on a une structure de Quillen exacte.
\end{itemize}

\dgrblock{dgrb683}{Exactness belongs to the chosen classes}
\begin{dcomment}{exposition}{p27}{Exactness belongs to the chosen classes}
The second definition rearranges the earlier axioms around the task they have to perform. The inclusions $TC\subset C\cap W$ and $TF\subset F\cap W$ distinguish designated trivial maps from all weak equivalences among the designated cofibrations and fibrations. This distinction is essential before retract closure has been imposed. In modern terminology, the closed case consists of two nested weak factorization systems; Grothendieck's $\text{M}''$ adds properness relative to the specified weak equivalences.

There are three different assertions. Exactness says that every weak equivalence has a factorization using the two designated trivial classes. Left exactness identifies the designated trivial cofibrations with $C\cap W$; right exactness identifies the designated trivial fibrations with $F\cap W$. Either identification implies exactness. For example, if $w=pi$ with $i\in C$ and $p\in TF$, then $w,p\in W$ imply $i\in W$, and left exactness gives $i\in TC$. The dual argument uses the other factorization.

Neither this argument nor the two-out-of-three axiom proves that the original chosen class $TF$ is retract closed. Footnote~20 tries to infer too much from retract closure of $C,F,W$. Its conclusion becomes correct after replacing the missing trivial class by its retract closure, but the chosen class need not already equal that closure. The next example makes the distinction concrete. For closed structures, by contrast, exactness alone forces both equalities by the retract argument developed later in the commentary on manuscript p.~57.
\end{dcomment}


\noindent\textbf{Définition 1 bis.} La structure de Quillen faible est dite \textit{structure de Quillen} (tout court) si elle est exacte à gauche et à droite, ou (comme on dit) “biexacte”, alors $(C, F)$ determine déjà $(C, TC, F, TF)$. (Ça correspond exactement aux “model categories” de \textsc{Quillen} \texttt{[propres]}.) On dit que la structure de Quillen faible est \textit{close} si on a % [sic] determine -> détermine, missing accent
\[
C \clarrow TF\ , \qquad\quad TC \clarrow F\ ,
\]
\textit{i.e.} $TF = C_*$, $C = (TF)^*$ et $F = (TC)_*$, $TC = F^*$ (\footnote[21]{Cela signifie aussi que $C$, $TF$, $TC$, $F$ sont stables par facteurs directs. On dit que la structure de Quillen faible est \textit{close à gauche} (resp. \textit{à droite}) si $C$, $TC$ (resp. $F$, $TF$) sont clos à gauche (resp. à droite).}). Une structure de Quillen faible à la fois biexacte (\textit{i.e.} pas faible du tout!) et close est appelée structure de Quillen \textit{parfaite}. (Ça correspond exactement à la “closed model category” \texttt{[propre]} de \textsc{Quillen}.)

\msp{page 29}
\noindent\textbf{Remarque.} Pour qu'une structure de Quillen faible soit parfaite, il suffit qu'elle soit biexacte et qu'on ait
\[
C = (TF)^*, \qquad\quad F = (TC)_*
\]
(ce qui implique déjà que $TC = F^*$, $TF = C_*$, \textit{i.e.} que $TC$ est clos à gauche et $TF$ clos à droite). C'est prouvé dans Quillen. Ça implique aussi que $W$ est fortement saturé -- c'est le pays de cocagne !

\dgrblock{dgrb684}{A retract condition missing from footnote 20}
\begin{dcomment}{development}{p28}{A retract condition missing from footnote 20}
Grothendieck wants the ordinary model-category identities to reappear as soon as retracts behave well. They do reappear after completing the designated classes, but not necessarily before. The following example isolates the missing step without involving any topological pathology.

\noindent\textbf{Left exactness without right exactness.}\ours\quad
Assume the axiom of choice. In $\mathrm{Set}$ put
\[
 W=\Fl(\mathrm{Set}),\qquad C=TC=\{\text{injections}\},
 \qquad F=\{\text{surjections}\},
\]
and let $TF$ consist of those surjections all of whose fibres are either singletons or infinite sets. These classes satisfy all the weak Quillen axioms on manuscript p.~27. The classes $C,F,W$ are retract closed, and the structure is left exact, but it is not right exact.

\emph{Proof.} Injections have the left lifting property against surjections: on the complement of the injected subset, choose a preimage in each relevant nonempty fibre. Both required lifting properties follow. For an arbitrary map $f:A\to B$, use
\[
 A\ \longrightarrow\ A\amalg(B\times\mathbb N)
 \ \xrightarrow{\ p\ }\ B,
 \qquad p|_A=f,\quad p(b,n)=b.
\]
The first map is injective and every fibre of the second is infinite; when $B$ is empty, necessarily $A$ is empty and the assertion is vacuous. This one factorization serves both required pairs.

Injections contain the isomorphisms and are stable under composition and pushout. Surjections have the corresponding pullback and composition properties. For the special class $TF$, pullback preserves each fibre up to bijection. A fibre of a composite is a disjoint union of nonempty fibres of the first map indexed by a fibre of the second. If the indexing fibre is a singleton, the result is a permitted fibre of the first map; if it is infinite, the disjoint union is infinite. Thus $TF$ also has the required stability. The inclusions $TC\subset C\cap W$ and $TF\subset F\cap W$ hold. Since all maps are weak equivalences, $\Cof_W$ and $\Fib_W$ are both the class of all maps, so $\text{M}''$ holds as well.

The usual elementwise arguments show that injections and surjections are retract closed in the arrow category, while $W$ trivially is. We have $TC=C\cap W$. However, the map from a two-element set to a singleton belongs to $F\cap W$ and does not belong to $TF$. It is a retract of $\mathbb N\to\{*\}$, which does belong to $TF$. This proves both the failure of right exactness and the failure of retract closure of the designated class.\hfill$\square$

The repaired assertion can be stated precisely. Suppose $C,F,W$ are retract closed and $TC=C\cap W$. Then $TC$ is retract closed. The factorization and lifting axioms give $TC={}^\perp F$ and $F=TC^\perp$. If $p\in F\cap W$, factor $p=qi$ with $i\in C$ and $q\in TF$. Two-out-of-three gives $i\in TC$, and lifting $i$ against $p$ exhibits $p$ as a retract of $q$. Conversely every retract of a map in $TF$ lies in $F\cap W$. Hence
\[
 \operatorname{Retr}(TF)=F\cap W.
\]
Replacing $TF$ by this class preserves lifting, factorization and properness and produces a perfect structure. The standard retract argument also gives $C={}^\perp(F\cap W)$ and $C^\perp=F\cap W$. The dual repair applies to a right exact structure. We checked the standard model-category retract arguments and Cisinski's treatments; the particular fibre-cardinality example is supplied here to settle this exact footnote, without a priority claim.
\end{dcomment}


\clearpage
\printedpage{18}
% head: A. Grothendieck : Dérivateurs, Ch. XIII, Édition des Universités de Montpellier II et Paris VII
% foot: Version 25 janvier 2004

\noindent\textbf{Scholie :} Le couple $TC \longleftrightarrow F$, satisfaisant à la condition de factorisation de Quillen, et satisfaisant $TC \subset W$, $F \subset \Fib_W$, permet déjà de développer une théorie satisfaisante des fibres $W$-homotopiques dans $\mathcal{M}$ -- mais sans propriété d'exactitude pour les carrés $W$-cartésiens. Dualement, la donnée de $C \longleftrightarrow TF$, avec factorisation, et $TF \subset W$, $C \subset \mathrm{Cofib}_W$, donne une bonne théorie des cofibres $W$-homotopiques, mais sans propriété d'exactitude pour les carrés $W$-cocartésiens. La structure de Quillen faible réunit ces deux données duales l'une de l'autre, avec en plus le lien supplémentaire entre ces deux donnés, par $TC \subset C$, $TF \subset F$ (stipulé dans $M'$). Cette donnée assure déjà que % [sic] donnés -> données, feminine noun as in "ces deux données duales"
\[
\underbrace{\mathcal{M}_{\mathrm{cf}}}_{\substack{\text{sous-catégorie pleine}\\ \text{de $\mathcal{M}$ formée des}\\ \text{objets fibrants et}\\ \text{cofibrants à la fois}}} W_{\mathrm{cf}}^{-1} \;\xrightarrow{\ \sim\ }\; \mathcal{M}W^{-1} = \mathrm{Ho}_W
\]
\msp{page 30}
est une équivalence, mais non pas que le foncteur localisation
\[
\mathcal{M}_{\mathrm{cf}} \longrightarrow \mathcal{M}_{\mathrm{cf}}W_{\mathrm{cf}}^{-1}
\]
est un simple passage au quotient sur les flèches, \textit{i.e.} que si $X$, $Y$ sont dans $\mathcal{M}_{\mathrm{cf}}$, alors
\[
\mathrm{Hom}_{\mathcal{M}}(X, Y) \;\longrightarrow\; \underbrace{\mathrm{Hom}_W(X, Y)}_{\overset{\text{déf}}{=}\,\mathrm{Hom}_{\mathrm{Ho}_W}(X, Y)} \tag*{$(*)$}
\]
soit surjectif. Ceci est assuré par l'axiome d'exactitude M 2 bis,
\[
W \;=\; TF \circ TC.
\]
Avec ce seul axiome, on n'a pas de caractérisation sympathique de la relation d'équivalence dans $\mathrm{Hom}_{\mathcal{M}}(X, Y)$ définie par la flèche $(*)$ (pour $X$, $Y$ dans $\mathcal{M}_{\mathrm{cf}}$).

\dgrblock{dgrb685}{Replacement controls arrows as well as objects}
\begin{dcomment}{exposition}{p29}{Replacement controls arrows as well as objects}
The passage to fibrant-cofibrant objects is the first point where the lifting axioms recover information which would be lost by merely choosing a weakly equivalent object. Grothendieck's assertion of an equivalence of localizations is correct under the weak axioms. The mechanism is worth making explicit because exactness will be needed only for the next assertion, that ordinary arrows represent all localized arrows.

Choose $q_X:QX\to X$ in $TF$ with $QX$ cofibrant. Given $f:X\to Y$, lift against $q_Y$ to obtain $Qf:QX\to QY$ with $q_YQf=fq_X$. This lift need not be unique as an arrow. It is unique in the localization of the cofibrant objects. Indeed, for two lifts $a,b:QX\to QY$, form $P=QY\times_YQY$. Its projections belong to $TF$ by pullback stability. Choose a cofibrant replacement $q_P:QP\to P$ and lift $(a,b)$ through it. The diagonal $QY\to P$ likewise lifts through $q_P$ because $QY$ is cofibrant. Call this lift $d:QY\to QP$. The two composites
\[
 QY\xrightarrow{d}QP\xrightarrow{\operatorname{pr}_j q_P}QY
\]
are identities. Each $\operatorname{pr}_j q_P$ belongs to $TF$, so $d\in W$. After localization both projections are therefore the inverse of $d$, and $a=b$ there. This uniqueness proves compatibility with identities and composition. Two-out-of-three shows that a weak equivalence $f$ gives a weak equivalence $Qf$. Thus $Q$ defines an inverse, after localization, to the inclusion of the cofibrant objects.

For the second step choose $r_X:X\to RX$ in $TC$, with $RX$ fibrant, when $X$ is cofibrant. Then $RX$ remains cofibrant. An arrow $X\to Y$ extends to $RX\to RY$ by lifting $r_X$ against $RY\to e$. To compare two extensions, form $P=RX\amalg_XRX$. The two inclusions $RX\to P$ are in $TC$. Choose $P\to RP$ in $TC$ with $RP$ fibrant. The fold $P\to RX$ extends to $d:RP\to RX$ because $RX$ is fibrant. If $v_0,v_1:RX\to RP$ are the two induced maps, then $dv_0=dv_1=1$, and $d,v_0,v_1$ are weak equivalences. The two extensions consequently agree after localization, by first extending their common pushout map $P\to RY$ to $RP\to RY$. This gives the inverse equivalence for the inclusion of fibrant-cofibrant objects into cofibrant objects.

No equality $TC=C\cap W$ or $TF=F\cap W$ entered these two arguments. What is indispensable is the ability to compare choices inside the relevant full subcategory. This is the same distinction already emphasized in the commentary on Chapter~VI, manuscript p.~87. For the usual model-category statement, see \citet[Theorem~2.2.15 and Corollary~2.2.18]{Cisinski2019HigherCategories}; the argument above records the weaker hypotheses actually used here.
\end{dcomment}


Si on a exactitude à gauche,
\[
TC = C \cap W \qquad\quad (\Longleftrightarrow\ TC = C \cap W^{\univ}),
\]
alors cette relation d'équivalence est celle engendrée par la relation de \textit{Q-homotopie à droite} (\footnote[22]{Il semblerait que la relation de Q-homotopie à droite est une relation d'équivalence, dès que $Y$ \texttt{[est]} dans $\mathcal{M}_{\mathrm{f}}$ et sans supposer la structure de Quillen faible exacte à droite ou à gauche.}, \footnote[23]{\textbf{NB} D'après \textsc{Baues}, il semblerait que sur $\mathcal{M}_{\mathrm{cf}}$, cette relation d'homotopie à droite équivaut à la relation d'homotopie à gauche, et que celle-ci ne dépend pas du cylindre choisi \ldots}) (impliquée par la Q-homotopie à gauche), et la catégorie quotient ainsi définie
\[
\pi_{\mathrm{d}}\mathcal{M}_{\mathrm{cf}} \;=\; \pi\mathcal{M}_{\mathrm{cf}} \;\simeq\; \mathrm{Ho}_W,
\]
s'envoie dans $\pi_{\mathrm{d}}\mathcal{M}_{\mathrm{c}}$,
\[
\pi\mathcal{M}_{\mathrm{cf}} \longrightarrow \pi_{\mathrm{d}}\mathcal{M}_{\mathrm{c}},
\]
et ce foncteur peut être interprété (sauf erreur)

\dgrblock{dgrb686}{Two path conventions and a nontransitive homotopy}
\begin{dcomment}{examples}{p30:alors cette relation d'équivalence est celle engendrée}{Two path conventions and a nontransitive homotopy}
Footnote~22 requires a distinction which the text does not explicitly make. A diagonal factorization
\[
 Y\xrightarrow{s}PY\xrightarrow{p=(d_0,d_1)}Y\times Y
\]
may be required to have $s\in TC$ and $p\in F$, or only $s\in W$ and $p\in F$. Call these, respectively, a \emph{designated} path object and a \emph{broad} path object. In standard model-category terminology the first is a very good path object and the second a good path object. The raw relation associated with a fixed path object declares $f,g:X\to Y$ related when $(f,g)=pH$ for some $H:X\to PY$. Taking the union of these relations over all broad path objects is another operation; it must not be silently identified with the relation from a designated one.

Every designated path object gives a reflexive and symmetric relation, for every source. Reflexivity uses $sf$. For symmetry, lift $s$ against $p$, with top map $s$ and bottom map $\tau p$, where $\tau$ interchanges the two factors of $Y^2$. The filler $r:PY\to PY$ satisfies $pr=\tau p$, and $rH$ reverses a homotopy. Moreover, if $Y\xrightarrow{s'}P'Y\xrightarrow{p'}Y^2$ is another designated path object, the square with vertical maps $s,p'$ and horizontal maps $s',p$ has a filler $PY\to P'Y$ preserving the endpoints. Interchanging the objects gives the converse comparison. Thus \emph{one fixed designated path object and all designated path objects define the same raw relation}, without fibrancy, cofibrancy or retract-closure assumptions. Transitivity is the additional issue.

\noindent\textbf{A counterexample with every object fibrant and cofibrant.}
The following is a dual application of the graph example of \citet[Definition~3.1, Observation~3.3 and Example~3.6]{KurzRosicky2005}. Let $\mathcal G$ be the category of reflexive directed multigraphs. An object consists of sets of vertices and edges, maps $s,t:E\rightrightarrows V$, and a distinguished loop $e:V\to E$ satisfying $se=te=1_V$. Morphisms preserve these data. This is a functor category, hence has all small limits and colimits. Let
\[
 L=\operatorname{Mono}(\mathcal G),\qquad R=L^{\square}.
\]
Here monomorphisms are injective on both vertices and edges. For precision, $R$ is defined by the literal lifting property, not by an abbreviated assertion about whether two image vertices are adjacent: in a multigraph a lifting problem specifies the target edge as well as its endpoints.

The pair $(L,R)$ is a weak factorization system. One generating map adds a vertex with its distinguished loop to the empty graph. The other adds a directed edge between the two vertices of a discrete two-vertex reflexive graph. In a pushout the two endpoints may coincide, so the latter generator also adds a nondistinguished loop. Every subgraph inclusion is a composite of pushouts of these maps: first adjoin its missing vertices, then its missing edges. Conversely such pushouts, their composites and their retracts are monomorphisms. The finite generating graphs are finitely presentable, so the small object theorem proved earlier gives exactly $(L,L^{\square})$.

Pass to $\mathcal M=\mathcal G^{\mathrm{op}}$ and put
\[
 W=C=\operatorname{Mor}(\mathcal M),\qquad TF=\operatorname{Iso}(\mathcal M),
 \qquad TC=R^{\mathrm{op}},\qquad F=L^{\mathrm{op}}.
\]
Both $(C,TF)$ and $(TC,F)$ are weak factorization systems. In particular their classes contain the isomorphisms and have the required composition, pushout or pullback stability, and even retract stability. The inclusions $TC\subset C\cap W$ and $TF\subset F\cap W$ hold. Since every arrow is weak, the intrinsic classes $\Cof_W$ and $\Fib_W$ also contain every arrow. Thus all the printed weak axioms, including $\text{M}''$, hold.

Every object is cofibrant because $C$ consists of all arrows. Every object is fibrant because its map to the terminal object of $\mathcal M$ is opposite to an inclusion of the empty graph, hence belongs to $L^{\mathrm{op}}$. These observations make the example independent of any failure to replace its source or target.

Let $V$ denote the graph with one vertex and its single distinguished loop. Let $D$ have two vertices $0,1$ and exactly one directed edge for each ordered pair of vertices, with the two diagonal edges distinguished. The inclusion $j:V\amalg V\to D$ belongs to $L$. The unique map $q:D\to V$ belongs to $R$: a graph map from a subgraph to $D$ extends by sending each additional vertex to $0$, and then sending each edge to the unique edge with the prescribed endpoints. This extends every lifting problem against a monomorphism. The equality $qj=\nabla_V$ therefore becomes a designated path factorization in the opposite category,
\[
 Y\xrightarrow{q^{\mathrm{op}}}D
     \xrightarrow{j^{\mathrm{op}}}Y\times Y,
 \qquad Y=V\quad\text{viewed in }\mathcal M.
\]

Now let $X$ be the reflexive graph with vertices $a,b,c$, distinguished loops, and four further edges
\[
 a\rightleftarrows b\rightleftarrows c,
\]
with no edge from $a$ to $c$ or from $c$ to $a$, again viewed as an object of $\mathcal M$. The three vertex maps $V\to X$ in $\mathcal G$ are three arrows $f_a,f_b,f_c:X\to Y$ in $\mathcal M$. A right homotopy through the displayed path object is precisely a graph map $D\to X$ with the specified vertex images. Such maps exist for $(a,b)$ and $(b,c)$, whereas none exists for $(a,c)$. Consequently
\[
 f_a\sim_r f_b,\qquad f_b\sim_r f_c,
 \qquad f_a\not\sim_r f_c.
\]
The comparison argument above shows that allowing every designated path object does not change this calculation. Raw designated right homotopy is not transitive even though every object is fibrant and cofibrant and all four lifting classes are retract closed.

This weak structure is not exact. Indeed $TF\circ TC=TC$, and the opposite of a vertex inclusion $V\to X$ is not in $TC$: the vertex inclusion fails the right lifting property against the map adding a vertex. Thus $TC\ne W$. There is also a useful exact variant: replace $C$ by $TC$ and $TF$ by $F$, retaining $W,TC,F$. This is a perfect model structure with all arrows weak. The same path calculation still fails transitivity for the same source, which is now noncofibrant. To see the latter directly, $X\to V$ in $\mathcal G$ does not lift the edge-adding generator with endpoints mapped to $a,c$. Hence even exactness does not make the raw relation from very good paths transitive for unrestricted sources.

\noindent\textbf{The broad relation for a fibrant target.}
Under the printed weak axioms, the relation defined by existence of a homotopy through \emph{some} broad path object is an equivalence relation whenever $Y$ is fibrant, with arbitrary source $X$. Here is the direct argument. Reflexivity uses the constant section of any designated, hence broad, path object. Symmetry uses the same broad object with its endpoints interchanged. For transitivity choose broad path objects
\[
 Y\xrightarrow{s_P}P\xrightarrow{(d_0,d_1)}Y^2,
 \qquad
 Y\xrightarrow{s_Q}Q\xrightarrow{(e_0,e_1)}Y^2
\]
and homotopies from $f$ to $g$ through $P$ and from $g$ to $h$ through $Q$. Form
\[
 K=P\mathbin{\times_{Y}}Q
 \quad\text{using }d_1\text{ and }e_0,
 \qquad
 t=(s_P,s_Q):Y\longrightarrow K.
\]
Each endpoint map $d_i,e_i$ belongs to $F$, since it is the composite of a path fibration with a projection $Y^2\to Y$, and the latter is a pullback of $Y\to e$. Each endpoint is also in $W$ by two-out-of-three applied to its composite with the constant section. The projection $K\to P$ is the pullback of $e_0\in W$ along $d_1\in F\subset\Fib_W$, so it is in $W$. Its composite with $t$ is $s_P\in W$, and hence $t\in W$.

The map $K\to Y^3$ recording the two outer endpoints and their common middle endpoint is the pullback of
\[
 P\times Q\longrightarrow Y^2\times Y^2
\]
along $(y_0,y_1,y_2)\mapsto(y_0,y_1,y_1,y_2)$. Products and pullbacks of $F$-maps are in $F$, so this map is in $F$. The projection $Y^3\to Y^2$ retaining the outer coordinates is again a pullback of $Y\to e$. Their composite $K\to Y^2$ is therefore in $F$. Thus $Y\xrightarrow{t}K\to Y^2$ is a broad path object. The two given homotopies induce $X\to K$, supplying the homotopy from $f$ to $h$. No exactness or retract closure has entered.

The missing implication for designated paths is now visible: this argument proves $t\in W$, but does not prove $t\in TC$. The graph example shows that the omission cannot be repaired from the weak axioms alone. In that example the broad path $Y\xrightarrow{\Delta}Y^2\xrightarrow{1}Y^2$ makes every pair related, whereas every designated path retains the nontransitive relation just calculated. Consequently footnote~22 has a positive answer for the existential broad convention and a negative answer for the designated convention used in the replacement constructions. Its status is a question of these precise conventions, not an unresolved general homotopy theorem.
\end{dcomment}


% ---- from Der-ChXIII-25-01-04_p021-025.tex
\setcounter{footnote}{23}

\printedpage{19}
% head: A. Grothendieck : Dérivateurs, Ch. XIII, Édition des Universités de Montpellier II et Paris VII
% foot: Version 25 janvier 2004
\msp{page 31}
comme un \textit{adjoint à droite} (pleinement fidèle) du foncteur de localisation
\[
\pi_{\mathrm{d}}\cM_{\mathrm{c}} \longrightarrow (\pi_{\mathrm{d}}\cM_{\mathrm{c}})W_{\mathrm{c}}^{-1} \;\simeq\; \mathrm{Ho}_W,
\]
et $\mathrm{Ho}_W$ se déduit de $\pi_{\mathrm{d}}\cM_{\mathrm{c}}$ par un \textit{calcul de fractions à gauche.} Dualement, si on a exactitude à droite, \textit{i.e.} $TF = F\cap W$ ($\Longleftrightarrow$ $TF = F\cap W_{\mathrm{univ}}$), alors la relation d'équivalence pertinente sur $\mathrm{Hom}_{\cM}(X,Y)$, pour $X$, $Y$ dans $\cM_{\mathrm{cf}}$, est la Q-homotopie à gauche (impliquée par la Q-homotopie à droite), et cette fois la flèche naturelle de $\pi_{\mathrm{g}}\cM_{\mathrm{cf}} = \pi\cM_{\mathrm{cf}} \simeq \mathrm{Ho}_W$ dans $\pi_{\mathrm{d}}\cM_{\mathrm{f}}$ \texttt{[plutôt }$\pi_{\mathtt{g}}\cM_{\mathtt{f}}$\texttt{]},
\[
\underbrace{\pi\cM_{\mathrm{cf}}}_{\simeq\,\mathrm{Ho}_W} \longrightarrow \pi_{\mathrm{d}}\cM_{\mathrm{f}} \qquad \texttt{[plutôt }\pi_{\mathtt{g}}\cM_{\mathtt{f}}\texttt{]}
\]
est un \textit{adjoint à gauche} (pleinement fidèle) du foncteur de localisation
\[
\pi_{\mathrm{d}}\cM_f \longrightarrow \cM_{\mathrm{f}}W_{\mathrm{f}}^{-1} \;\simeq\; \mathrm{Ho}_W \qquad \texttt{[plutôt }\pi_{\mathtt{g}}\cM_f \longrightarrow \cM_{\mathtt{f}}W_{\mathtt{f}}^{-1} \;\simeq\; \mathrm{Ho}_W\texttt{]}, % [sic] \cM_f (italic f, twice) -> \cM_{\mathrm{f}}, elsewhere the index f is upright
\]
et la localisée s'explicite par un calcul de fractions à droite.

\dgrblock{dgrb687}{Exactness and the congruence generated by both homotopies}
\begin{dcomment}{development}{p30}{Exactness and the congruence generated by both homotopies}
Grothendieck has now isolated the hypothesis which allows a weak equivalence between good objects to be inverted by an ordinary arrow. His observation admits a precise completion even before one-sided exactness identifies a preferred homotopy relation.

A chosen cylinder of $X$ is a factorization
\[
 X\amalg X\xrightarrow{(j_0,j_1)}IX\xrightarrow{q}X,
 \qquad (j_0,j_1)\in C,\quad q\in TF.
\]
A chosen path object of $Y$ is a factorization
\[
 Y\xrightarrow{s}PY\xrightarrow{(p_0,p_1)}Y\times Y,
 \qquad s\in TC,\quad(p_0,p_1)\in F.
\]
Left and right Q-homotopies are maps $IX\to Y$ and $X\to PY$ with the prescribed endpoints. We allow all such factorizations in these definitions.

\noindent\textbf{A quotient description under exactness alone.}\ours\quad
Assume the weak Quillen axioms and $W=TF\circ TC$. Let $\equiv$ be the smallest congruence on $\mathcal M_{\mathrm{cf}}$ containing both left and right Q-homotopy. Then
\[
 \mathcal M_{\mathrm{cf}}/{\equiv}\ \simeq\ \mathcal M[W^{-1}].
\]
In particular every morphism between fibrant-cofibrant objects in the localization is represented by an ordinary morphism.

\emph{Proof.} Both kinds of homotopy become equalities after localization: for a cylinder, $qj_0=qj_1=1$ and $q\in W$; for a path object, $p_0s=p_1s=1$ and $s\in W$. It therefore suffices to prove that the quotient inverts every weak equivalence between fibrant-cofibrant objects.

Let $w:X\to Y$ be such a weak equivalence. Exactness gives $w=pi$ with $i:X\to Z$ in $TC$ and $p:Z\to Y$ in $TF$. The intermediate object $Z$ is cofibrant because $TC\subset C$ and $X$ is cofibrant; it is fibrant because $TF\subset F$ and $Y$ is fibrant. Lifting $i$ against $X\to e$ gives $r:Z\to X$ with $ri=1_X$. Lifting $\varnothing\to Y$ against $p$ gives $t:Y\to Z$ with $pt=1_Y$.

To compare $ir$ with $1_Z$, choose a path object $Z\xrightarrow{s}PZ\xrightarrow{v}Z\times Z$. The square with left edge $i$, right edge $v$, top edge $si$ and bottom edge $(ir,1_Z)$ commutes. The lifting axiom $TC\perp F$ supplies a right homotopy from $ir$ to $1_Z$. To compare $tp$ with $1_Z$, choose a cylinder $Z\amalg Z\xrightarrow{j}IZ\xrightarrow{q}Z$. The square with left edge $j$, right edge $p$, top edge $(tp,1_Z)$ and bottom edge $pq$ commutes. The lifting axiom $C\perp TF$ supplies a left homotopy from $tp$ to $1_Z$.

Consequently $i$ and $p$ are invertible in the quotient, with inverses $r$ and $t$, and $w$ has inverse $rt$. The universal property of localization now identifies the quotient with $\mathcal M_{\mathrm{cf}}[W_{\mathrm{cf}}^{-1}]$. The replacement argument in the preceding block identifies that category with $\mathcal M[W^{-1}]$.\hfill$\square$

The congruence formulation deliberately keeps the two homotopies together. Under left exactness a left homotopy implies right homotopy: for a cofibrant source, either cylinder endpoint is in $C\cap W=TC$, so the familiar endpoint lifting square against a path fibration produces the right homotopy. Dually, right exactness converts a right homotopy into a left homotopy. Thus the proposition also proves the manuscript's description by the relation \emph{generated} by right homotopy under left exactness, and its dual. We have checked Cisinski's model-category and derivable-category treatments; their standard statements identify a single homotopy relation under their own hypotheses. The joint-congruence formulation is supplied here for the exact weak axioms on these pages, without claiming that it is new.
\end{dcomment}


Dans le cas biexacte (\textit{i.e.} celui des véritables structures de Quillen), les deux relations de Q-homotopie dans $\cM_{\mathrm{cf}}$ (plus généralement dans $\mathrm{Hom}(X,Y)$, si $X$ dans $\cM_{\mathrm{c}}$ et $Y$ dans $\cM_{\mathrm{f}}$)
\msp{page 32}
coïncident, et deviennent même particulièrement simples. (Car si $f, g : X \rightrightarrows Y$, avec $X$ cofibrant et $Y$ fibrant, alors $f\sim g$ peut se tester sur \textit{n'importe quel} “objet cylindre” pour $X$, ou sur \textit{n'importe quel} “objet chemins” pour $Y$.)

\dgrblock{dgrb688}{Why the localization has a fully faithful adjoint}
\begin{dcomment}{exposition}{p31}{Why the localization has a fully faithful adjoint}
The proposed adjoint is a way to recognize fibrant replacement from a universal property in the homotopy quotient. Assume left exactness, and write $\pi_{\mathrm d}\mathcal M_{\mathrm c}$ for the category of cofibrant objects modulo the equivalence relation generated by right Q-homotopies. Right homotopy is compatible with both compositions: precomposition is immediate, while postcomposition follows by extending a map between the constant paths across $s:Y\to PY$ using $TC\perp F$. Thus taking its generated equivalence relation really defines a category.

Choose $r_X:X\to RX$ in $TC$ with $RX$ fibrant. For any fibrant-cofibrant $Y$, restriction along $r_X$ gives a bijection
\[
 \operatorname{Hom}_{\pi_{\mathrm d}\mathcal M_{\mathrm c}}(RX,Y)
 \ \xrightarrow{\ \sim\ }\
 \operatorname{Hom}_{\pi_{\mathrm d}\mathcal M_{\mathrm c}}(X,Y).
\]
Surjectivity follows by lifting $r_X$ against $Y\to e$. For injectivity, a right homotopy between $ur_X$ and $vr_X$ extends across $r_X$ by lifting against the path fibration $PY\to Y^2$, giving a right homotopy between $u$ and $v$. For a chain of homotopies, choose extensions of all intermediate maps and apply this argument successively. The same proof works for chains with either orientation.

The previous block identifies the homotopy quotient on fibrant-cofibrant objects with $\mathrm{Ho}_W$. Moreover its inclusion in the larger quotient is fully faithful: any relation in the larger quotient gives equality after localization, and on fibrant-cofibrant objects such equality is already the generated right-homotopy relation. The displayed bijection therefore says exactly that the localization
\[
 L:\pi_{\mathrm d}\mathcal M_{\mathrm c}\longrightarrow\mathrm{Ho}_W
\]
is left adjoint to the inclusion of the fibrant-cofibrant homotopy category. The unit is represented by $r_X$. This proves the suggested fully faithful right adjoint without any claim that a chosen lift is unique before homotopy.

It also proves the stated calculus of fractions. Write $S$ for the images of the arrows of $W$ in this quotient; identities and compositions belong to $S$. Precomposition by $s:X\to X'$ in $S$ is bijective on maps into any fibrant-cofibrant object, since the displayed bijection identifies these Hom sets with those in the localization. Given $f:X\to Y$ and $s:X\to X'$ in $S$, choose the denominator $r_Y:Y\to RY$. That bijectivity supplies a map $g:X'\to RY$ with
\[
 gs=r_Yf,
\]
which is the common-denominator condition. For cancellation, suppose $f,g:X\to Y$ and $s:X'\to X$ in $S$ satisfy $fs=gs$. The same bijectivity, now with target $RY$, gives $r_Yf=r_Yg$. Thus the denominator $r_Y$ also supplies the cancellation condition. Every localized arrow consequently has a representative $X\to RY\leftarrow Y$, with the right-hand arrow in $S$; this specifies the manuscript's left-fraction convention.

Applying the whole argument to the opposite category proves the right-exact statement and the fully faithful left adjoint on the quotient of fibrant objects. The editors' changes from $\pi_{\mathrm d}$ to $\pi_{\mathrm g}$ on this page are therefore mathematically required.

For the cofibrant sources and fibrant targets considered here, the biexact case gives a single homotopy relation, computable using any cylinder or path object \citep[Lemma~2.2.12]{Cisinski2019HigherCategories}. Footnote~22 omits the cofibrancy assumption on the source and leaves the path convention implicit. The block \emph{Two path conventions and a nontransitive homotopy} resolves that distinction. For designated $TC$--$F$ paths, raw right homotopy can fail transitivity even when every object is fibrant and cofibrant if the weak structure is not exact; it can also fail in a perfect structure for a noncofibrant source. Both failures come from the graph example of \citet[Example~3.6]{KurzRosicky2005}. In contrast, the existential relation obtained by allowing all $W$--$F$ path objects is an equivalence relation for a fibrant target and an arbitrary source, by direct concatenation. The generated relation used in the adjoint argument above is therefore a deliberate and sufficient formulation.
\end{dcomment}


Dans le cas où la structure de Quillen faible est exacte à gauche, \textit{i.e.} où la relation d'équivalence pertinente dans $\mathrm{Fl}\,\cM_{\mathrm{cf}}$ est la Q-homotopie à droite, alors l'axiome d'exactitude est vérifié pour les carrés $W$-cartésiens, mais je doute qu'il en soit de même pour les carrés $W$-cocartésiens (\footnote{Il semble pourtant que c'est le cas, d'après \textsc{Baues}, cf. annotation marginale correspondante page 30.}). (Si on ne fait une hypothèse supplémentaire, p. ex. l'exactitude à droite de la structure de Quillen.) Si on n'a que l'exactitude (mais ni à gauche ni à droite), il semble qu'on ne peut en conclure l'exactitude ni des carrés $W$-cartésiens, ni des carrés $W$-cocartésiens, donc pas celle des suites de cosuspension (loop sequences), ni des suites de suspension. On pourrait donc penser que la condition d'exactitude, à elle seule, ne sert pas à grand
\msp{page 33}
chose pour le formalisme homotopique dans les catégories de fractions $\mathrm{Ho}_W$ etc. Mais on a pourtant ceci : soit $\Psi = {}$\begin{tikzcd}[row sep=tiny, column sep=small, baseline=(\tikzcdmatrixname-2-1.base)]
 & 1 \\
0 \arrow[ru] \arrow[rd] & \\
 & 2
\end{tikzcd}, et supposons que $(\cM,W)$ \textit{et} $(\cM(\Psi),W(\Psi))$ \texttt{[où pour toute petite catégorie }$I$\texttt{, }$\cM(I)$\texttt{ désigne la catégorie des préfaisceaux sur }$I$\texttt{ à valeurs dans }$\cM$\texttt{, et }$W(I)$\texttt{ est la partie de }\texttt{Fl}$(\cM(I))$\texttt{ formée des morphismes}
\printedpage{20}
% head: A. Grothendieck : Dérivateurs, Ch. XIII, Édition des Universités de Montpellier II et Paris VII
% foot: Version 25 janvier 2004
\texttt{de préfaisceaux qui sont dans }$W$\texttt{ "argument par argument"]} admettent une structure de Quillen faible exacte. Alors dans $\cM$ les carrés $W$-cartésiens sont $W$-exactes, donc on a dans $\cM$ des suites exactes longues de cosuspension (exact loop sequences) (\footnote{il semble que je me suis canulé ici, et je doute que le résultat évoqué soit exact.}). (L'hypothèse sur $(\cM(\Psi),W(\Psi))$ sert uniquement à assurer que toute flèche de $\mathrm{Ho}_{W(\Psi)} = \cM(\Psi)W(\Psi)^{-1}$ est isomorphe à une flèche provenant d'une flèche de $\cM(\Psi)$ lui-même. Mais c'est là une propriété apparemment très délicate, et je ne connais d'autre façon de l'obtenir que via l'existence d'une structure de Quillen faible exacte.) Dualement, si $(\cM,W)$ et $(\cM(\Phi),W(\Phi))$ admettent des structures de Quillen faibles exactes, alors les carrés $W$-cocartésiens dans $\cM$ sont $W$-exacts, et on trouve (si $\cM$ est ponctué) des suites exactes longues de suspension. % [sic] W-exactes -> W-exacts, agreement with "carrés" (cf. "W-exacts" below)

\dgrblock{dgrb689}{Exactness permits completion of the trivial classes}
\begin{dcomment}{development}{p32}{Exactness permits completion of the trivial classes}
Grothendieck's hesitation about exactness can be resolved by changing the designated trivial classes while keeping the weak equivalences, cofibrations and fibrations. This is a different operation from asserting that the original designated classes were already complete. It is therefore compatible with the counterexample to footnote~20.

\noindent\textbf{Completion under exactness.}\ours\quad
Suppose $(C,TC,F,TF)$ is an exact weak Quillen structure. Then
\[
 \bigl(C,\ C\cap W,\ F,\ F\cap W\bigr)
\]
is a biexact Quillen structure, with the same $C,F,W$. No retract-closure hypothesis on these classes is needed for this assertion. If all four original lifting classes were closed, the completion makes no change.

\emph{Proof.} The only new lifting problems involve $C\cap W$ against $F$, and $C$ against $F\cap W$. Let $j\in C\cap W$. Exactness gives $j=qi$ with $i\in TC$ and $q\in TF$. Since $j\in C$ and $q\in TF$, the lifting axiom exhibits $j$ as a retract of $i$ in the arrow category. Lifting properties are preserved by retracts, so $j$ has the left lifting property against every map in $F$. Dually, if $p\in F\cap W$, factor $p=qi$ with $i\in TC$ and $q\in TF$. Lifting $i$ against $p$ exhibits $p$ as a retract of $q$, so it has the right lifting property against $C$.

The original factorizations still serve, because $TC\subset C\cap W$ and $TF\subset F\cap W$. The classes $C$ and $F$ have not changed. The classes $C\cap W$ and $F\cap W$ have the required pushout or pullback stability by $\text{M}''$ and the universal-equivalence criterion of manuscript pp.~10--11; composition and identities follow from two-out-of-three. Thus every axiom holds, and both exactness equalities hold by definition. When $TC$ and $TF$ are retract closed, the two retract arguments also show $C\cap W\subset TC$ and $F\cap W\subset TF$, giving equality with the original classes.\hfill$\square$

This completion is a possibly nonclosed model structure in Quillen's original terminology. It suffices for the ordinary homotopy theory of cofibrant and fibrant objects: the standard endpoint arguments now use $C\cap W\perp F$ and $C\perp F\cap W$. In particular both homotopies classify the localized maps from a cofibrant object to a fibrant object, and any good chosen cylinder or path object from the original weak structure remains available. The existence of different designated classes therefore does not obstruct the usual homotopy calculation once exactness holds. This also explains the applicability of the nonclosed Quillen case in \citet[Example~5.2]{Cisinski2010}.

We checked the quoted Quillen and derivable-category frameworks and did not locate this completion stated using the exact four-class axioms of the manuscript. The proof is a direct application of the standard retract argument; its purpose is to settle the hypotheses on these pages, not to claim a new model-category theorem.
\end{dcomment}


\msp{page 34}
Ceci montre que si, pour une catégorie de modèles $(\cM,W)$ donnée, on s'intéresse à “faire de l'algèbre homotopique” pour tous les couples $(\cM(I),W(I))$ à la fois, avec $I$ dans Cat, il suffit (si on veut avoir une bonne théorie des fibres et cofibres $W$-homotopiques dans ces catégories, y inclus la propriété de $W$-exactitude pour les carrés $W$-cartésiens et pour les carrés $W$-cocartésiens) de disposer pour chaque $(\cM(I),W(I))$ d'une structure de Quillen faible exacte sans plus.

\dgrblock{dgrb690}{Exactness of squares after completing the classes}
\begin{dcomment}{exposition}{p33}{Exactness of squares after completing the classes}
Grothendieck introduces $\mathcal M(\Psi)$ to recover relations among maps while representatives remain available. Exactness in that diagram category indeed says that every localized arrow is isomorphic, in the arrow category, to an actual natural transformation: replace its source and target by fibrant-cofibrant objects and apply the quotient description above. His footnote correctly warns that this representability argument has not, by itself, established compatibility with the chosen homotopy pullback.

There is a more direct repair. The completion just proved turns an exact weak structure on $\mathcal M$ into a possibly nonclosed proper Quillen structure with the same $W,C,F$. Both homotopy-cartesian and homotopy-cocartesian exactness follow from that structure. For this conclusion the additional exact structure on $\mathcal M(\Psi)$ is unnecessary. Thus the asserted conclusion on p.~33 is valid, while the proposed route through arrow representability leaves a step unproved.

Here is the local lifting calculation behind the cartesian assertion. After taking the usual fibrant representatives, write the homotopy pullback as the ordinary pullback $P=E\times_SY$, where $p:E\to S$ is a fibration and $E,Y,S$ are fibrant. Let $T$ be cofibrant. Suppose maps $a:T\to E$ and $b:T\to Y$ have homotopic composites to $S$. Represent the homotopy by $H:IT\to S$, with endpoints $pa$ and $qb$. The endpoint inclusion $j_0:T\to IT$ is in $C\cap W$: it is a cofibration because $T$ is cofibrant, and a weak equivalence because the cylinder projection is one. The completed lifting property against $p$ gives a lift $\widetilde H:IT\to E$. Its other endpoint $a'=\widetilde H j_1$ is homotopic to $a$ and satisfies $pa'=qb$. The pair $(a',b)$ therefore gives a map $T\to P$. This proves the required surjectivity onto pairs of homotopy classes with equal composites. The dual argument gives cocartesian exactness and the corresponding suspension sequences in the pointed case.

The geometry of the calculation is visible in the cospan $*\to S^1\leftarrow *$. Its homotopy pullback is $\Omega S^1$, not the one-point strict pullback. Lifting a based loop to $\mathbb R$ with initial point $0$ assigns its endpoint in $\mathbb Z$; two loops are homotopic relative endpoints precisely when those integers agree. Hence $\pi_0\Omega S^1=\mathbb Z$. The path between the two composites is part of the homotopy-pullback data. The lifting argument preserves that datum while making the square strictly commute. This is the compatibility that mere representability of an arrow does not explain.
\end{dcomment}


Ceci nous montre donc l'intérêt de développer des critères, pour une $(\cM,W)$, assurant qu'il en est bien ainsi. À ce sujet, je note que je ne sais comment vérifier l'exactitude d'une structure de Quillen \textit{que si} je dispose de l'exactitude à gauche, ou à droite. C'est dire que je dois dégager en fait des conditions sur $(\cM,W)$ qui assurent (disons) que les $(\cM(I),W(I))$ admettent des structures de Quillen faibles exactes à gauche.

\msp{page 35}
Les structures de Quillen faibles pour les $(\cM(I),W(I))$ seront déduites d'une structure de Quillen faible pour $(\cM,W)$, de la façon suivante : elles seront toutes \textit{“closes à droite”,} \textit{i.e.}
\[
F(I) \;=\; (TC(I))_*, \qquad TF(I) \;=\; (C(I))_*, \tag*{(1)}
\]
en particulier (pour $I=e$)
\[
F \;=\; (TC)_*, \qquad TF \;=\; C_*,
\]
ce qui implique déjà (puisque $C\subset \mathrm{Cof}_W$, $TC\subset W^{\mathrm{univ}}$)
\[
(\mathrm{Cof}_W)_* \;\subset\; W, \qquad (W^{\mathrm{univ}})_* \;\subset\; \mathrm{Fib}_W \tag*{(2)}
\]
(d'où $(\mathrm{Cof}_W)_*\subset W_{\mathrm{univ}}$). En vertu de (1), pour définir les structures de Quillen faibles pour les $(\cM(I),W(I))$, il suffit de définir $C(I)$, $TC(I)$, et même, comme on veut l'exactitude à gauche, \textit{i.e.}
\[
TC(I) \;=\; C(I)\cap W(I) \tag*{(3)}
\]
\printedpage{21}
% head: A. Grothendieck : Dérivateurs, Ch. XIII, Édition des Universités de Montpellier II et Paris VII
% foot: Version 25 janvier 2004
(en particulier
\[
TC \;=\; C\cap W \quad ), \tag*{(4)}
\]
il reste à définir $C$. On définira $C(I)$ “argument par argument”, \textit{i.e.} pour $f\in\mathrm{Fl}(\cM(I))$
\[
f\in C(I) \qquad\Longleftrightarrow\qquad \alpha_s^*(f)\in C \ \text{ pour tout } s\in\mathrm{Ob}\,I, \tag*{(5)}
\]
où
\[
\alpha_s : e\longrightarrow I \qquad \texttt{[est]}\ \text{défini par } \alpha_s(e)=s.
\]
Il en résulte d'ailleurs aussitôt, par (3),
\[
f\in TC(I) \qquad\Longleftrightarrow\qquad \alpha_s^*(f)\in TC \ \text{ pour tout } s\in\mathrm{Ob}\,I. \tag*{(6)}
\]

\msp{page 36}
On suppose donc que $(\cM,W)$ admet une structure de Quillen faible exacte à gauche et close à droite $(C,TC,F,TF)$, définie en termes de
\[
C\;\subset\;\mathrm{Fl}(\cM),
\]
par
\[
TC \;=\; C\cap W, \qquad F \;=\; (TC)_*, \qquad TF = C_*.
\]
Pour que cela donne bel et bien une structure de Quillen faible (forcément exacte à gauche et close à droite), il faut simplement vérifier les conditions suivantes.

\dgrblock{dgrb691}{The injective choice of diagram structure}
\begin{dcomment}{exposition}{p35}{The injective choice of diagram structure}
Grothendieck now makes a choice which determines the variance of the whole construction. Cofibrations and weak equivalences will be tested at each object of the indexing category. Fibrations will be characterized by lifting against these objectwise classes. In modern terminology this is the injective choice. The asymmetry is productive: every restriction of diagrams preserves the classes specified objectwise, so its right adjoint becomes a candidate for a right Quillen functor.

Keep the variance explicit: $\mathcal M(I)=\operatorname{Fun}(I^{\mathrm{op}},\mathcal M)$. If $C(I)$ is the objectwise class associated with $C$, then
\[
 TC(I)=C(I)\cap W(I)=(C\cap W)(I)
\]
is an identity of classes, not an existence theorem. Once $F(I)=TC(I)_*$ and $TF(I)=C(I)_*$ are defined, lifting is automatic and the right classes have their required composition and pullback stability. Objectwise pushouts show the stability of $C(I)$; the corresponding statement for $TC(I)$ uses the universal weak-equivalence property already obtained in the base category.

The remaining two tasks have different characters. The inclusions $TF(I)\subset W(I)$ and $F(I)\subset\Fib_{W(I)}$ can be checked by evaluation and its adjoints; pp.~37--39 give that argument. Factorization in the diagram category requires a construction of coherent fillers and cannot be inferred merely by factoring every component. Even functorial componentwise factorizations do not make both lifting classes objectwise. The walking-arrow example below shows the distinction in the smallest nontrivial indexing category.
\end{dcomment}


\begin{itemize}[leftmargin=3em]
\item[\circled{i}] Factorisation pour $(C, TF=C_*)$ et $(TC=C\cap W, F=(TC)_*)$ – on reviendra sur ce point.
\item[\circled{ii}] Stabilité de $C$ par composition et cochangement de base, et $C$ contient les isomorphismes (d'où résultera la stabilité similaire de $TC = C\cap W$, compte tenu que $C\subset\mathrm{Cof}_W$) (NB Les stabilités idoines pour $F$, $TF$ sont triviales par définition $F=(TC)_*$, $TF=C_*$.)
\item[\circled{iii}] $C\subset\mathrm{Cof}_W$, $\underbrace{C_*}_{=\,TF}\subset W$, $(C\cap W)_*\subset\mathrm{Fib}_W$.
\end{itemize}

Pour vérifier, de même, que $(C(I),TC(I),F(I),TF(I))$ soit une structure de Quillen faible pour $W(I)$, il faut vérifier les conditions similaires à \circled{i}\,\circled{ii}\,\circled{iii} . Or (\textit{ii}) pour $C(I)$ résulte aussitôt de la condition (\textit{ii}) pour $C$. Il reste à examiner (\textit{iii}) et (\textit{i}).

\printedpage{22}
% head: A. Grothendieck : Dérivateurs, Ch. XIII, Édition des Universités de Montpellier II et Paris VII
% foot: Version 25 janvier 2004
\msp{page 37}
Les inclusions de (\textit{iii}).

L'inclusion
\[
C\;\subset\;\mathrm{Cof}_W
\]
implique que toute $f\in C(I)$ est dans $\mathrm{Cof}_W$ argument par argument. Cela implique que $f$ est dans $\mathrm{Cof}_{W(I)}$, comme on voit aussitôt sur la définition – je l'explicite en le

\textbf{Lemme 1 :} \textit{Soit $f : X_\bullet\longrightarrow Y_\bullet$ une flèche de $\cM(I)$. Pour que $f$ soit dans} $\mathrm{Cof}_{W(I)}$ (resp. \textit{dans} $\mathrm{Fib}_{W(I)}$), \textit{il} suffit \textit{que pour tout $s$ dans} $\mathrm{Ob}\,I$, \textit{$\alpha_s^*(f)$ soit dans} $\mathrm{Cof}_W$ (resp. \textit{dans} $\mathrm{Fib}_W$)\textit{. Itou (pour mémoire) pour une condition pour que $f\in W(I)$ – là il faut et il suffit que $\alpha_s^*(f)\in W\ \forall\, s\in\mathrm{Ob}\,I$.}

Nous allons appliquer ce même lemme pour établir les inclusions
\[
C(I)_* \;\subset\; W(I), \qquad (C(I)\cap W(I))_* \;\subset\; \mathrm{Fib}_{W(I)}.
\]
Donc il suffit de voir que
\[
\begin{aligned}
\alpha_s^*(C(I)_*) \;&\subset\; C_* &&\qquad (\subset\; W)\\
\alpha_s^*((\underbrace{C(I)\cap W(I)}_{(C\cap W)(I)})_*) \;&\subset\; (C\cap W)_* &&\qquad (\subset\; \mathrm{Fib}_W),
\end{aligned}
\]
et plus généralement
\[
\alpha_s^*(\Phi(I)_*) \;\subset\; \Phi_* \tag*{(7)}
\]
pour toute ( ?) partie $\Phi$ de $\mathrm{Fl}(\cM)$, définissant donc $\Phi(I)\subset\mathrm{Fl}(\cM(I))$ terme à terme.

\msp{page 38}
En explicitant la chose, on trouve la condition
\[
\alpha_{s\,!}(\Phi)\subset\Phi(I) \qquad \text{pour tout } s\in\mathrm{Ob}\,I,
\]
pour que ça marche, \textit{i.e.} on veut que
\[
i\in\Phi \qquad\Longrightarrow\qquad \alpha_{s\,!}(i)\in\Phi(I), \quad \textit{i.e.}\ \ \alpha_t^*(\alpha_{s\,!}(i))\in\Phi\ \ \forall\, t\in\mathrm{Ob}\,I.
\]
Or le foncteur $\alpha_t^*\alpha_{s\,!}$ s'explicite comme
\[
A \longmapsto \underbrace{\mathrm{Hom}_I(t,s)\times A}_{\substack{\text{somme de } \mathrm{Hom}_I(t,s)\\ \text{copies de } A \text{ dans } \cM}},
\]
donc il faut supposer $\Phi$ \textit{stable par sommes directes} (du type envisagé $\mathrm{Hom}_I(t,s)\times A$). C'est bien le cas pour $\Phi = C$ et $\Phi = C\cap W$ et pour des sommes directes \textit{finies,} grâce à (\textit{ii}). On trouve donc le

\printedpage{23}
% head: A. Grothendieck : Dérivateurs, Ch. XIII, Édition des Universités de Montpellier II et Paris VII
% foot: Version 25 janvier 2004
\textbf{Lemme 2.}

\begin{itemize}[leftmargin=3em]
\item[a)] \textit{Soit $I$ dans} Cat \textit{telle que les sommes directes de type} $\mathrm{Hom}_I(t,s)$ \textit{existent dans $\cM$. (P. ex. on suppose que $I$ est} finie, \textit{plus généralement que les} $\mathrm{Hom}_I(t,s)$ \textit{sont des ensembles finis ; ou que les “petites” sommes directes existent dans $\cM$.) Soit $\Phi$ une partie de} $\mathrm{Fl}(\cM)$\textit{, stable par ce type de sommes directes, et soit} $\Phi(I)\subset\mathrm{Fl}(\cM(I))$ \textit{la partie correspondante. Alors pour tout $s\in\mathrm{Ob}\,I$, on a}

\msp{page 39}
\[
\alpha_s^*(\Phi(I)_*) \;\subset\; \Phi_*.
\]
\item[b)] \textit{Supposons que les} $\mathrm{Hom}_I(t,s)$ \textit{soient finis (p. ex. $I$ finie). Alors l'hypothèse de} a) \textit{s'applique pour $\Phi=C$ et pour $\Phi=C{\cap}W$, donc la condition} (iii) \textit{p.} 36 \textit{est satisfaite pour la structure $(W(I),C(I),TC(I),F(I),TF(I))$ dans $\cM(I)$. Même conclusion sans hypothèse de finitude sur les} $\mathrm{Hom}_I(t,s)$\textit{, quand on suppose $\cM$ stable par “petites” sommes directes, et $C$ et $W$ stables également par de telles sommes directes.}
\end{itemize}

\dgrblock{dgrb692}{Evaluation transports the lifting problem}
\begin{dcomment}{exposition}{p38}{Evaluation transports the lifting problem}
The copower formula is the essential calculation in the two lemmas. Let $s$ be an object of $I$. The left adjoint to evaluation at $s$ is
\[
 (\alpha_{s!}A)_t=\coprod_{\operatorname{Hom}_I(t,s)}A.
\]
For an arrow $t'\to t$, precomposition sends $\operatorname{Hom}_I(t,s)$ to $\operatorname{Hom}_I(t',s)$ and defines the presheaf restriction map. A natural transformation $\alpha_{s!}A\to X$ is uniquely determined by the map from the summand indexed by $1_s$ to $X_s$; all other summands are obtained by restriction. This proves the adjunction and fixes the direction of the Hom set.

Suppose $\Phi$ is stable under the copowers appearing in this formula. Then $\alpha_{s!}(i)\in\Phi(I)$ for every $i\in\Phi$. For $p\in\Phi(I)_*$, a lifting square from $i$ to $p_s$ transposes to a square from $\alpha_{s!}i$ to $p$. Its diagonal transposes back to the required lift. Hence
\[
 \alpha_s^*(\Phi(I)_*)\subset\Phi_*.
\]
The statement on p.~37 for an arbitrary class, accompanied by a question mark, therefore needs exactly the additional closure used in Lemma~2. We are not assuming that evaluation itself preserves right orthogonals without an adjoint test.

Finite copowers require no new hypothesis for $C$ or $TC$. For maps $i:A\to B$ and $j:A'\to B'$, the coproduct map factors as
\[
 A\amalg A'\longrightarrow B\amalg A'\longrightarrow B\amalg B',
\]
where each factor is a pushout of one of the original maps. The empty copower is an identity on the initial object. Induction proves the finite statement. For arbitrary copowers one must either assume their stability or obtain it from transfinite composition by adjoining the summands in a well-order. The required colimits must exist.

Finally, finite limits and colimits of presheaves are computed objectwise. A double pushout diagram testing a $W(I)$-cofibration evaluates to the corresponding test in $\mathcal M$, and an objectwise weak equivalence remains one after all these tests. This proves Lemma~1 and its dual. Combining it with the lifting argument proves all the inclusions on p.~36 under the finite-Hom hypothesis, while deliberately leaving factorization for the next stage.
\end{dcomment}


\textbf{NB} Vu le scholie plus haut, il est parfaitement raisonnable de se borner au besoin aux catégories $I$ \textit{finies} – vu que la catégorie de ces $I$ est stable par les produits $I\longmapsto I\times\Psi$, $I\longmapsto I\times\Phi$.

Il reste, pour être heureux, à vérifier la condition de factorisation \circled{i} pour $(C,TF)$, $(TC,F)$, mais aussi et surtout pour les $(C(I),TF(I),TC(I),F(I))$. C'est dire que de partir même d'une structure de Quillen close (pays de cocagne)
\msp{page 40}
ne nous avance guère à lui tout seul, faute de savoir que l'axiome de factorisation se conserve, en passant de $(\cM,W,\ldots)$ à $(\cM(I),W(I),\ldots)$.

\dgrblock{dgrb693}{A complete calculation for the walking arrow}
\begin{dcomment}{examples}{p39}{A complete calculation for the walking arrow}
The formulas become especially transparent when $I=[1]$, with its arrow $0\to1$. A presheaf is a map $X_1\to X_0$. Let $\Phi$ contain identities and be stable under finite coproducts, and put $R=\Phi_*$. Write $x:X_1\to X_0$ and $y:Y_1\to Y_0$ for the structural maps. For a natural transformation $p:X\to Y$, consider the matching map
\[
 m_p:X_1\longrightarrow Y_1\times_{Y_0}X_0.
\]
Then $p$ has the right lifting property against every objectwise $\Phi$-map if and only if
\[
 p_0:X_0\to Y_0\ \in R,
 \qquad m_p\in R.
\]

To prove sufficiency, take a square from an objectwise $\Phi$-map $A\to B$ to $p$. First lift its component at $0$ against $p_0$, obtaining $B_0\to X_0$. The component at $1$, together with this already chosen map composed with $B_1\to B_0$, gives a lifting square against $m_p$. A lift provides $B_1\to X_1$ and enforces precisely the naturality relation. Thus the two component lifts fit together.

For necessity, evaluation at $0$ and its left adjoint show $p_0\in R$. To test $m_p$ against $i:A\to B$ in $\Phi$, write a square as maps $a:A\to X_1$, $b_1:B\to Y_1$, and $b_0:B\to X_0$ with
\[
 p_1a=b_1i,\qquad xa=b_0i,\qquad p_0b_0=yb_1.
\]
Use the map of arrow diagrams
\[
 (A\xrightarrow{i}B)\longrightarrow(B\xrightarrow{1}B),
\]
whose components are $i$ at $1$ and $1_B$ at $0$. This is objectwise in $\Phi$. A lift of its induced square against $p$ is exactly a solution to the original square against $m_p$. This proves necessity.

Apply the calculation in sets with $\Phi$ the injections and $R$ the surjections, assuming choice. Let $E_1=\{a\}$, $E_0=\{0,1\}$, and let $E_1\to E_0$ send $a$ to $0$. The transformation from $E$ to the terminal arrow diagram is surjective at both objects. Its matching map is the non-surjective map $\{a\}\to\{0,1\}$, so it is not in $\Phi(I)_*$. Explicitly, a proposed lift choosing the element $1$ at vertex $0$ cannot be extended to vertex $1$. This is a failure of naturality, despite the existence of lifts at each object separately.

For the injective model structure, taking $\Phi=C$ gives the characterization of trivial fibrations, and taking $\Phi=TC$ gives the characterization of fibrations. It follows in particular that an injectively fibrant arrow has fibrant target and a fibration as its structural arrow. The general admissible-edge result later in the manuscript is an extension of this relative lifting calculation, rather than a consequence of pointwise fibrancy alone.
\end{dcomment}


C'est ici que le théorème de factorisation du no.~1 doit s'avérer providentiel. Mais pour pouvoir l'appliquer aux couples $(C(I),C(I)_*)$, $(TC(I),TC(I)_*)$, il faut (en plus de l'accessibilité de $\cM$, ce qui est une condition fort anodine) des conditions de stabilité très fortes pour $C(I)$, $TC(I)$ (\footnote{Plus les \textit{petits} systèmes de générateurs pour $C$, $TC$.}). À vrai dire, il reste la stabilité de $C$ et de $C\cap W$ par limites inductives ordinales (ce qui impliquera la même stabilité pour les $C(I)$, $(C{\cap}W)(I) = C(I){\cap}W(I)$).

À dire vrai, je vois peu de chances de trouver une structure de Quillen faible $(C,TC,F,TF)$ associé à $W$, et satisfaisant toutes les conditions dites, y compris celle de stabilité de $C$ et de $C\cap W$ par $\varinjlim$ ordinales, sans supposer que $W$ elle-même ne soit déjà stable par (petites) limites inductives filtrantes. % [sic] associé -> associée, agreement with "structure"

\msp{page 41}
On peut songer à utiliser une factorisation \textit{fonctorielle} des flèches dans $\cM$ en
\[
\begin{aligned}
f = pi \quad \text{avec}\quad & i\in TC, && p\in F\\
\text{ou}\quad & i\in C, && p\in TF,
\end{aligned}
\]

\dgrblock{dgrb694}{The factorization problem remains a coherence problem}
\begin{dcomment}{later}{p40}{The factorization problem remains a coherence problem}
The final paragraph should be read as a choice of a sufficient strategy, not as a necessity theorem about weak equivalences. Transfinite closure of cofibrations and trivial cofibrations is part of the small-object method; closure of every weak equivalence under every filtered colimit is a convenient, stronger condition from which Grothendieck hopes to obtain it. Nothing on this page proves that such closure of $W$ is necessary for diagram model structures.

The later injective existence theorem in the combinatorial setting answers the construction problem with explicit generation hypotheses; its precise application belongs to the discussion of manuscript pp.~42--44. There is also a different route through weaker diagram-compatible structures: \citet[Theorem~6.22]{Cisinski2010} proves stability under presheaf categories for homotopically complete model categories in his sense; the theorem also covers homotopically complete and cocomplete weak model categories. These are two ways to obtain the coherence which the walking-arrow calculation detects. Neither licenses treating the pointwise right class as the required right lifting class.
\end{dcomment}


% ---- from Der-ChXIII-25-01-04_p026-030.tex
\printedpage{24}
% head: A. Grothendieck : Dérivateurs, Ch. XIII, Édition des Universités de Montpellier II et Paris VII
pour déduire une factorisation similaire pour les flèches $f$ de $\M(I)$, en
\[
\begin{aligned}
f = pi \quad \text{avec}\quad & i \in TC(I), && p \in F(I)\\
\text{ou}\quad & i \in C(I), && p \in TF(I),
\end{aligned}
\]
où $C(I)$, $TC(I)$, $F(I)$, $TF(I)$ sont ici définis via $C$, $TC$, $F$, $TF$ “argument par argument”. Mais l'ennui, c'est que en procédant de cette façon naïve, on n'aura plus nécessairement $C(I) \longleftrightarrow TF(I)$, $TC(I) \longleftrightarrow F(I)$. Il faut considérablement rapetisser $TF(I)$ et $F(I)$, pour avoir une propriété RLP par rapport à $C(I)$ resp. $TC(I)$.

Donc je ne vois guère d'autre moyen de s'en tirer, dans ce contexte général, qu'en supposant la stabilité de $W$ par $\limind$ filtrantes (ce qui n'est après tout pas très méchant).

\noindent\msp{page 42}

\dgrblock{dgrb695}{Natural factorizations and natural fillers}
\begin{dcomment}{exposition}{p41}{Natural factorizations and natural fillers}
Grothendieck has reached the obstruction which makes a theory of diagrams more than a repetition of the theory of objects. A functorial factorization in $\mathcal M$ does produce a factorization of every natural transformation. Its middle objects and their transition maps fit together because the original factorization is functorial on the arrow category. But a choice of a filler for each component of a lifting problem need not commute with the transition maps. The missing datum is a compatible family of fillers, and the factorization functor supplies no such family.

The walking-arrow calculation in the commentary at manuscript pp.~38--39 makes this failure completely explicit. For diagrams $X_1\to X_0$ and $Y_1\to Y_0$, right lifting against all objectwise maps in a left class $L$ requires both
\[
 X_0\longrightarrow Y_0\in R,
 \qquad X_1\longrightarrow Y_1\mathbin{\times}_{Y_0}X_0\in R,
 \qquad R=L_*.
\]
The second condition remembers the compatibility of the lift with the arrow. Requiring only $X_1\to Y_1\in R$ forgets it. This is the first matching-object condition of the later Reedy formalism; the general condition appears in \citet[Definition~5.2.2 and Theorem~5.2.5]{Hovey1999}.

There are consequently two different ways to organize the model structure of a diagram category. In the injective structure the cofibrations and weak equivalences are objectwise, and fibrations are defined by a global right lifting property. In the projective structure the fibrations and weak equivalences are objectwise, and the cofibrations carry the compatibility condition. The manuscript is pursuing the injective structure. Its decision to keep $C(I)$ and $TC(I)$ objectwise therefore forces the replacement of the naive classes of fibrations by the smaller orthogonal classes $(TC(I))_*$ and $C(I)_*$.

This choice has a purpose beyond obtaining factorizations. Restriction along any functor preserves the two objectwise left classes. Its right adjoint will therefore preserve the two right classes, which is exactly the input needed to construct homotopy right Kan extensions. Thus the apparent asymmetry of the construction already prepares the three variances of the next section.
\end{dcomment}


Bon gré, mal gré, on revient à la situation envisagée dans le paragraphe 2. On va réénoncer l'essentiel de ce qui a été pressenti alors :

\textbf{Théorème 2} (encore un peu conjectural). {\itshape Soit $(\M,W)$ une catégorie de modèles, avec $W$ (comme toujours) faiblement saturée, mais de plus on suppose $\M$ stable par petites $\limind$ et par $\limproj$ finies, et $W$ stable par $\limind$ filtrantes, enfin $\M$ accessible. Considérons les conditions suivantes :

\begin{enumerate}
\item[\textup{a)}] Il existe une structure de Quillen faible pour $W$, exacte à gauche, close à droite, $(C,TC,F,TF)$.
\item[\textup{a$'$)}] Comme \textup{a)}, mais avec de plus $C$, $TC$ clos à gauche (\textup{i.e.} on a une structure de Quillen faible close), et $C$, $TC$ de la forme $\widetilde{C_0}$, $\widetilde{TC_0}$, où $C_0$ et $TC_0$ sont petits.
\item[\textup{a$''$)}] Il existe $C_0 \subset \Cof_W$, $TC_0 \subset W^{\univ}$, parties petites, telles que $C_{0*} \subset W$, $(TC_0)_* \subset \Fib_W$.
\item[\textup{b)}] On a
\[
(\Cof_W)_* \;\subset\; W_{\univ}, \qquad\quad (W^{\univ})_* \;\subset\; \Fib_W
\]
(d'où résulte, $\Cof_W$ et $W^{\univ}$ étant clos à gauche, que
\[
(W_{\univ})^* \subset \Cof_W, \qquad\quad (\Fib_W)^* \subset W^{\univ} \qquad\quad \text{\itshape ).}
\]
\end{enumerate}

\noindent\msp{page 43}

Ces conditions impliquent que pour toute $I$ dans \textup{Cat}, il existe une structure de Quillen faible exacte à gauche et close pour $W(I)$ dans $\M(I)$. Une telle structure se définit, en termes d'une structure similaire $C$, $TC = C \cap W$, $F = (TC)_*$, $TF = C_*$ pour $W$, par}
\[
C(I) \;=\; \Bigl\{ f \in \Fl(\M(I)) \Bigm| \alpha_s^*(f) \in C \ \ \forall\, s \in \Ob I \Bigr\},
\]

\printedpage{25}
% head: A. Grothendieck : Dérivateurs, Ch. XIII, Édition des Universités de Montpellier II et Paris VII
\noindent\textit{puis}
\[
\begin{aligned}
TC(I) \;&=\; C(I)\cap W(I) \;=\; \Bigl\{ f \in \Fl(\M(I)) \Bigm| \alpha_s^*(f) \in TC \ \ \forall\, s \in \Ob I \Bigr\}\\
F(I) \;&=\; (TC(I))_*\\
TF(I) \;&=\; C(I)_*\,.
\end{aligned}
\]

\textbf{Corollaire-Scholie.} \textit{Ces conditions impliquent que pour tout $I$ dans} Cat, \textit{$(\M(I),W(I))$ donne lieu à une bonne théorie des fibres et cofibres $W(I)$-homotopiques, avec carrés $W(I)$-cartésiens $W(I)$-exacts, et carrés $W(I)$-cocartésiens $W(I)$-exacts – donc avec des suites} exactes \textit{longues de cosuspension et de suspension, dans le cas où $\M$, donc aussi $\M(I)$, est une catégorie avec objet neutre.}

\noindent\msp{page 44}

\dgrblock{dgrb696}{The injective existence theorem and its precise hypothesis}
\begin{dcomment}{later}{p43}{The injective existence theorem and its precise hypothesis}
Grothendieck wants the factorization theorem to work uniformly after adding any small indexing category. One part of this proposal is now a standard existence theorem: a combinatorial model category has injective and projective model structures on each small diagram category. In the terminology of the present pages, a closed structure exact on the left is already an ordinary closed model structure, and condition $a'$ supplies the two sets of generators. Since an accessible cocomplete category is locally presentable, $a'$ places us in the combinatorial situation.

For the injective construction, the set of generators on the diagram category is the substantive point. One cannot simply use the free diagrams on the generators at each object: those give the projective left class. A theorem of \citet[Corollary~3.4, author preprint of 21 November 2013]{MakkaiRosicky2014Cellular} instead says that the pointwise left class associated with a set-generated, retract-closed cellular structure on a locally presentable category is again set-generated. Apply it separately to $C$ and $TC$. The two resulting small-object factorizations have left classes $C(I)$ and $TC(I)$ and right classes $C(I)_*$ and $(TC(I))_*$.

Here is why the compatibility of these two factorizations survives. Evaluation takes either right class into its corresponding right class in $\mathcal M$: its left adjoint is the copower by the representable presheaf, and copowers preserve $C$ and $TC$. Consequently $C(I)_*\subset W(I)$, while an element of $(TC(I))_*$ which belongs to $W(I)$ is objectwise in $TF$. To obtain the reverse lifting assertion without confusing componentwise and natural fillers, factor such a map $f$ using $(C(I),C(I)_*)$, say $f=pi$. The maps $p$ and $f$ are weak equivalences, so $i\in C(I)\cap W(I)=TC(I)$. The lifting property of $f$ against $i$ exhibits $f$ as a retract of $p$. Hence $f\in C(I)_*$. This proves the second trivial-class equality, and the model axioms follow.

Properness also passes to these structures. Pushouts and pullbacks are objectwise; injective cofibrations are objectwise cofibrations, and injective fibrations are objectwise fibrations by the preceding evaluation argument. The base properness axioms therefore apply at every object. These conclusions are recorded in \citet[Theorems~1.17 and 1.19, Propositions~1.21--1.22]{Barwick2007}; the numbering here is that of the 2007 preprint.

This settles the diagram construction under $a'$. It does not establish every proposed recognition implication involving $a$, $a''$ or $b$. Closure under filtered colimits alone is not the solution-set condition in a model-structure recognition theorem. The standard Smith theorem requires an accessible, accessibly embedded class of weak equivalences in the arrow category, as well as a specified set of generating cofibrations; see \citet[Proposition~1.7]{Barwick2007}. No assertion that the whole conjectural theorem has been proved is justified by citing the injective existence theorem. Conversely, once a combinatorial model structure is given, its diagram structures do not require preservation of \emph{every} filtered colimit by the weak equivalences. That stronger hypothesis belongs to Grothendieck's proposed route to recognition.
\end{dcomment}


Je m'occuperai de la démonstration du théorème précédent, ainsi que du théorème de factorisation du paragraphe 1, quand j'aurai mes notes sur les catégories et foncteurs accessibles et filtrations cardinales. En attendant, j'ai envie déjà d'expliciter un peu plus le formalisme auquel on parvient moyennant les hypothèses faites ci-dessus. (Je suis à peu près sûr tout au moins que la condition $a''$) suffit à impliquer toutes les autres. Et c'est bien sous la forme $a''$) que ces conditions vont se vérifier en pratique, je présume, notamment dans le cas où $\M = $ Cat.)

\dgrblock{dgrb697}{The variance is fixed by presheaves}
\begin{dcomment}{exposition}{p44:En attendant, j'ai envie déjà d'expliciter un peu plus}{The variance is fixed by presheaves}
Grothendieck now turns the diagram model structures into operations. Throughout these pages
\[
 \mathcal M(I)=\operatorname{Fun}(I^{\mathrm{op}},\mathcal M),
 \qquad u:I'\longrightarrow I.
\]
This convention determines the comma categories in both pointwise Kan formulas. If $X':(I')^{\mathrm{op}}\to\mathcal M$, then
\[
 \begin{aligned}
 (u_!X')(i)&=\operatorname*{colim}_{(i\downarrow u)^{\mathrm{op}}}X'(i'),\\
 (u_*X')(i)&=\operatorname*{lim}_{(u\downarrow i)^{\mathrm{op}}}X'(i').
 \end{aligned}
\]
The manuscript writes $i\backslash I'$ and $I'/i$, with the contravariant convention understood. Making the opposites explicit explains the otherwise surprising occurrence of an \emph{initial} object in footnote~27: it becomes terminal in the indexing category of the colimit. It also prevents interchanging the two comma categories when passing to modern covariant notation.

The lifting argument behind Proposition~2 is a precise form of the principle that adjunction transfers extension problems. Suppose $L:\mathcal A\rightleftarrows\mathcal B:R$ and let $i:A\to A'$ and $p:B\to B'$. By the adjunction, the two horizontal maps in a square from $Li$ to $p$ correspond to the two horizontal maps in a square from $i$ to $Rp$. Naturality of the adjunction bijection identifies the commutativity conditions. A diagonal $LA'\to B$ corresponds to a diagonal $A'\to RB$, and each of the two triangle identities required of a filler is carried to the corresponding identity. The sets of fillers are therefore naturally bijective. In particular
\[
 Li\mathrel{\perp}p\quad\Longleftrightarrow\quad i\mathrel{\perp}Rp.
\]
Applying this first to $u^*\dashv u_*$ proves the preservation of $F$ and $TF$ by $u_*$. Applying it to $u_!\dashv u^*$ proves the two equivalences in (3). This is the same elementary mechanism as \citet[Lemma~1.3.4]{Hovey1999}.

For a useful special case, suppose every component of $i\downarrow u$ has an initial object $e_\alpha$. The left Kan formula reduces to
\[
 (u_!X')(i)\cong\coprod_\alpha X'(e_\alpha).
\]
A left orthogonal class is closed under coproducts: a square out of a coproduct is a family of squares, and a family of fillers assembles to a filler. Thus $u_!$ preserves $C$ and $TC$ under this hypothesis, and $u^*$ preserves $F$ and $TF$. No claim about preservation of general colimits by the weak equivalences is involved. The comma-category hypothesis has replaced that issue by a coproduct calculation.
\end{dcomment}


\section*{4\quad Les variances $u^*$, $u_*$, $u_{!}$ pour $(\M(I),W(I))$ (I variable)}

Je suppose les conditions du paragraphe précédent vérifiées pour $(\M,W)$ (cf. théorème p. 42), et choisis une structure de Quillen faible close et exacte à gauche $(C,TC,F,TF)$ pour $W$, d'où pour tout $I$ dans Cat

\noindent\msp{page 45}

\noindent une structure de Quillen faible close et exacte à gauche $(C(I),TC(I),F(I),TF(I))$ pour $W(I)$, dans $\M(I)$, décrite page 43. Soit maintenant
\[
u \;:\; I' \longrightarrow I
\]
une flèche dans Cat, d'où une flèche
\[
u^* \;:\; (\M(I),W(I)) \;\longrightarrow\; (\M(I'),W(I')).
\]
On sait, vu que $\M$ est stable par $\limind$ quelconques et aussi (étant de plus accessible) par $\limproj$ quelconques, que $u^*$ admet des foncteurs adjoints à gauche et à droite,
\[
u_{!},\ u_* \;:\; \M(I') \;\longrightarrow\; \M(I),
\]
\printedpage{26}
% head: A. Grothendieck : Dérivateurs, Ch. XIII, Édition des Universités de Montpellier II et Paris VII
et je me propose ici de résumer les propriétés les plus immédiates de $u^*$ et de $u_{!}$, $u_*$ par rapport aux structures de Quillen faibles envisagées.

\textbf{Proposition 1 :} \textit{Avec les notations précédentes, on a les relations :}
\[
u^*(C(I)) \subset C(I'), \qquad u^*(TC(I)) \subset TC(I'). \tag*{(1)}
\]
\[
u_*(F(I')) \subset F(I), \qquad u_*(TF(I')) \subset TF(I). \tag*{(2)}
\]
\[
\begin{aligned}
u^*(F(I)) \subset F(I') &\iff u_{!}(TC(I')) \subset TC(I)\\
u^*(TF(I)) \subset TF(I') &\iff u_{!}(C(I')) \subset C(I).
\end{aligned} \tag*{(3)}
\]

D'autre part, je ne sais rien dire de positif pour les inclusions
\[
\begin{aligned}
u_*(C(I')) \;&\subset\; C(I), \;\; u_*(TC(I')) \;\subset\; TC(I)\\
u_{!}(F(I')) \;&\subset\; F(I), \;\; u_{!}(TF(I')) \;\subset\; TF(I).
\end{aligned}
\]
Si ce n'est que pour un $u$ tant soit peu général ($u$ pas un isomorphisme, disons), $u_*$ n'a aucune envie

\noindent\msp{page 46}

\noindent d'être compatible avec les structures $C$, $TC$ (propriété d'exactitude homotopique à droite, moralement, vu que $u_*$ est exact à gauche et en général pas du tout à droite), ni $u_{!}$ avec les structures $F$, $TF$ (commentaire dual à donner, $u_{!}$ est exacte à \textit{droite}, et n'a pas envie d'être “homotopiquement exact à gauche”). % [sic] exacte -> exact, accord avec « foncteur » / « exact à gauche »

\textsc{Démonstration} de (1), (2), (3). Les relations (1) sont évidentes par définition. Les relations (2), (3) le sont à peine moins, vu la propriété soritale générale suivante, dont la place serait au paragraphe 1 :

\textbf{Proposition 2.} \textit{Soit}
\[
\begin{tikzcd}
C' \arrow[r, shift left, "u_{!}"] & C \arrow[l, shift left, "u^*"]
\end{tikzcd}
\]
\textit{un couple de foncteurs adjoints $(u_{!},u^*)$, et soit $\Phi' \subset \Fl(C')$, $\Psi \subset \Fl(C)$. Alors}
\[
(u_{!}(\Phi') \longleftrightarrow \Psi) \qquad\quad \iff \qquad\quad (\Phi' \longleftrightarrow u^*(\Psi)).
\]

\textbf{Corollaire.} (C'est le même énoncé, avec seulement des notations changées.) \textit{Supposons que $u^* : C \longrightarrow C'$ admette un adjoint à} droite \textit{$u_*$ (pas forcément un adjoint à gauche)}
\[
\begin{tikzcd}
C' \arrow[r, shift right, "u_*"'] & C, \arrow[l, shift right, "u^*"']
\end{tikzcd}
\]

\printedpage{27}
% head: A. Grothendieck : Dérivateurs, Ch. XIII, Édition des Universités de Montpellier II et Paris VII
\noindent\textit{soit $\Phi \subset \Fl(C)$, $\Psi' \subset \Fl(C')$, alors}
\[
(u^*(\Phi) \longleftrightarrow \Psi') \qquad\quad \iff \qquad\quad (\Phi \longleftrightarrow u_*(\Psi')).
\]

C'est la proposition 2 pour \begin{tikzcd} C'_0 \arrow[r, shift left, "u_{0\,!}"] & C_0 \arrow[l, shift left, "u_0^*"] \end{tikzcd}, avec $C'_0 = C$, $C_0 = C'$, $(u_{0\,!} : C \longrightarrow C') = u^*$, et $(u_0^* : C' \longrightarrow C) = u_*$.

\noindent\msp{page 47}

Pour les inclusions (2), on utilise le corollaire de la proposition 2 pour les couples $\Phi = TC(I)$, $\Psi' = F(I')$, resp. $\Phi = C(I)$, $\Psi' = F(I')$ \texttt{[plutôt} $\Psi' = TF(I')$\texttt{]}, et on applique (1) qui donne $u^*(\Phi) \longleftrightarrow \Psi'$, d'où $\Phi \longleftrightarrow u_*(\Psi')$, ce qui n'est autre que (2). Pour les équivalences (3), on utilise la proposition 2, en faisant $\Psi = F(I)$, $\Phi' = TC(I')$ resp. $\Psi = TF(I)$, $\Phi' = C(I')$.

\textbf{Corollaire 1} (de la proposition 1) : \textit{Supposons $u : I' \longrightarrow I$ telle que pour tout $i \in \Ob I$, la catégorie $i\backslash I'$ soit discrète} (\footnote[27]{Plus généralement, il suffit que toute composante connexe $Z_\alpha$ de $i\backslash I'$ ait un objet initial $e_\alpha$, de sorte que $\limind\limits_{i\backslash I'} X'(i') \simeq \coprod_{\alpha\in\pi_0(i\backslash I')} X'(e_\alpha)$. C'est le cas notamment si $I' \longrightarrow I$ est un isomorphisme local, p. ex. une immersion ouverte.}) \textit{(p. ex. $I'$ discrète). Alors on a}
\[
\left\{
\begin{aligned}
u_{!}(C(I')) \;&\subset\; C(I), \;\; u_{!}(TC(I')) \;\subset\; TC(I)\\
u^*(F(I)) \;&\subset\; F(I'), \;\; u^*(TF(I)) \;\subset\; TF(I').
\end{aligned}
\right.
\]

En effet, les relations de la deuxième ligne équivalent (à cause de (3)) à celles de la première. Il faut prouver que si $f$ dans $\Fl\M(I')$ est dans $C(I')$ resp. $TC(I')$, $u_{!}(f)$ est dans $C(I)$ resp. $TC(I)$, ce qui signifie aussi que $\alpha_s^*(u_{!}(f))$ est dans $C$ resp. $TC$, pour tout $s \in \Ob I$. Or le foncteur
\[
\alpha_s^* u_{!} \;:\; \M(I') \;\longrightarrow\; \M
\]
s'explicite par la formule bien connue
\[
\alpha_s^*(u_{!}(X')) \;\overset{\text{déf}}{=}\; u_{!}(X')(s) \;\simeq\; \limind\limits_{i'\in s\backslash I'} X'(i').
\]

Comme

\noindent\msp{page 48}

\noindent $s\backslash I'$ est discrète, on a donc l'isomorphisme fonctoriel en $X'$ dans $\M(I')$
\[
\alpha_s^*(u_{!}(X')) \;\simeq \coprod_{(s',\, x \,:\, s \,\longrightarrow\, u(s')) \text{ dans } \Ob\,(s\backslash I')} X'(s'),
\]
donc si $f \,:\, X' \longrightarrow Y'$, $\alpha_s^*(u_{!}(f))$ s'explicite comme une somme directe dans $\underline{\mathrm{Fl}}(\M)$ de flèches de la forme
\[
f_{s'} \;:\; X'(s') \;\longrightarrow\; Y'(s'),
\]
\printedpage{28}
% head: A. Grothendieck : Dérivateurs, Ch. XIII, Édition des Universités de Montpellier II et Paris VII
lesquelles sont dans $C$ resp. $TC$ puisque $f$ est dans $C(I')$ resp. $TC(I')$. Comme $C$, $TC$ sont stables par sommes directes quelconques ($C$, $TC$ étant clos à gauche), on gagne.

\textbf{Corollaire 2.} \textit{Soit $I$ dans} Cat, \textit{et $f \in \Fl(\M(I))$. Pour que $f$ soit dans $F(I)$} (resp. \textit{dans $TF(I)$\,), il} faut \textit{que $\forall\, s \in \Ob I$, $f_s : X(s) \longrightarrow Y(s)$ le soit.} (Mais cette condition n'a pas du tout l'air d'être suffisante !) \textit{En particulier, soit $X \in \Ob\M(I)$, alors $X \in \Ob\M(I)_f \Longrightarrow X(s) \in \Ob\M_f \ \ \forall\, s \in \Ob I$.}

\dgrblock{dgrb698}{Ordinary limits and the two homotopical exactness properties}
\begin{dcomment}{examples}{p48:Mais cette condition n'a pas du tout l'air}{Ordinary limits and the two homotopical exactness properties}
Grothendieck's hesitation on pp.~49--50 identifies a necessary distinction. A right Kan extension can preserve fibrations and trivial fibrations without preserving weak equivalences between arbitrary diagrams. Its derived functor is computed on fibrant diagrams, and it is this functor which preserves homotopy limits. The distinction already matters for a finite poset.

Work in simplicial sets, and let $K$ be the nerve of the groupoid with two objects and exactly one morphism between every ordered pair. The simplicial set $K$ is Kan and contractible. Consider the cospan
\[
 *\xrightarrow{0}K\xleftarrow{1}*,
\]
where the two maps select distinct vertices. Its strict pullback is empty. Collapsing $K$ to a point gives an objectwise weak equivalence from this cospan to the constant one-point cospan, whose strict pullback is a point. Hence the ordinary limit functor does not preserve objectwise weak equivalences, even though every value of both diagrams is fibrant. The square formed by this weak equivalence over the identity of the terminal diagram is homotopy cartesian; its image under the strict limit is not, since it would require the map from the empty space to the point to be a weak equivalence. Fibrancy of values has not imposed the matching condition. The homotopy pullback of the first cospan is the space of paths from $0$ to $1$ in $K$, which is contractible, as it must be.

There is an equally concrete colimit example. Form the simplicial span
\[
 \Delta^1\longleftarrow\partial\Delta^1\longrightarrow\Delta^1
\]
using the two boundary inclusions. Its pushout is a simplicial circle. Collapsing both intervals gives an objectwise weak equivalence to $*\leftarrow\partial\Delta^1\to *$, whose strict pushout is a point. The first span consists of cofibrations and calculates the homotopy pushout; the second does not. Thus objectwise cofibrancy of all simplicial sets does not make every span a cofibrant diagram for computing colimits.

Even after deriving the right Kan extension, its preservation of homotopy pushouts is a different assertion and is generally false. Take $u$ from the discrete two-object category to the terminal category and work with pointed spaces. Then $\mathbf R u_*$ is the product. In the category of pairs, the coproduct of $(S^1,*)$ and $(*,S^1)$ is $(S^1,S^1)$. Applying the product gives the torus $S^1\times S^1$, whereas taking the coproduct after applying the product gives $S^1\vee S^1$. The comparison induces
\[
 \mathbb Z*\mathbb Z\longrightarrow\mathbb Z\times\mathbb Z
\]
on fundamental groups and is not an isomorphism. Since a pointed homotopy coproduct is a homotopy pushout over the zero object, this supplies the finite-poset example requested in the manuscript.

Restriction has a different behavior. Once the diagram localizations form a derivator, restrictions between shifted derivators preserve homotopy Kan extensions in the independent diagram variable, and hence both cartesian and cocartesian squares; this is \citet[Proposition~2.5]{Groth2013}. The two claims about restriction are therefore compatible. The asymmetry belongs to a general left or right Kan extension, and disappears in stable situations for \emph{finite} homotopy limits and colimits, where the two notions coincide.
\end{dcomment}


On applique le corollaire 1 à $\alpha_s : e \longrightarrow I$.

\textbf{Scholie.} La proposition 1 (relation (1)) montre que le foncteur $u^*$ transforme carrés $W$-cocartésiens en itou, suites exactes longues de suspension en itou – c'est “l'exactitude à droite homotopique” de $u^*$. On s'attend à ce qu'il en soit de même pour les carrés cartésiens et les suites exactes

\noindent\msp{page 49}

\noindent longues de cosuspension, mais pour le prouver il faudra un travail, qu'on fera dans la suite.

Les relations (2) montrent que le foncteur $u_*$ transforme carrés $W$-cartésiens en itou, donc suites exactes longues de cosuspension en itou. Je ne m'attends \textit{pas} à ce qu'il en soit de même en général pour les carrés $W$-cocartésiens – \textit{i.e.} que $u_*$ soit “homotopiquement exact à droite”, et non seulement à gauche. Mais il faudrait au moins que je fasse des exemples (avec $I'$, $I$ des catégories finies ordonnées).

Quant à $u_{!}$, la proposition ne dit rien de bien palpable à son sujet. Néanmoins on s'attend à ce que $u_{!}$ transforme carrés $W$-cocartésiens en itou – c'est la moindre des choses ! Ou peut-être je fais erreur, en confondant les foncteurs $u_*$, $u_{!}$ envisagés ici, avec les foncteurs similaires
\[
\begin{tikzcd}
\HO_W(I') \arrow[r, shift left, "u_{!}^W"] \arrow[r, shift right, "u_*^W"'] & \HO_W(I),
\end{tikzcd}
\]
adjoints à gauche resp. à droite de
\[
\begin{tikzcd}
\HO_W(I') & \HO_W(I), \arrow[l, "u_W^*"']
\end{tikzcd}
\]
induit par
\[
u^* \;:\; (\M(I),W(I)) \;\longrightarrow\; (\M(I'),W(I')).
\]
Mais en général $u_*$, $u_{!}$ ne sont pas compatibles

\noindent\msp{page 50}

\noindent avec les localiseurs $W(I')$, $W(I)$, donc ils ne doivent pas être confondus avec les foncteurs $u_{!}^W$, $u_*^W$. Donc j'ai été peut-être hâtif, dans ce que j'ai dit au sujet de $u_*$ et de $u_{!}$. Je serai obligé d'étudier la situation de très près, quand j'essayerai de construire les foncteurs $u_*^W$ et $u_{!}^W$. Pour le moment, je n'ai encore aucune assurance que les conditions envisagées pour $(\M,W)$ suffisent à assurer l'existence de $u_*^W$, $u_{!}^W$ – pas même pour $I$, $I'$ des ensembles ordonnés finis !

\textbf{Corollaire 1.} \textit{Soit $X \in \M(I)_{\mathrm{f}}$, alors pour toute $\alpha : s \longrightarrow t$ dans $I$ qui est} \texttt{[un]} monomorphisme, \textit{$X(\alpha) : X(t) \longrightarrow X(s)$ est dans $F$.}

% foot: Version 25 janvier 2004

% ---- from Der-ChXIII-25-01-04_p031-035.tex
\printedpage{29}
% head: A. Grothendieck : Dérivateurs, Ch. XIII, Édition des Universités de Montpellier II et Paris VII
(On applique \texttt{[le]} corollaire 1 (p. 47) à $\tilde{\alpha} : \boldsymbol{\Delta}_1 \longrightarrow I$.)

\textbf{Proposition 3.} (\footnote[28]{Comparer avec le lemme p. 63 et sa démonstration (plus exactement, le corollaire du lemme).}) \textit{Soit $X$ dans $\mathcal{M}(\boldsymbol{\Delta}_1)$, $X = (X_0 \xleftarrow{\alpha} X_1)$. Alors $X$ est dans $\mathcal{M}(\boldsymbol{\Delta}_1)_{\mathrm{f}}$ si et seulement si $X_0$, $X_1$} \texttt{[sont]} \textit{dans $\mathcal{M}_{\mathrm{f}}$ et $\alpha \in F$.}

\textbf{Corollaire 2.} \textit{Soit $I' \overset{u}{\hookrightarrow} I$} \texttt{[une]} immersion ouverte. \textit{Alors pour $X \in \operatorname{Ob}\mathcal{M}(I)_{\mathrm{f}}$, $\varprojlim\limits_{I} X(i) \longrightarrow \varprojlim\limits_{I'} X(i)$ est dans $F$.}

\section*{5\quad Catégories de modèles quillénisables}

Je reviens au cas général d'une catégorie de modèles $(\mathcal{M},W)$.

\dgrblock{dgrb699}{The two finite lifting arguments}
\begin{dcomment}{exposition}{p50:Je reviens au cas général d'une catégorie de modèles}{The two finite lifting arguments}
Grothendieck isolates consequences of diagram fibrancy which can be proved before any large existence theorem. The walking-arrow criterion of Proposition~3 is the special case $Y=*$ of the matching-map calculation discussed above: $X=(X_1\to X_0)$ is injectively fibrant precisely when $X_0$ is fibrant and $X_1\to X_0$ is a fibration. The fibrancy of $X_1$ then follows by composition. This proves the proposition without an appeal to the later lemma cited in footnote~28.

The monomorphism hypothesis in Corollary~1 has an exact combinatorial role. Let $\alpha:s\to t$ be monic and let $a:[1]\to I$ classify it. For any object $i$, the category $i\downarrow a$ has objects given by maps $i\to s$ or $i\to t$. A map of the first kind is joined to its composite with $\alpha$; a map of the second kind has at most one such predecessor because $\alpha$ is monic. Thus every component is either a singleton or a single arrow, and has an initial object. Footnote~27 shows that $a^*$ preserves injective fibrancy. Applying Proposition~3 to $a^*X$ gives $X(t)\to X(s)\in F$. No assumption that $I$ is itself a poset is needed.

For the restriction of limits, let $j:J\hookrightarrow I$ be a sieve, the presheaf convention for the manuscript's open immersion, and let $X$ be injectively fibrant. We prove directly that
\[
 \lim_I X\longrightarrow\lim_J j^*X
\]
has the right lifting property against $TC$. Take a square against $a:A\to B\in TC$. The upper map is a cone from $A$ to $X$ over $I$, while the lower map is a cone from $B$ over $J$ extending it. Since $j$ is a sieve, left extension $j_!$ is extension by the initial object outside $J$. The two cones combine to a map into $X$ from
\[
 T=(\operatorname{const}_I A)
 \mathbin{\amalg}_{j_!\operatorname{const}_J A}
 j_!\operatorname{const}_J B.
\]
The canonical map $T\to\operatorname{const}_I B$ is the identity at objects of $J$ and is $a$ elsewhere. It belongs to $TC(I)$. The fibrancy of $X$ supplies a lift $\operatorname{const}_I B\to X$, which is exactly the missing cone. This proves the corollary whenever the displayed limits and the required finite colimits exist. The construction exhibits the compatibility information in the lifting problem instead of suppressing it in a pointwise calculation.
\end{dcomment}


% marginal vertical bar printed beside the following Définition 1
\textbf{Définition 1.} On dit que $(\mathcal{M},W)$ est \textit{quillénisable} s'il existe une structure de Quillen faible \textit{exacte} $(C,TC,F,TF)$ pour $\mathcal{M}$.

C'est là une condition autoduale, mais que je ne sais vérifier en pratique que moyennant l'exactitude soit à gauche, soit à droite -- donc par une condition plus forte qui détruit la symétrie. C'est pour la sauvegarder dans la mesure du possible que j'introduis la définition précédente, ainsi que la

\textbf{Définition 2.} Soit, en plus de $(\mathcal{M},W)$, Diag $\subset$ Cat une sous-catégorie pleine de Cat. On dit que
\msp{page 51}$W$ est \textit{quillénisable sur} Diag, si pour tout $I$ dans Diag, $(\mathcal{M}(I),W(I))$ est quillénisable.

En pratique, j'utiliserai cette définition dans le cas où Diag $=$ Cat, ou alors, qu'il satisfasse au moins les propriétés de stabilité utiles pour le domaine d'un dérivateur : stable par $\varprojlim$ finies, par sommes finies -- eventuellement par $\varinjlim$ finies, etc (\footnote[29]{Je supposerai que Diag contient tout au moins les catégories provenant d'ensembles ordonnés finis.}). % [sic] eventuellement -> éventuellement, missing accent

\dgrblock{dgrb700}{What quillenizability promises}
\begin{dcomment}{exposition}{p51:En pratique, j'utiliserai cette définition}{What quillenizability promises}
Grothendieck wants a condition on the relative category $(\mathcal M,W)$, independent of the auxiliary choices of cofibrations and fibrations. Quillenizability says that some exact weak structure exists. Quillenizability on a domain of diagrams repeats this existence statement for every $I$ in the domain. It does not itself choose compatible structures. The exact-completion argument in the commentary to pp.~32--34 does allow one to enlarge the two trivial classes to $C\cap W$ and $F\cap W$, recovering the original, possibly nonclosed Quillen axioms without changing $C$, $F$ or $W$. Thus exactness has real force even before closure is imposed. The adjective ``weak'' must nevertheless retain the definition from p.~27. It is not a synonym for the later notions of a semi-model category or of a weak model category.

There are three separate questions. First, can one form the localization $\mathcal M(I)[W(I)^{-1}]$? Second, does it carry the fibre and cofibre operations constructed by replacement? Third, do restriction and Kan extension relate these categories with the exactness and pointwise formulas required of a derivator? Existence of localizations answers the first question, an exact weak structure addresses the second, and the uniform constructions aim at the third. Treating these three tasks separately explains why the author returns to examples from ordinary homological algebra rather than simply assuming arbitrary colimits from the outset.

His quotation marks around ``triangulée'' are also essential. In a pointed unstable theory there are suspension and loop functors, and long exact sequences obtained from fibres or cofibres. The suspension need not be an equivalence, the homotopy category need not be additive, and these data are not a triangulated category in Verdier's sense. Stable derivators do yield canonical triangulations, with exact restriction and Kan extension functors; the stability hypothesis is the additional mathematical content in \citet[Theorem~4.16]{Groth2013}. For complexes, shifting is invertible on the unbounded derived category, whereas nonnegative chain complexes provide a pointed but generally unstable example. This is precisely the distinction the next example needs.
\end{dcomment}


\dgrblock{dgrb701}{Quillenisability and the strength of its qualifiers}
\begin{dcomment}{exposition}{p51:En pratique, j'utiliserai cette définition}{Quillenisability and the strength of its qualifiers}
Grothendieck is looking for hypotheses on a category with weak equivalences which survive passage to diagrams. The point is not merely to construct another presentation of its homotopy category. A derivator must retain enough of the original category to restrict a diagram, replace it, and eventually take its homotopy limit. His succession of adjectives records different parts of this programme; reading them all as the modern expression ``model category'' would erase the question he is asking.

Throughout these pages we use his variance:
\[
 \mathcal M(I)=\operatorname{Fun}(I^{\mathrm{op}},\mathcal M),
 \qquad W(I)=\{f\mid f_i\in W\text{ for every }i\}.
\]
Thus the ordered category $0\to1$ indexes an arrow $X_1\to X_0$ in $\mathcal M$. A terminal object of $I$ computes the ordinary limit of such a diagram. This explains several formulas which would have the opposite orientation with covariant diagrams.

The glossary supplied with the chapter makes the terminology precise. A weak Quillen structure has the two factorizations $(C,TF)$ and $(TC,F)$, their lifting properties, and $TC\subset C\cap W$, $TF\subset F\cap W$. The axiom called exactness is
\[
                         W=TF\circ TC.
\]
Exactness on the left is the stronger-looking equality $TC=C\cap W$; exactness on the right is $TF=F\cap W$. Closedness says that the four classes are determined by their respective lifting properties. The additional axiom $\mathbf M''$ concerns preservation of weak equivalences under base change along fibrations and cobase change along cofibrations. In a usual model structure these are the properness requirements. Finally, $\mathbf M0$ asks only for finite limits and finite colimits. Small completeness must therefore be supplied when an infinite Kan extension is used.

There is a useful simplification precisely when \emph{closedness and exactness are both present}, as they will be in \S6. It is the retract argument, not an extra smallness theorem. Suppose $f\in C\cap W$. Exactness supplies $f=pi$ with $i\in TC$ and $p\in TF$. The left lifting property of $f$ against $p$ makes $f$ a retract of $i$ in the arrow category. A class defined by a left lifting property is closed under retracts, so $f\in TC$. For $f\in F\cap W$, use the same factorization and lift $i$ against $f$; this exhibits $f$ as a retract of $p$, and hence puts it in $TF$. Consequently
\[
              TC=C\cap W,\qquad TF=F\cap W
\]
under the closed exact hypotheses. This application of the familiar retract argument is particularly consequential here: see \citet[Lemma 1.1.9]{Hovey1999}. The proof uses neither arbitrary products nor a prior assertion that $W$ is closed under retracts.

Quillenisability by itself, in Definition 1, does not include closedness. We therefore do not substitute the stronger hypothesis into \S5. Once a uniform closed exact quillenisability has been assumed in \S6, however, the preceding equalities are available in every diagram category. They will allow the apparently separate questions about finite posets, products, towers and group actions to be treated by one argument. Keeping this distinction visible preserves both the breadth of the definition and the actual force of the theorem's hypotheses.
\end{dcomment}


\textbf{Exemple 1 :} Celui des paragraphes 3, 4, où on a dégagé des conditions sur $(\mathcal{M},W)$, avec $\mathcal{M}$ une ``grosse'' catégorie et $\mathcal{M}$ accessible, pour que $W$ soit quillénisable sur Cat tout entier. C'est le cas qui me paraît à présent le plus intéressant.

\textbf{Exemple 2 :} Soient $\mathcal{A}$ une catégorie abélienne, et $\mathcal{M} = K^{\bullet}(A)$ la catégorie des complexes de cochaînes à valeurs dans $\mathcal{A}$, $W = W_{\mathcal{A}}$ l'ensemble des quasi-isomorphismes, donc % [sic] K^\bullet(A) -> K^\bullet(\mathcal{A}), plain italic A where calligraphic A is used everywhere else
\[ \mathcal{M}W^{-1} \;=\; K^{\bullet}(\mathcal{A})W_{\mathcal{A}}^{-1} \;\simeq\; D^{\bullet}(\mathcal{A}), \qquad \text{catégorie dérivée de $\mathcal{A}$.} \]
On définit aussi $\mathcal{M}^{+}$, $\mathcal{M}^{-}$, $\mathcal{M}^{\mathrm{b}}$, avec les localiseurs induits $W_{\mathcal{A}}^{+}$, $W_{\mathcal{A}}^{-}$, $W_{\mathcal{A}}^{\mathrm{b}}$, et les catégories de fractions $D^{+}(\mathcal{A})$, $D^{-}(\mathcal{A})$, $D^{\mathrm{b}}(\mathcal{A})$. Ou
\msp{page 52}encore, pour $I$ dans Cat,
\[ \mathcal{M}(I) \;=\; K^{\bullet}(\mathcal{A}(I)) \]
% foot: Version 25 janvier 2004
\printedpage{30}
% head: A. Grothendieck : Dérivateurs, Ch. XIII, Édition des Universités de Montpellier II et Paris VII
\[ W(I) \;=\; \underbrace{W_{\mathcal{A}(I)}}_{\text{quasi-isomorphismes}} \;\subset\; \mathrm{Fl}(K^{\bullet}(\mathcal{A}(I))), \]
et de même pour $\mathcal{M}^{\alpha}(I)$, $W^{\alpha}(I)$ $(\alpha \in \{+,-,\mathrm{b}\})$, du moins \textit{si $I$ est finie.} Il y a aussi les sous-catégories $K^{\geq 0}(\mathcal{A})$ et $K^{\leq 0}(\mathcal{A})$ des complexes de cochaînes resp. de chaînes au sens strict (avec $K^{i} = 0$ si $i \geq 0$, resp. si $i \leq 0$ \texttt{[plutôt si $i < 0$, resp. si $i > 0$]}). On aura pour $I$ quelconque
\[ \left\{ \begin{aligned} K^{\geq 0}(\mathcal{A})(I) &\;\simeq\; K^{\geq 0}(\mathcal{A}(I)) \\ W_{\mathcal{A}}^{\geq 0}(I) &\;=\; W_{\mathcal{A}(I)}^{\geq 0}, \end{aligned} \right. \]
et de même pour $K^{\leq 0}$, $W^{\leq 0}$.

D'après \textsc{Quillen} page 1.2, $(K^{+}(\mathcal{A}),W_{\mathcal{A}}^{+})$ est quillénisable si $\mathcal{A}$ a assez d'injectifs, $(K^{-}(\mathcal{A}),W_{\mathcal{A}}^{-})$ l'est si $\mathcal{A}$ a assez de projectifs. Il ne dit rien pour $(K^{\bullet}(\mathcal{A}),W_{\mathcal{A}})$ sans plus -- pourtant on fait de l'algèbre cohomologique avec $D^{\bullet}(\mathcal{A})$ sans plus, et même sans faire d'hypothèse d'existence d'injectifs ou de projectifs. J'ai un peu oublié ce que \textsc{Verdier} a fait à ce sujet -- je devrais quand même regarder sa thèse (!), ou bien l'exposé de \textsc{Hartshorne}.

\dgrblock{dgrb702}{Two meanings of a bounded diagram of complexes}
\begin{dcomment}{examples}{p52:je devrais quand même regarder sa thèse}{Two meanings of a bounded diagram of complexes}
Grothendieck's qualification ``du moins si $I$ est finie'' protects an important distinction. We use $\operatorname{Ch}(\mathcal A)$ for his $K^\bullet(\mathcal A)$, which here denotes complexes and chain maps, before quotienting by chain homotopy. For unbounded complexes there is a genuine identification
\[
 \operatorname{Fun}(I^{\mathrm{op}},\operatorname{Ch}(\mathcal A))
 \cong\operatorname{Ch}(\operatorname{Fun}(I^{\mathrm{op}},\mathcal A)).
\]
Both sides specify objects $X_i^n$, differentials, and structure maps satisfying the same commutation relations. Quasi-isomorphisms are objectwise because kernels, cokernels, and hence cohomology in an abelian functor category are computed objectwise. The same argument works for a fixed half-line of degrees, such as $n\geq0$.

For $\operatorname{Ch}^+$ there are two possible bounds. An object of $\operatorname{Fun}(I^{\mathrm{op}},\operatorname{Ch}^+(\mathcal A))$ has a lower bound depending on $i$. An object of $\operatorname{Ch}^+(\mathcal A(I))$ has one lower bound valid for every $i$. A finite object set allows us to take the minimum of the finitely many lower bounds, so the two categories agree. An infinite object set need not. For $I=\mathbb N$ discrete and a nonzero abelian group $A$, the family $A[n]$, $n\geq0$, has one nonzero term in degree $-n$ and gives an immediate counterexample. There are analogous upper-bound and two-sided-bound distinctions.

This also explains footnote~31. In $\operatorname{Ch}^+(\mathbf{Ab})$ take the filtered system
\[
 X_m=\bigoplus_{0\leq n\leq m}\mathbb Z[n],\qquad
 X_m\longrightarrow X_{m+1}
\]
with the evident inclusions and zero differentials. Its colimit in unbounded complexes has a copy of $\mathbb Z$ in every nonpositive degree. It does not belong to $\operatorname{Ch}^+$. Thus the assertion ``$K^\bullet(\mathcal A)$ etc. have the same properties'' cannot include bounded-below or bounded-above complexes without a common bound; the manuscript itself flags this point.

More seriously, the difference can destroy an adjoint. Suppose the right adjoint to the constant-diagram functor
\[
 D^+(\mathbf{Ab})\longrightarrow
 \prod_{n\geq0}D^+(\mathbf{Ab})
\]
existed at the family $(\mathbb Z[n])_{n\geq0}$, and call its value $P$. Applying the adjunction to $\mathbb Z[k]$ gives
\[
 \begin{aligned}
 H^{-k}(P)&\cong\operatorname{Hom}_{D(\mathbf{Ab})}(\mathbb Z[k],P)\\
 &\cong\prod_{n\geq0}
 \operatorname{Hom}_{D(\mathbf{Ab})}(\mathbb Z[k],\mathbb Z[n])\\
 &\cong\mathbb Z\qquad(k\geq0).
 \end{aligned}
\]
The last equality uses the projectivity of $\mathbb Z$: only the factor $n=k$ is nonzero. No bounded-below complex can have these cohomology groups. Therefore this right adjoint does not exist. For a discrete indexing category, the localization of the objectwise bounded complexes is indeed the product of the bounded derived categories: choose a bounded-below injective resolution separately in each component and represent morphisms by chain maps between these resolutions.

Accordingly, the familiar $\mathbf R u_*:D^+(\mathcal A(I'))\to D^+(\mathcal A(I))$ refers to a \emph{uniformly} bounded-below input. It must not be used without comment for the generally larger localization of $\operatorname{Fun}((I')^{\mathrm{op}},\operatorname{Ch}^+(\mathcal A))$. This failure concerns the requested adjoint, not a proof that the latter relative category has no possible weak quillenization. The commentary on Chapter~VI, manuscript p.~131, already records a uniform-bound obstruction; here the obstruction is located directly in the proposed domain of a derived Kan extension.
\end{dcomment}


\msp{page 53}Toujours est-il qu'on a là des exemples de type tout différent de ceux envisagés dans 1°. Ainsi, il n'y a pas à supposer que $\mathcal{M}$ soit stable par petites $\varinjlim$ -- $\mathcal{M}$ peut même être une petite catégorie, elle l'est si et seulement si $\mathcal{A}$ elle même l'est. Si $\mathcal{A}$ est stable par (petites) $\varinjlim$ filtrantes, et est accessible (p. ex. a une petite famille génératrice, et les $\varinjlim$ filtrantes sont exactes), alors $K^{\bullet}(\mathcal{A})$ etc. (\footnote[30]{!!}) ont les mêmes propriétés, et $(K^{+}(\mathcal{A}),W_{\mathcal{A}}^{+})$ est quillénisable sur Cat tout entier (\footnote[31]{attention, $K^{+}(\mathcal{A})$ n'est \textit{pas} stable par $\varinjlim$ filtrantes infinies -- donc on ne peut appliquer le tapis des paragraphes 2, 3 directement à $K^{+}$.}). (C'est le cas où on sait que $\mathcal{A}$ a assez d'injectifs, et l'algèbre cohomologique se fait surtout à coups de résolutions injectives \dots) Dans ce cas on sait que les % [sic] elle même -> elle-même, missing hyphen
\[ u_{*} \;:\; D^{+}(I',\mathcal{A}) \longrightarrow D^{+}(I,\mathcal{A}) \]
existent comme adjoints à droite des
\[ u^{*} \;:\; D^{+}(I,\mathcal{A}) \longrightarrow D^{+}(I',\mathcal{A}), \]
mais je ne suis plus trop sûr si on sait définir $u_{*}$ aussi sur $D(I',\mathcal{A})$ tout entier, et je doute que $u_{!}$ existe sans hypothèse sur $u$ (genre : dimension cohomologique finie). Sans doute on doit toujours avoir
\[ u_{!} \;:\; D^{-}(I',\mathcal{A}) \longrightarrow D^{-}(I,\mathcal{A}), \]
\msp{page 54}adjoint à gauche de
\[ u^{*} \;:\; D^{-}(I,\mathcal{A}) \longrightarrow D^{-}(I',\mathcal{A}). \]
% foot: Version 25 janvier 2004

\dgrblock{dgrb703}{A bound on each component is not a bound on the diagram}
\begin{dcomment}{examples}{p53}{A bound on each component is not a bound on the diagram}
Grothendieck turns to complexes because they offer a familiar theory of derived coefficients without the large-colimit hypotheses of the preceding sections. His notation $K^\bullet(\mathcal A)$ denotes complexes and chain maps here. To avoid conflict with the now customary use of $K(\mathcal A)$ for the chain-homotopy category, we write $\operatorname{Ch}(\mathcal A)$ for that category of complexes. Localization at quasi-isomorphisms gives the derived category in either convention.

For unbounded complexes the exchange of the two diagram directions is literal:
\[
 \operatorname{Fun}(I^{\mathrm{op}},\operatorname{Ch}(\mathcal A))
 \simeq \operatorname{Ch}(\operatorname{Fun}(I^{\mathrm{op}},\mathcal A)).
\]
Both sides are the same collection of objects, differentials, transition maps and commutativities. Quasi-isomorphisms agree because kernels, cokernels and cohomology in the functor category are computed at each object. The fixed truncation conditions $X^n=0$ for $n<0$, or for $n>0$, also commute with arbitrary diagram categories. The editor's correction to the reversed inequalities on manuscript p.~52 is essential.

The movable bound in $\operatorname{Ch}^+$ behaves differently. A functor into $\operatorname{Ch}^+(\mathcal A)$ gives a lower bound $b_i$ for each component. An object of $\operatorname{Ch}^+(\mathcal A(I))$ requires one bound valid for all components. These conditions agree for a finite set of objects, since one can take $\min_i b_i$, but need not agree otherwise. The same distinction applies to upper bounds and to bounded complexes. This is why the words ``si $I$ est finie'' belong to the argument, rather than merely to its notation.

Here is the obstruction in its simplest form. Let $k$ be a field and put $X_n=k[n]$, with our cochain convention $X[n]^j=X^{j+n}$. Thus $X_n$ has a single nonzero term in degree $-n$. Each $X_n$ is bounded below. The degreewise product
\[
                         \prod_{n\geq0} k[n]
\]
has nonzero cohomology in every degree $0,-1,-2,\ldots$ and is not bounded below. This is not merely a defect of that chosen complex. If a product $P$ of the family existed in $D^+(k)$, the defining universal property, tested on $k[m]$, would give
\[
 H^{-m}(P)=\operatorname{Hom}_{D(k)}(k[m],P)
 \cong\prod_{n\geq0}\operatorname{Hom}_{D(k)}(k[m],k[n])
 \cong k
\]
for every $m\geq0$. No object of $D^+(k)$ has this property. Thus the projection from an infinite discrete category does not have a right adjoint on the theory of \emph{objectwise} bounded-below diagrams with values in $D^+(k)$.

The filtered-colimit warning in footnote 31 has an equally concrete witness. Let $P_r=\bigoplus_{0\leq n\leq r}k[n]$ and include $P_r$ into $P_{r+1}$. Each $P_r$ is bounded, whereas their degreewise filtered colimit has nonzero terms arbitrarily far to the left. Hence $\operatorname{Ch}^+$, and likewise $\operatorname{Ch}^{\mathrm b}$, is not closed under arbitrary filtered colimits in $\operatorname{Ch}$. The dual example uses $k[-n]$ to treat $\operatorname{Ch}^-$. The unbounded category and the two fixed half-line categories do not have this problem.

Enough injectives do give the familiar bounded-below resolutions, and enough projectives give bounded-above resolutions. What the argument on these pages does not establish is quillenisability on every infinite diagram shape for the category allowing a different bound at every object. The accessibility theorem quoted from the preceding sections cannot be applied directly to that category. We record this as a missing justification, not as a counterexample to quillenisability itself: the product calculation concerns a Kan extension, which a finite-limit model structure is not obliged to possess. Passing to unbounded complexes is the clean way to obtain the unrestricted diagram theory.
\end{dcomment}


\printedpage{31}
% head: A. Grothendieck : Dérivateurs, Ch. XIII, Édition des Universités de Montpellier II et Paris VII
Mais je ne sais plus trop s'il faut s'attendre à pouvoir définir $u_{!}$ sur $D(I',\mathcal{A})$ tout entier, en dehors d'hypothèses genre dimension cohomologique finie. La triste vérité, c'est que j'ai $\pm$ oublié la théorie des foncteurs dérivés ! (Il faudrait que je demande encore à \textsc{Jean Malgoire} de m'envoyer des textes canoniques.)

\dgrblock{dgrb704}{Unbounded resolutions remove the dimensional restriction}
\begin{dcomment}{later}{p54:La triste vérité, c'est que}{Unbounded resolutions remove the dimensional restriction}
Grothendieck's uncertainty about unbounded derived functors has a precise modern answer for Grothendieck abelian categories. An unbounded complex needs a resolution which is injective as a complex, in the homotopical sense, rather than merely injective in each degree. A complex $J$ is $K$-injective when $\operatorname{Hom}^\bullet(A,J)$ is acyclic for every acyclic complex $A$. This condition makes homotopy classes of maps into $J$ invariant under quasi-isomorphisms in the source.

\citet[Theorems~C and D]{Spaltenstein1988} constructed unbounded resolutions for modules and for sheaves of modules. For an arbitrary Grothendieck abelian category, the general existence result is stated and proved in \citet[Theorem~19.12.6, tag 079P]{StacksKInjectives}. One may also equip $\operatorname{Ch}(\mathcal A)$ with the proper injective model structure whose weak equivalences are quasi-isomorphisms and whose cofibrations are monomorphisms; see \citet[Theorem~1.2]{CisinskiDeglise2009}, where the existence theorem is attributed to Beke.

For a small category $I$, the abelian category $\mathcal A(I)$ is again Grothendieck. Exactness of filtered colimits is objectwise. A generating family is obtained by applying the left adjoints to evaluation to a generating family of $\mathcal A$; these left adjoints are coproducts indexed by hom-sets. Applying the injective and projective diagram constructions to unbounded complexes consequently gives
\[
 \mathbf L u_!\dashv u^*\dashv\mathbf R u_*
 \quad\text{on}\quad
 D(\mathcal A(I'))\ \text{and}\ D(\mathcal A(I))
\]
for every functor of small indexing categories. Here ``projective diagram structure'' refers to the diagram model structure and does not assert that $\mathcal A$ has enough projective objects. No finite cohomological dimension of $u$ is required for existence on the unbounded categories.

Finite dimensionality becomes relevant when one asks these functors to preserve a chosen bounded subcategory. This can be seen without any sophisticated example. Let $G$ be cyclic of prime order $p$, let $k=\mathbb F_p$, and consider constant coefficients $k$ over $BG$. Write $\epsilon=g-1$ in $k[G]\cong k[\epsilon]/(\epsilon^p)$. The periodic free resolution of the trivial module has differentials alternating multiplication by $\epsilon$ and $\epsilon^{p-1}$:
\[
 \cdots\xrightarrow{\epsilon}k[G]
 \xrightarrow{\epsilon^{p-1}}k[G]
 \xrightarrow{\epsilon}k[G]\longrightarrow k\longrightarrow0.
\]
It is exact because the kernel of multiplication by $\epsilon$ is the image of multiplication by $\epsilon^{p-1}$, and conversely. After tensoring with $k$, or applying $\operatorname{Hom}_{k[G]}(-,k)$, every differential is zero. Thus
\[
 H^{-n}(\mathbf L\operatorname{colim}_{BG}k)\cong k,
 \qquad
 H^n(\mathbf R\lim_{BG}k)\cong k
 \quad(n\geq0).
\]
The left derived colimit of this degree-zero input is not bounded below, and the right derived limit is not bounded above. Both nevertheless exist as unbounded complexes. The source's worries about existence and about boundedness therefore have different resolutions. For an arbitrary abelian category, existence of the Verdier localization alone still does not supply either unbounded resolutions or all these Kan extensions.
\end{dcomment}


\dgrblock{dgrb705}{From a presentation to a homotopy theory of diagrams}
\begin{dcomment}{later}{p54:Mais je ne sais plus trop s'il faut}{From a presentation to a homotopy theory of diagrams}
The uncertainty about derived direct images brings Grothendieck back to the
reason for keeping all indexing categories at once. A construction on one
localized category can conceal the choices needed to make it functorial.
Passing first to diagrams and then localizing preserves the information on
which the construction depends. In particular, a coherent square is an
object of the localization of the category of squares; it is not merely a
commutative square in the ordinary homotopy category.

Renaudin makes the comparison with Quillen presentations precise for
combinatorial model categories. In the published version, the bicategorical
localization of left proper combinatorial model categories at Quillen
equivalences maps locally equivalently to the 2-category of derivators with
adjunctions as morphisms; restricting to derivators of small presentation
gives a biequivalence \citep[Theorems~3.3.2 and~3.4.4]{Renaudin2009}.
The choice of morphisms is part of this statement. It is not an assertion
that every prederivator, or every transformation of prederivators, comes from
an arbitrary relative category in this way.

The analogous construction for an infinity-category $\mathcal C$ has the
particularly transparent form
\[
 I\longmapsto
 \operatorname{Ho}\operatorname{Fun}(NI^{\mathrm{op}},\mathcal C).
\]
Here $N$ is the nerve, and the opposite records the presheaf convention of
this chapter. Limits and colimits in $\mathcal C$ give the homotopy Kan
extensions of this system. Riehl and Verity express the relevant comparison
law as a Beck--Chevalley theorem for comma squares, and explain the resulting
derivator of a complete and cocomplete quasi-category
\citep[Proposition~5.3.9 and Remark~5.3.10]{RiehlVerity2017}.

This description also clarifies the role of the classes $C$ and $F$ in the
preceding calculations. They supply replacements that compute these
intrinsic operations. The right replacement for a presheaf is generally
more than a list of fibrant values: the transition maps carry compatibility
conditions. Grothendieck's arrow example on manuscript page~50 already
exhibits this fact. On the other hand, forgetting to derive an ordinary Kan
extension destroys precisely the homotopy invariance that the intrinsic
description requires. The withdrawal on manuscript pages~49--50 corrects
the construction while preserving its objective.
\end{dcomment}


\textbf{Scholie.} Si $(\mathcal{M},W)$ est quillénisable sur Diag $\subset$ Cat (et même sans qu'il soit quillénisable), il définit un prédérivateur $\mathbf{D}_{\mathcal{M},W}$ sur Diag. On posera parfois
\[ \mathrm{HO}_{\mathcal{M},W}(I) \;\text{ ou }\; \mathrm{HO}_{W}(I) \;=\; \mathbf{D}_{\mathcal{M},W}(I). \]
L'hypothèse quillénisable implique, sur chacune de ces catégories homotopiques, la structure (relativement faible) ``triangulée'', provenant des suites \textit{exactes} de suspension et de cosuspension (\footnote[32]{Dans le cas où $\mathcal{M}$ est ponctuée par \texttt{[un]} objet neutre. Sinon, il y a la structure associée aux carrés cartésiens et cocartésiens. Il serait assez intéressant d'axiomatiser une telle espèce de structure.}). (Le caractère quillénisable assure l'exactitude de ces suites \dots) Les foncteurs $u^{*} : \mathrm{HO}_{W}(I) \longrightarrow \mathrm{HO}_{W}(I')$
\msp{page 55}sont ``homotopiquement exacts'' dans le sens des ces structures ``triangulées'' (non commutatives), \textit{i.e.} transforment suites exactes de suspension resp. de cosuspension en itou. % [sic] des ces -> de ces, grammar

\dgrblock{dgrb706}{An answer already present in Joyal's letter}
\begin{dcomment}{context}{p54}{An answer already present in Joyal's letter}
The questions about complexes and uniform diagram structures have a particularly close antecedent in the correspondence surrounding \emph{Pursuing Stacks}. In the letter begun on 11 April 1984, Joyal describes Heller's two model structures on simplicial presheaves and explicitly connects them with derived right and left Kan extensions for an arbitrary functor of indexing categories. In the continuation dated 15 April, Theorem 2 states the injective model structure on unbounded complexes in an arbitrary Grothendieck abelian category, with monomorphisms as cofibrations and homology isomorphisms as weak equivalences \citep[pp.~3 and 10]{Joyal1984Letter}.

This locates the present hesitation accurately. The manuscript records a return to an earlier set of questions, not the historical first appearance of the unrestricted complex example. The later proofs by Beke and the constructions of $K$-injective resolutions make those statements accessible in systematic published forms. The letter is also unusually direct evidence for the role of Heller in the diagram-category approach; it supports that precise connection without requiring us to infer a chronology from similarities between later axiomatizations.
\end{dcomment}


\dgrblock{dgrb707}{Unbounded resolutions answer the derived-image question}
\begin{dcomment}{later}{p54}{Unbounded resolutions answer the derived-image question}
Grothendieck's hesitation about unbounded derived images concerns the existence of resolutions on which an ordinary functor can safely be evaluated. A complex of injective objects is not automatically suitable when it is unbounded: what is needed is a condition on the whole complex. A complex $J$ is $K$-injective when $\operatorname{Hom}^\bullet(A,J)$ is acyclic for every acyclic complex $A$. Equivalently, maps into $J$ in the homotopy category already compute maps into its image in the derived category.

Spaltenstein's work was already available when these manuscript pages were written. His Theorem C treats unbounded complexes of modules, Theorem D treats sheaves of modules, and Theorem 4.5 gives the latter existence construction; these results should not be retrospectively attributed to the present text \citep{Spaltenstein1988}. For an arbitrary Grothendieck abelian category the general result is associated with \citet{Serpe2003}. A detailed modern statement, including a functorial monomorphic quasi-isomorphism into a termwise injective $K$-injective complex, is Theorem 19.12.6, Tag 079P, of the Stacks Project \citep{StacksKInjective079P}.

The application to right Kan extensions is short once this existence theorem has been supplied. Let $u:I'\to I$ be a functor of small categories and let $\mathcal A$ be Grothendieck abelian. The category $\mathcal A(I')$ is again Grothendieck: exactness of filtered colimits is checked componentwise, and free diagrams on a set of generators give a set of generators. Restriction $u^*$ is exact, and its ordinary right adjoint $u_*$ exists. If $J$ is $K$-injective in $\operatorname{Ch}(\mathcal A(I'))$, then for any acyclic $A$ in $\operatorname{Ch}(\mathcal A(I))$ the ordinary adjunction gives
\[
 \operatorname{Hom}^\bullet(A,u_*J)
   \cong\operatorname{Hom}^\bullet(u^*A,J).
\]
The right-hand side is acyclic because $u^*A$ is acyclic. Thus $u_*J$ is $K$-injective. Replacing $X'$ by such a $J$ defines $\mathbf Ru_*X'=u_*J$ on the entire unbounded derived category, and the same Hom identity proves $u^*\dashv\mathbf Ru_*$ there.

The corresponding unbounded left Kan extensions also exist for Grothendieck coefficients. The precise model-theoretic input is \citet[Proposition 3.13]{Beke2000}: unbounded complexes have a cofibrantly generated model structure with monomorphisms as cofibrations and quasi-isomorphisms as weak equivalences. The underlying category is locally presentable, so this is a combinatorial model; its projective diagram structures and the Brown argument supply all left Kan extensions. Neither construction requires a finite cohomological dimension for the indexing functor. Such dimension restrictions instead govern preservation of bounded ranges. Since restriction is exact for the standard degreewise $t$-structures, its right adjoint preserves lower bounds and its left adjoint preserves upper bounds. Thus the usual statements on globally bounded-below and globally bounded-above diagrams remain meaningful, while a bounded object can acquire cohomology in arbitrarily many degrees on the other side.

For this reason the letter $u_!$ here must be read as the left adjoint to restriction of diagrams. It is not a direct image with proper supports on a space. Those functors have their own geometric hypotheses; importing them would introduce difficulties unrelated to the question asked on pp.~53--54. The example of a finite group below will display the genuine boundedness issue without any such ambiguity.
\end{dcomment}


Quels sont les axiomes des dérivateurs qui sont vérifiés sous la seule hypothèse quillénisable ? \textbf{Der} 1 \texttt{[voir chapitre I, p. 43]} l'est sans hypothèse aucune sur la catégorie de modèles $(\mathcal{M},W)$ :
\[ \mathrm{HO}_{W}(\coprod_{\alpha} I_{\alpha}) \;\simeq\; \prod_{\alpha} \mathrm{HO}_{W}(I_{\alpha}) \]
pour une famille \textit{finie} d'objets $I_{\alpha}$ dans Diag. L'axiome de localisation \textbf{Der} 2 \texttt{[voir chapitre I, p. 43]} l'est sous l'hypothèse que $W$ soit \textit{fortement saturé,} et quillénisable. On utilise le fait que toute flèche de $\mathrm{HO}_{W}(I)$ provient, à isomorphisme près, d'une flèche dans $\mathcal{M}(I)$. (C'était là un des points qui restait en suspens dans VI, dans la théorie de $\mathrm{HOT}_{W}$, $W \subset \mathrm{Fl}(\mathrm{Cat})$ un localiseur fondamental \dots)

\dgrblock{dgrb708}{The two elementary derivator axioms}
\begin{dcomment}{exposition}{p55:On utilise le fait que toute flèche de}{The two elementary derivator axioms}
Grothendieck is now testing how much of the derivator is already present before constructing homotopy Kan extensions. The two assertions about finite coproducts of shapes and detection of isomorphisms have different proofs, and the latter uses a homotopical input which cannot be replaced by saturation alone.

For a finite family of shapes there is an isomorphism of relative categories
\[
 \mathcal M\Bigl(\coprod_{r=1}^n I_r\Bigr)
 \cong\prod_{r=1}^n\mathcal M(I_r)
\]
with componentwise weak equivalences. Localization commutes with this finite product. One way to see the point is to separate each arrow $(f_1,\ldots,f_n)$ into the $n$ coordinate arrows, inserting identities in the other positions. Coordinate arrows commute, and this commutation remains true after inverting weak equivalences. Finite zigzags in the individual coordinates can therefore be padded with identities and assembled into a finite zigzag in the product. The same operation on the elementary relations between zigzags gives full faithfulness as well as essential surjectivity. The empty family gives the terminal category. Infinite products require an additional argument: one cannot pad zigzags of unbounded length to a single finite length. This is why the statement on the page is deliberately finite.

For detection of isomorphisms, use the representation property of the exact weak structure proved in the commentary to pp.~29--34: every arrow $\xi$ in $\mathcal M(I)[W(I)^{-1}]$ is isomorphic in the arrow category to the localization of an actual map $f:X\to Y$. If every evaluation of $\xi$ is invertible, the same is true of every $f_i$ in $\mathcal M[W^{-1}]$. Strong saturation of $W$ then gives $f_i\in W$ for all $i$, hence $f\in W(I)$, and $\xi$ is invertible. The role of representability is to put all components into one honest natural transformation before applying saturation. The commentary on Chapter~VI, manuscript p.~130, already explains why saturation by itself is insufficient.

These arguments establish the two stated axioms with their stated qualifications. They do not produce the pointwise Kan-extension formula. In the modern model-category treatment, that remaining structure is available: \citet[Proposition~1.30]{Groth2013} proves the strong derivator theorem in the combinatorial case, and his Theorem~1.32 records Cisinski's theorem for model categories in general. The word ``model'' in that latter statement has the ordinary model-category meaning. The exact-completion lemma above first recovers Quillen's original axioms from the weak structure; closedness, size conventions and the required small limits and colimits must still be tracked when invoking a modern formulation.
\end{dcomment}


L'axiome \textbf{Der} 3, nouvelle mouture \texttt{[voir chapitre I, pp. 111-112]}, implique un localiseur fondamental $\mathbf{W} \subset \mathrm{Fl}(\mathrm{Cat})$ (pas forcément $\mathbf{W}_{\infty}$) (\footnote[33]{Il est naturel ici de supposer que $\mathbf{W}$ satisfait à tous les ``bons axiomes'' sur les localiseurs fondamentaux W (1, 2, 3, 4, 5, 6, 7) \texttt{[voir chapitre XII]}.}). La forme la plus forte serait celle-ci : si
\[ u \;:\; I' \longrightarrow I \]
\msp{page 56}est dans $\mathbf{W}$, on veut que le foncteur $u^{*}$ induise une \textit{équivalence}
\[ u_{\mathrm{lc}}^{*} \;:\; \mathrm{HO}_{W}^{\mathrm{lc}}(I) \xrightarrow{\;\sim\;} \mathrm{HO}_{W}^{\mathrm{lc}}(I'). \]
Mais je doute qu'on soit armé pour vérifier un axiome aussi fort, sans avoir au moins une théorie des opérations $u_{!}$ ou $u_{*}$. On pourrait à tout hasard poser cet axiome en prenant $\mathbf{W} = \mathbf{W}_{\alpha} \overset{\text{déf}}{=}$ le plus petit localiseur fondamental satisfaisant W (1 à 7) (le plus grand de
% foot: Version 25 janvier 2004
\printedpage{32}
% head: A. Grothendieck : Dérivateurs, Ch. XIII, Édition des Universités de Montpellier II et Paris VII
ces localiseurs étant $\mathbf{W}_{\infty}$). Pour l'instant, la façon de vérifier cet axiome (en termes de la donnée $(\mathcal{M},W)$) est pour moi un mystère $\pm$ total. (Mais je sais le vérifier pour $\mathrm{HOT}_{W}$, quand $W \subset \mathrm{Fl}(\mathrm{Cat})$ satisfait W (1, 2, 3, 4, 5, \dots).)

\dgrblock{dgrb709}{Localization axioms, stability, and locally constant coefficients}
\begin{dcomment}{exposition}{p55}{Localization axioms, stability, and locally constant coefficients}
Grothendieck now asks which parts of the derivator formalism are already forced by a good calculus of weak equivalences. Two separate gains should be distinguished. Localizing categories of diagrams preserves their coherent origin; obtaining homotopy Kan extensions requires additional functors between these localized categories. The prederivator exists before the second gain has been made.

Finite disjoint unions are formal. A diagram on $I_1\amalg\cdots\amalg I_r$ is a tuple of diagrams, and weak equivalences are componentwise. Localization commutes with this finite product: finite zigzags may be padded by identities, and the relations may be checked in one coordinate at a time. This explains the absence of a quillenisability hypothesis in the stated finite version of $\mathbf{Der}1$. The same argument is not an argument for arbitrary products of unrelated relative categories, since there may be no uniform finite zigzag length. No such infinite assertion is needed here.

For $\mathbf{Der}2$, the point is to represent a localized arrow by a genuine arrow after replacing its endpoints. Let $f$ be an arrow of $\mathrm{HO}_W(I)$ whose evaluations are all invertible. Under the stated homotopy calculus, choose a representative $a:A\to B$ in $\mathcal M(I)$, with the endpoint identifications weak equivalences. Every $a_i$ is invertible after localization in $\mathcal M$. Strong saturation of $W$ implies $a_i\in W$, so $a\in W(I)$ and $f$ is invertible. The source is correctly using both ingredients: componentwise detection in the original relative category and representability of arrows in the localized diagram category. Strong saturation alone would not supply the latter.

The word ``triangulée'' has a deliberately broader meaning in this passage than the Verdier definition. In a pointed homotopy theory one has cofiber sequences, fiber sequences, suspension and loop objects. Suspension need not be invertible, and there is no reason for the homotopy category to be additive. Pointed spaces give the basic example. An actual triangulated category emerges in the stable case, where suspension and loop are inverse equivalences and cocartesian squares agree with cartesian squares. For strong stable derivators, the canonical triangulations and exactness of the induced functors are established in \citet[Theorem 4.16 and Proposition 4.18]{Groth2013}. This is a later completion of the stable part of the programme, not an assumption hidden in the quoted scholium.

The locally constant condition on p.~56 is likewise a condition on a coherent diagram. Its transition arrows become isomorphisms after localization, but the object still belongs to $\mathrm{HO}_W(I)$, not merely to $\operatorname{Fun}(I^{\mathrm{op}},\mathrm{HO}_W(e))$. The higher coherence carries information. For ordinary model categories, the precise later answer is \citet[Theorem 1.3]{Cisinski2009LocallyConstant}: a functor inducing a weak equivalence of nerves induces an equivalence between the corresponding categories of locally constant diagrams. One can interpret this as dependence only on the infinity-groupoid obtained by inverting all arrows of the indexing category. It does not assert the result for every relative category merely satisfying the definition of quillenisability in \S5; the necessary homotopy theory must still be available.

The axiom numbers are historical. The present $\mathbf{Der}3$ concerns a localizer and locally constant objects, whereas $\mathbf{Der}4$ and $\mathbf{Der}5$ below concern existence and pointwise computation of right images. Modern texts commonly number the existence and pointwise axioms as Der3 and Der4, with strength as an additional Der5. A reference to a modern numbered axiom must therefore be translated by its content.
\end{dcomment}


\dgrblock{dgrb710}{Locally constant diagrams depend on the homotopy type}
\begin{dcomment}{later}{p56:Pour l'instant, la façon de vérifier cet axiome}{Locally constant diagrams depend on the homotopy type}
Grothendieck wants a homotopy-invariance theorem for coefficients, stronger than invariance of a single homotopy limit. A locally constant diagram sends every arrow of the indexing category to a weak equivalence. Its transition maps consequently remember transport along paths, together with the higher compatibility data carried by the original diagram before localization. This is why the whole homotopy type of the indexing category, rather than only its set of connected components, enters.

For an ordinary model category $\mathcal V$, \citet[Theorem~1.3]{Cisinski2009LocallyConstant} proves the precise statement anticipated here: if $N(u)$ is a weak homotopy equivalence, restriction induces an equivalence between the homotopy categories of locally constant $\mathcal V$-valued diagrams. Corollary~1.28 gives the converse when all coefficient model categories are allowed; simplicial sets already suffice. Apply the theorem to the opposite indexing categories to match the presheaf convention of these pages. This is an established answer for the usual localizer of nerve weak equivalences, not a conjecture awaiting the construction of $u_!$ or $u_*$.

In the language of higher categories, the reason is that a locally constant diagram factors through the groupoid completion of the index. That completion is the homotopy type represented by its nerve. An equivalence of those homotopy types induces an equivalence of their categories of functors into any fixed homotopy theory. The model-category theorem makes this explanation precise while retaining strict diagrams and their objectwise weak equivalences.

Two qualifications keep this comparison faithful to the text. The source's ``Der~3, nouvelle mouture'' is an axiom about local systems and a chosen fundamental localizer; it is not the modern axiom usually numbered (Der3), which asserts existence of homotopy Kan extensions. Also, a fixed coefficient theory may detect fewer distinctions between shapes than all coefficient theories together. Cisinski's converse uses the latter quantifier. We refer to the commentary on Chapter~VI, manuscript pp.~161--165, for the detailed comparison of those two forms of detection, rather than repeating it here.
\end{dcomment}


Je vais m'en occuper, une fois que j'aurai une maîtrise raisonnable sur le formalisme de $u_{*}$, $u_{!}$ relatif aux $u : I' \longrightarrow I$ dans Diag. L'hypothèse quillénisable seule semble trop grossière pour assurer l'existence de $u_{!}$ ou de $u_{\star}$ dans des cas intéressants. % [sic] u_\star -> u_*, star symbol where u_* is used elsewhere

\noindent\msp{page 57}

\section*{6\quad Catégorie de modèles quillénisables, et foncteurs $u_{!}$, $u_{*}$}

\textbf{Définition 1 :}

\begin{itemize}[leftmargin=2.5em]
\item[\circled{a}] On dit que $(\mathcal{M},W)$ est \textit{quillénisable uniformément à gauche} sur Diag, si on peut trouver une structure de Quillen faible $(C,TC,F,TF)$ pour $\mathcal{M}$ qui soit \textit{close,} et qui de plus ait la propriété suivante : pour $I$ dans Cat, soient $C(I)$, $TC(I) \subset \mathrm{Fl}(\mathcal{M}(I))$ définis par
\[ \begin{aligned} &(u : X \longrightarrow Y) \in C(I) \text{ (resp. } TC(I)) \text{ si et seulement si } \ \forall\, s \in \operatorname{Ob} I, \\ &u_{s} : X_{s} \longrightarrow Y_{s} \text{ est dans } C(I) \text{ (resp. } TC(I)). \end{aligned} \]
Soit de plus
\[ F(I) \;=\; TC(I)_{*}, \qquad TF(I) \;=\; C(I)_{*}. \]
La condition envisagée, c'est que pour tout $I$ \textit{dans} Diag, $(C(I),TC(I),F(I),TF(I))$ est une structure de Quillen faible \textit{close exacte} pour $W(I)$ (\footnote[34]{$(C,TC,F,TF)$ est appelé alors une \textit{quillénisation uniforme à gauche} de $(\mathcal{M},W)$ sur Diag.}).

\item[\circled{b}] On dit que $(\mathcal{M},W)$ est \textit{quillénisable sur Diag uniformément et exactement} \texttt{[à]} \textit{gauche}, si on a de plus $TC = C \cap W$, \textit{i.e.} si $(C,TC,F,TF)$ est exacte à gauche, donc défini par $C$ seul, via
\[ TC = C \cap W, \qquad F = (TC)_{*}, \qquad TF = C_{*} \qquad \text{(\footnotemark[35])} \]
\footnotetext[35]{On dit que $(C,\ldots)$, ou simplement $C$, définit une \textit{quillénisation uniforme et exacte à gauche} de $(\mathcal{M},W)$ sur Diag.}
\end{itemize}

On a vu au paragraphe 3 que si $\mathcal{M}$ est stable par petites $\varinjlim$ et accessible, et $W$ stable par petites $\varinjlim$ filtrantes, enfin si
% foot: Version 25 janvier 2004
\printedpage{33}
% head: A. Grothendieck : Dérivateurs, Ch. XIII, Édition des Universités de Montpellier II et Paris VII
\msp{page 58}$W$ admet une structure de Quillen faible $(C,TC,F,TF)$ qui soit close et exacte à gauche -- alors les conditions de la définition 1 sont satisfaites.

\dgrblock{dgrb711}{Why uniform left structures produce right images}
\begin{dcomment}{exposition}{p57}{Why uniform left structures produce right images}
Grothendieck's word ``uniform'' puts the same replacement mechanism on every diagram category. With cofibrations and trivial cofibrations defined objectwise, restriction preserves both classes for the elementary reason that it merely selects components. This makes restriction the left-hand member of the relevant Quillen adjunction, even though its geometric name is ``inverse image''. Its right adjoint is therefore the functor on which fibrant replacement should be performed.

Write $u:I'\to I$ and assume that the ordinary right Kan extension exists. The lifting adjunction immediately gives
\[
 u_*F(I')\subset F(I),\qquad u_*TF(I')\subset TF(I).
\]
For example, lifting a map of $\mathcal M(I)$ in $C(I)$ against $u_*p$ is equivalent to lifting its restriction, which lies in $C(I')$, against $p\in TF(I')$. The argument for $F$ uses $TC$ instead. Since $u_*$ preserves the terminal object, it also preserves fibrant objects. These are the facts summarized before the theorem-scholie; they do not yet appeal to a formula for a homotopy limit.

The pointwise formula determines which ordinary limits must exist:
\[
 (u_*X')(s)=\varprojlim_{(i',\,u(i')\to s)\in I'/s}X'(i').
\]
Here the displayed limit is the limit of the contravariant diagram on $I'/s$. A finite category has finitely many morphisms as well as finitely many objects, so these are finite limits when $I$ and $I'$ are finite. The phrase ``$\varinjlim$ finies'' on p.~59 is a source slip: the relevant half of $\mathbf M0$ is finite \emph{limits}. For an infinite comma category, neither $\mathbf M0$ nor the word quillenisable creates the missing limit.

There are two ways to continue. The manuscript isolates preservation of weak equivalences between fibrant diagrams as a new condition called admissibility, then proves it by finite gluing. This is a sound and instructive construction of finite homotopy limits. The closed exact hypotheses also permit the general Brown argument given below, in the block ``The Brown argument makes admissibility automatic''. It shows that admissibility is automatic whenever the ordinary right Kan extension exists. Consequently the remaining substantive issue is existence of appropriate diagram structures and ordinary limits, rather than an independent homotopy-invariance property of their right adjoints.

For comparison, a combinatorial model category has both the injective and projective model structures on every small diagram category. The injective structure is the modern counterpart of the objectwise-cofibration construction used here, and the projective structure of its dual. The existence statement and its derivator application are respectively Theorem 1.28 and Proposition 1.30 of \citet{Groth2013}. The original finite-limit formulation is still useful when arbitrary small limits do not exist, as in bounded complexes; it then yields only the Kan extensions for which the underlying limit problem can be solved.
\end{dcomment}


D'autre part, notons que la proposition 1 page 45 et ses corollaires, y inclus la proposition 3 (page 50) et ses corollaires, sont valables dès lors que les quillénisations des $(\mathcal{M}(I),W(I))$ envisagés sont définies ``uniformément à gauche'' par une quillénisation uniforme à gauche (pas nécessairement exacte à gauche) de $(\mathcal{M},W)$ sur Diag.

\dgrblock{dgrb712}{Closure and exactness already imply both trivial-class equalities}
\begin{dcomment}{exposition}{p58:D'autre part, notons que la proposition}{Closure and exactness already imply both trivial-class equalities}
Grothendieck introduces uniform quillenization to remove the dependence of the cofibration classes on the shape. The exact definition on p.~57 has a useful consequence which shortens the later discussion. In each diagram category it asks for a structure which is both \emph{closed} and \emph{exact}. Together these hypotheses force the two equalities which define biexactness.

Indeed, suppose $i\in C\cap W$. Exactness factors this very map as
\[
 i=pj,\qquad j\in TC,\quad p\in TF.
\]
Since $i\in C$ and $p\in TF$, the square with top map $j$ and bottom map the identity has a lift. The usual retract argument makes $i$ a retract of $j$ in the arrow category. Closedness of $TC$ gives $i\in TC$. Thus $C\cap W=TC$. Dually, factoring $p\in F\cap W$ as $qj$ with $j\in TC$, $q\in TF$ makes $p$ a retract of $q$, so $F\cap W=TF$. These are applications of the ordinary retract argument, as in \citet[Lemma~1.1.9]{Hovey1999}; no transfinite construction is involved.

Consequently a closed exact weak structure is perfect in the manuscript's terminology. In particular, if the domain contains the terminal category, as stipulated on p.~51, Definition~1(a) already implies the additional equality in Definition~1(b). The advertised distinction between uniform quillenizability and uniform \emph{exact} left quillenizability is therefore redundant under the printed hypotheses. This is a logical consequence of those hypotheses, not a proposed emendation of the French text. If ``exact'' were removed from Definition~1(a), there would be a different and genuinely weaker notion, but that is not the definition supplied here.

The formula in Definition~1(a) has a separate typographical error: a component $u_s:X_s\to Y_s$ belongs to $C$ or $TC$, not to $C(I)$ or $TC(I)$. The latter classes contain natural transformations, so the printed condition is not well-typed. The correction ledger records the error as a slip of the edition and preserves the printed text. After reading the component condition correctly, the uniform structures are precisely injective model structures with the same underlying cofibration classes, subject here also to the properness axiom M$''$ and to the available finite or small limits.
\end{dcomment}


La quillénisation uniforme gauche nous servira à établir l'existence des images directes homotopiques $u_{*}$, pour certains $u$. La quillénisation uniforme droite, dualement, aura comme conséquence principale l'existence des $u_{!}$ homotopiques pour certains $u : I' \longrightarrow I$ dans Diag.

Je vais résumer ce que j'ai pu établir jusqu'à présent dans ce sens, dans le

\msp{page 59}\textbf{Théorème-scholie :} \textit{Soit $(C,TC,F,TF)$ une quillénisation uniforme gauche de $W$ sur} Diag $\subset$ Cat.

\begin{itemize}[leftmargin=3em]
\item[1)] \textit{Soit $u : I' \longrightarrow I$ une flèche dans} Diag. \textit{On suppose} Diag \textit{accordé à $\mathcal{M}$ de façon que pour $u$, flèche dans} Diag, \textit{les $\varprojlim$ dans $\mathcal{M}$ (de type $I'/i$) requises pour assurer l'existence de}
\[ u_{*} \;:\; \mathcal{M}(I') \longrightarrow \mathcal{M}(I) \]
\textit{existent dans $\mathcal{M}$. (C'est O.K. si $I$, $I'$} \texttt{[sont]} \textit{finies, $\mathcal{M}$ étant supposée stable par $\varinjlim$ finies de toutes façons.) On sait} (prop. 1, page 45) \textit{que $u_{*}$ induit}
\[ u_{*}^{\mathrm{f}} \;:\; \mathcal{M}(I')_{\mathrm{f}} \longrightarrow \mathcal{M}(I)_{\mathrm{f}}\,. \tag*{$(*)$} \]
\textit{On dit que $u$ est} admissible (à gauche) pour la quillénisation uniforme gauche donnée de $W$ sur Diag, \textit{si $u_{*}^{\mathrm{f}}$ est compatible avec les localiseurs $W(I')_{\mathrm{f}}$, $W(I)_{\mathrm{f}}$,} i.e. \textit{si $f' : X' \longrightarrow Y'$ dans $\mathcal{M}(I')$, $X'$, $Y'$ fibrants, $f' \in W(I')$ implique $u_{*}(f') \in W(I)$. On dit que $Z$ dans} Diag \textit{est} admissible (pour la quillénisation uniforme gauche donnée de $W$ sur Diag \dots) \textit{si le foncteur $Z \longrightarrow e$ est admissible. Ceci posé :}
\begin{itemize}[leftmargin=2.5em]
\item[a)] \textit{Pour que $u$ soit admissible, il suffit que les $I'/s$ pour $s \in \operatorname{Ob} I$, le soient. NB il y a une réciproque, quand $I$ est ordonné, donc les $I'/s \longrightarrow I'$ des immersions ouvertes.}
\item[b)] \msp{page 60}\textit{Tout ensemble ordonné fini est admissible. Plus généralement (sauf erreur) tout ensemble ordonné $I$ tel que l'espace topologique associé à $I^{\circ}$ soit noethérien (}i.e. \textit{toute suite croissante d'ouverts de $I^{\circ}$ est stationnaire, ou encore toute suite croissante de fermées de $I$ est stationnaire), et si de plus toute partie fermée irréductible de $I^{\circ}$ admet un point générique,} i.e. \textit{$I^{\circ}$ est (non seulement un espace noethérien, mais aussi) un espace topologique} sobre. \textit{(Pour $I$, cela signifie que toute partie ouverte filtrante croissante de $I$ admet un plus grand élément.)} % [sic] fermées -> fermés, agreement with masculine (ouverts / fermés)
\end{itemize}
\end{itemize}
% foot: Version 25 janvier 2004


\dgrblock{dgrb713}{The Brown argument makes admissibility automatic}
\begin{dcomment}{later}{p59}{The Brown argument makes admissibility automatic}
The purpose of admissibility is to ensure that a limit of a suitably fibrant diagram depends only on its homotopy class. Under the closed exact hypotheses actually imposed in Definition 1, this is already a consequence of the lifting and factorization axioms. We give the argument explicitly, since it both strengthens the theorem-scholie and settles the later questions without assuming the existence of limits which the source has not provided.

First recall the consequence of closedness and exactness proved above: $F\cap W=TF$. Let $T:\mathcal C\to\mathcal E$ be a functor, where $\mathcal C$ has the source's closed exact structure and $\mathcal E$ has a class $V$ containing the identity maps and satisfying two-out-of-three. Assume that $T$ carries trivial fibrations between fibrant objects into $V$. If $f:X\to Y$ is a weak equivalence between fibrant objects, factor its graph as
\[
 X\xrightarrow{i} Z\xrightarrow{p}X\times Y,
 \qquad i\in TC,\quad p\in F,\quad pi=(1_X,f).
\]
Put $p_X=\operatorname{pr}_Xp$ and $p_Y=\operatorname{pr}_Yp$. The two projections are fibrations because $X$ and $Y$ are fibrant. Consequently $Z$ is fibrant and $p_X,p_Y\in F$. The identities $p_Xi=1_X$ and $p_Yi=f$, together with $i,f\in W$, imply $p_X,p_Y\in W$. Both maps therefore lie in $TF$. Their images under $T$ belong to $V$. Since $T(p_X)T(i)=1$, two-out-of-three first gives $T(i)\in V$ and then gives
\[
                         T(f)=T(p_Y)T(i)\in V.
\]
This is the dual form of Ken Brown's lemma \citep[Lemma 1.1.12]{Hovey1999}. We have displayed its proof to identify the exact hypotheses used: finite products, factorization, composition and pullback stability of $F$, two-out-of-three, and the equality $TF=F\cap W$. There is no appeal to smallness, accessibility, or filtered colimits.

Apply this to $T=u_*$. Uniformity makes $u^*$ preserve $C(I)$, so the lifting adjunction makes $u_*$ preserve $TF(I')$. It also preserves fibrant objects. The lemma gives precisely
\[
 f\in W(I'),\quad X',Y'\in\mathcal M(I')_{\mathrm f}
 \quad\Longrightarrow\quad u_*f\in W(I).
\]
Thus every functor $u$ for which the ordinary right Kan extension exists is admissible under the stated uniform closed exact assumptions. The finite-poset theorem remains correct, but its restriction on shapes is unnecessary for this particular homotopy-invariance assertion.

The conclusion should not be enlarged in a different direction. It does not manufacture an ordinary right adjoint in a category lacking the required limits. Nor does it, by itself, prove the pointwise homotopy formula: that formula additionally needs fibrant restriction to the comma categories, as discussed at p.~72. These are genuine existence and compatibility questions. The question of preserving weak equivalences \emph{after the relevant fibrant replacements exist} is answered by the short argument above. We attribute it to the established Brown lemma, rather than present the application as a new result.
\end{dcomment}


\dgrblock{dgrb714}{Every available right Kan extension is admissible}
\begin{dcomment}{later}{p60}{Every available right Kan extension is admissible}
Grothendieck's admissibility condition asks for the exact homotopical property needed to derive $u_*$. Under the uniform hypotheses actually printed on p.~57, it holds for every functor for which the ordinary right Kan extension exists. There is no restriction to finite posets. The reason is the dual of Ken Brown's lemma, whose standard statement is \citet[Lemma~1.1.12]{Hovey1999}.

We give the short proof in the present notation. By the preceding retract argument, each uniform structure is perfect. Proposition~1 says that $u_*$ preserves $F$ and $TF$; as a right adjoint it preserves the terminal object, and hence fibrant objects. Let $f:X\to Y$ be a weak equivalence between fibrant diagrams. Factor its graph as
\[
 X\xrightarrow{j}Z\xrightarrow{q}X\times Y,
 \qquad j\in TC,\quad q\in F.
\]
The composites $q_X:Z\to X$ and $q_Y:Z\to Y$ are fibrations because $X$ and $Y$ are fibrant. Their identities $q_Xj=1_X$ and $q_Yj=f$, together with $j,f\in W$, show that both belong to $W$. Biexactness gives $q_X,q_Y\in TF$. Their images under $u_*$ therefore lie in $W$. The equality $(u_*q_X)(u_*j)=1$ and two-out-of-three show $u_*j\in W$; hence $u_*f=(u_*q_Y)(u_*j)\in W$. This is precisely admissibility.

Choosing a fibrant replacement $X\to RX$ now defines
\[
 \mathbf R u_*[X]=[u_*RX].
\]
Brown's lemma ensures independence of the replacement in the localized category. The adjunction with the derived left functor of $u^*$ is the usual derived Quillen adjunction; since restriction already preserves all objectwise weak equivalences, that left derived functor is $u^*$ itself. Thus $u^*\dashv\mathbf R u_*$. The dual argument with uniformly right-defined structures gives $\mathbf L u_!\dashv u^*$. This is exactly the construction developed in \citet[Definition~1.3.6 and Lemma~1.3.10]{Hovey1999}.

The pointwise criterion in part~(a) remains useful for calculations. In the right Kan formula, the projection $I'/s\to I'$ is a discrete fibration. When the requisite fibre coproducts exist, in particular for finite indexing categories, the left Kan extension of presheaves along such a projection is computed by those coproducts: each relevant comma component has an initial object obtained by lifting the arrow to the fibre. It therefore preserves $TC$, so restriction preserves fibrant diagrams. The admissibility of the slice then makes the pointwise limit preserve weak equivalences. If $I$ is a poset, the slice projection identifies its domain with a sieve, as the author says. Under the full uniform hypotheses, the global Brown argument has already established both sides of the proposed converse, without requiring these auxiliary coproducts.

This proof uses finite products and factorizations in the diagram categories, plus the existence of the indicated right adjoint. The parenthesis on p.~59 should invoke finite \emph{limits}; it prints finite colimits. Both are included in M~0, so this slip does not invalidate the finite-index assertion. Up to manuscript p.~60 the theorem is only stated; the proof above supplies its admissibility conclusion independently of the continuation.
\end{dcomment}


\dgrblock{dgrb715}{The noetherian sober poset condition is finite}
\begin{dcomment}{exposition}{p60}{The noetherian sober poset condition is finite}
Grothendieck tries to enlarge the finite-poset case by replacing finite induction with topological noetherian induction. For the usual Alexandrov topology attached to a poset, however, the two hypotheses he writes down already force finiteness. The observation concerns this topology, in which arbitrary intersections of open sets are open; it would not be true for arbitrary noetherian sober spaces.

To make the order convention immaterial, let $P$ be the specialization poset and take its Alexandrov opens to be the upper sets. Closed sets are lower sets. An irreducible nonempty closed set is then exactly a directed lower set. For one direction, if $D$ is directed and meets two upper open sets, a common upper bound in $D$ lies in their intersection. Conversely, if a lower set $D$ is irreducible and $x,y\in D$, the open sets $\mathord\uparrow x$ and $\mathord\uparrow y$ both meet $D$, so their intersection meets $D$; that is a common upper bound in $D$. Sobriety says that each such directed lower set is $\mathord\downarrow z$ for a unique $z$.

It follows that $P$ has no strictly increasing infinite sequence. Otherwise, from $x_0<x_1<\cdots$ form the directed lower set $D=\bigcup_n\mathord\downarrow x_n$. If $D=\mathord\downarrow z$, then $z\leq x_N$ for some $N$, whereas $x_{N+1}\leq z$, a contradiction. This is the sobriety criterion of \citet[Chapter~II, Proposition~6.4.4]{PicadoPultr2012}.

Every nonempty subset of $P$ consequently has a maximal element: starting from an element with no maximal element above it would produce such an increasing sequence. A maximal element $m$ of a subset $S$ is isolated in $S$, since $(\mathord\uparrow m)\cap S=\{m\}$. The Alexandrov space is therefore scattered, meaning that every nonempty subspace has an isolated point.

Now use noetherianity only once. Among the finite open subsets choose a maximal one, $U$; it exists because a strictly increasing sequence of finite open subsets is forbidden. If the complement $P\setminus U$ were nonempty, scatteredness would give an isolated point $x$ in that subspace. There would be an open $V$ of $P$ with $V\cap(P\setminus U)=\{x\}$. Then $U\cup V=U\cup\{x\}$ would be a strictly larger finite open set, contradicting maximality. Hence $P=U$ is finite.

Applied to $I^{\mathrm{op}}$, the proof shows that the class on p.~60 consists exactly of finite posets. Both chain conditions are used: the order condition controls sobriety, and topological noetherianity controls opens. The sequence arguments use classical choice.
\end{dcomment}
\begin{otherlanguage}{english}
\dgroverview{XIII-a}{Chapter XIII \textperiodcentered\ \textup{XIII-a} \textperiodcentered\ Overview}
\selectlanguage{english}

\noindent\textbf{Chapter XIII, first part: building the machinery that can later be forgotten.}

These sixty manuscript pages turn from the intrinsic calculus of a
localizer to the auxiliary structures that make it calculable. Grothendieck
calls a pair $(\mathcal M,W)$ a category of models when $W$ contains the
isomorphisms and satisfies two-out-of-three. This is much less than what
``model category'' ordinarily means today. He seeks classes of cofibrations
and fibrations that provide replacements, homotopy fibres and cofibres, and
eventually the two derived Kan extensions. Their purpose is to establish
properties of the localized diagram categories
\[
 \operatorname{HO}_W(I)
   =\mathcal M^{I^{\mathrm{op}}}[W(I)^{-1}],
\]
where $W(I)$ is defined objectwise. Once these constructions and their
universal properties have been established, the auxiliary classes should
cease to be part of the intrinsic answer
\citep[manuscript pp.~12--14]{Grothendieck1990XIII}.

The chapter opens with the algebra of lifting properties. In current
notation, Grothendieck's $\Phi_*$ is $\Phi^{\square}$, his $\Psi^*$ is
${}^{\square}\Psi$, and his ``couple de Quillen'' is a weak factorization
system. It is not what is now called a Quillen adjunction. The Galois
connection explains why left classes are stable under pushouts, retracts,
and continuous transfinite compositions, and why right classes have the
dual properties. The small-object theorem then supplies factorizations
from a set of generating arrows in a locally presentable category. Our
commentary proves the construction with its cardinal bound and separates
the cellular class from its retract closure. This distinction matters
because factorization can exist before the two classes are closed under
retracts.

The intrinsic classes $\operatorname{Cof}_W$ and
$\operatorname{Fib}_W$ express preservation of weak equivalences under
cobase change and base change. They are not defined by a lifting property.
Their relation to the universally trivial arrows is established using the
factorization hypotheses, not just two-out-of-three. The base-change
argument on manuscript pages~10--11 is particularly economical: factor the
change of base, apply the fibration condition along its weak-equivalence
part, and use two-out-of-three. The commentary expands this proof and gives
a three-element-poset example showing why the factorization assumption
cannot simply be omitted.

The conjectural theorem of manuscript pages~14--26 asks whether the
intrinsic classes contain small collections of tests sufficient to build
the desired structures. Here three different forms of smallness enter:
closure under filtered colimits, accessibility of a full subcategory of the
arrow category, and generation of a lifting class by a set. They do not
mean the same thing. In particular, even an accessible left lifting class
can fail to be generated by a set in the ordinary small-object sense; the
order-embedding example is worked out in the commentary. Conversely, one
specific conjecture in these pages has a precise later answer: intersections
of two set-generated left classes in a locally presentable category are
again set-generated, by the limit theorem for combinatorial cellular
categories \citep[Theorem~3.2 and Remark~3.8]{MakkaiRosicky2014Cellular}.

Smith's recognition theorem gives a useful repair of the larger programme.
One fixes a set of generating cofibrations, assumes that its right
orthogonal consists of weak equivalences, requires the trivial cofibrations
to be stable under pushouts and transfinite composition, and imposes a
solution-set condition on $W$. Under these hypotheses one obtains an
ordinary cofibrantly generated model structure
\citep[Theorem~1.7]{Beke2000}. With the intrinsic inclusions envisaged here,
the resulting structure is proper. This is a theorem with explicit extra
hypotheses, not a proof that the manuscript's filtered-colimit assumption
alone yields every assertion in its conjectural equivalence. The audit
records the individual implications and their dependencies rather than
assigning one status to the whole discussion.

The second formulation, beginning on manuscript page~27, isolates a
different issue. A weak Quillen structure is called exact when
$W=TF\circ TC$; it is left exact when $TC=C\cap W$ and right exact when
$TF=F\cap W$. These distinctions are meaningful for classes that have not
been closed under retracts. The marginal claim that retract closure of
$C,F,W$ alone makes the two one-sided exactness conditions equivalent is
false. Our example in sets takes $C=TC$ to be the monomorphisms, $F$ the
surjections, $W$ all maps, and $TF$ the surjections whose fibres are singletons
or infinite. It satisfies the weak axioms and left exactness, but excludes
the map from a two-element set to a singleton from $TF$. The missing
retract condition is visible directly.

Exactness nevertheless permits a useful completion without any retract
closure assumption. If $j\in C\cap W$, factor $j=qi$ with $i\in TC$ and
$q\in TF$. The lifting axiom makes $j$ a retract of $i$, so $j$ has the left
lifting property against every fibration. The dual argument applies to
$F\cap W$. Replacing the two designated trivial classes by $C\cap W$ and
$F\cap W$ therefore gives a biexact Quillen structure with the same
$C,F,W$. Properness supplies the required stability under base and cobase
change. This is a completion of the designated classes; the counterexample
shows why it cannot be replaced by an assertion that they were equal
already. When all four original classes are closed, the completion makes
no change, and a \emph{closed exact} weak structure is already perfect.
We use this observation again at the end of these pages to remove an
unnecessary admissibility restriction on right Quillen functors.

The discussion of homotopy on manuscript pages~29--34 concerns the actual
morphisms in the localization. Replacing objects by objects that are both
cofibrant and fibrant is not yet the same as proving that all localized
morphisms are represented by single arrows. Exactness supplies such
representatives. We first prove a description by the congruence generated
jointly by left and right homotopies. The completion then identifies each
of the usual homotopies with equality after localization for a cofibrant
source and a fibrant target. It also resolves the doubts on manuscript
pages~32--34: homotopy-cartesian and homotopy-cocartesian squares are exact
under the weak exact axioms themselves. An endpoint-lifting argument makes
the first assertion explicit; the extra hypothesis about the diagram
category used in the source's proposed proof is unnecessary for this
conclusion. Pointedness gives the suspension and loop sequences. Stability
is the additional condition needed for a triangulated structure. The
separate suggestion about raw right homotopy depends on the path-object
convention. A graph example, adapted from \citet[Example~3.6]{KurzRosicky2005},
shows that paths supplied by the designated factorization system need not
give a transitive relation without exactness, even on fibrant-cofibrant
objects. Allowing some weak-equivalence--fibration path object gives an
equivalence relation for a fibrant target by the concatenation argument
proved in the commentary.

The passage to presheaves supplies the chapter's central technical test.
Objectwise cofibrations lead to an injective diagram structure. Its
fibrations are defined by lifting against all objectwise trivial
cofibrations, and need not be merely objectwise fibrations. The calculation
\[
 (\alpha_{s!}A)(t)=\coprod_{\operatorname{Hom}_I(t,s)} A
\]
explains the coproduct hypothesis in the evaluation lemma. The subsequent
failure of componentwise factorization to solve natural lifting problems
is exhibited by explicit examples. Modern existence results for injective
structures apply in the combinatorial setting, while the cellular-category
theorem also accounts for the set-generation of the pointwise left
classes \citep[Corollary~3.4]{MakkaiRosicky2014Cellular}.

The variances then become transparent. Restriction preserves the
objectwise cofibrations and trivial cofibrations, so its right adjoint is
right Quillen. Its left adjoint is not automatically left Quillen for these
same injective structures. Above all, an ordinary Kan extension and its
derived functor are different constructions. Grothendieck explicitly
withdraws the premature identification on manuscript pages~49--50. The
commentary supplies counterexamples and the appropriate derived statements.
Under the uniform closed exact hypotheses finally adopted on manuscript
page~57, every existing ordinary right Kan extension preserves weak
equivalences between injectively fibrant diagrams by Brown's lemma.
Consequently the final admissibility conclusion holds for every such
indexing functor, beyond the finite posets singled out in the text.

Complexes give a complementary warning about size. For an infinite
indexing category, a diagram of individually bounded-below complexes need
not have a common lower bound. It therefore need not be a bounded-below
complex in the abelian category of diagrams. A family with nonzero
cohomology in successively lower degrees makes the resulting failure of a
product in $D^+(\mathbf{Ab})$ explicit. Unbounded derived operations in
Grothendieck abelian categories are instead available through suitable
resolutions and model structures; the boundedness issue concerns the
chosen subcategory and cannot be removed by renaming it.

The later derivator constructions realize the intended independence from
presentation. Cisinski supplies homotopy right Kan extensions for general
closed model categories with small limits
\citep[Theorem~6.11]{Cisinski2003}; Renaudin compares the localized
2-category of suitable combinatorial presentations with derivators of
small presentation \citep[Theorems~3.3.2 and~3.4.4]{Renaudin2009}; and the
infinity-categorical treatment recovers the same pointwise comparison
laws from Kan extensions in the underlying infinity-category
\citep[Proposition~5.3.9 and Remark~5.3.10]{RiehlVerity2017}. These results
give precise answers in their respective settings. The final
order-topological extension also becomes precise: a noetherian sober
Alexandrov $T_0$ space is finite, so under the ordinary poset topology the
apparently broader class of shapes on manuscript page~60 adds no new cases.
For the weakest recognition assertions of the manuscript, what was
established here is kept apart from what was not located in the checked
literature. An unlocated proof is not certified as an open problem.
\end{otherlanguage}
\dgrchunkend
